% =============================================================================
% Mémoire de fin d'études — Master en Intelligence Artificielle & Big Data
%
% Conception et implémentation d'une application intelligente pour
% l'automatisation du traitement et de la réponse aux dossiers d'appels
% d'offres dans le secteur du BTP — cas de SITINFRA Cameroun
%
% Compilation : ./compile.sh      (ou : latexmk -pdf main.tex)
% =============================================================================
\documentclass[12pt,a4paper,oneside,openany]{report}

% =============================================================================
% Préambule — Mémoire BTP RAG (rapport_btp)
% Palette : bleu clair institutionnel
% =============================================================================

% -----------------------------------------------------------------------------
% Encodage, langue, typographie
% -----------------------------------------------------------------------------
\usepackage[T1]{fontenc}
\usepackage[utf8]{inputenc}
\usepackage[english, main=french]{babel}
\usepackage{lmodern}
\usepackage[protrusion=true, expansion=false]{microtype}
\usepackage{csquotes}
\usepackage{textcomp}

\frenchsetup{ThinSpaceInFrenchNumbers=true}

% -----------------------------------------------------------------------------
% Mise en page
% -----------------------------------------------------------------------------
\usepackage[
  top=2.6cm,
  bottom=2.6cm,
  left=2.9cm,
  right=2.9cm,
  headheight=26pt,
  headsep=14pt,
  footskip=28pt
]{geometry}

\usepackage{setspace}
\onehalfspacing

\usepackage{emptypage}
\usepackage{ragged2e}

% -----------------------------------------------------------------------------
% Graphiques, tableaux, listes
% -----------------------------------------------------------------------------
\usepackage{graphicx}
\usepackage{float}
\usepackage{caption}
\usepackage{subcaption}
\usepackage{booktabs}
\usepackage{longtable}
\usepackage{tabularx}
\usepackage{array}
\usepackage{multirow}
\usepackage{makecell}
\usepackage{colortbl}
\usepackage{enumitem}
\usepackage{amsmath, amssymb}
\usepackage{pdflscape}

% Ornements et palette (définit btpblue, btpsky, btplight, btppale, ...)
% =============================================================================
% decorations.tex — Ornements du mémoire (palette bleu clair)
%
% Style : ornements gravés « vectorian » (pgfornament), déclinés en bleu clair.
% Inspiré des mémoires classiques à cadre ornemental, avec un vocabulaire
% graphique propre : cadre à double filet + coins fleuronnés + bandeau répété,
% filets à fleuron pour les chapitres et les sections.
%
% Toutes les commandes publiques sont préfixées \btp...
% =============================================================================

\usepackage{tikz}
\usetikzlibrary{calc}
\usepackage{pgfornament}
\usepackage{eso-pic}
\usepackage{lettrine}

% -----------------------------------------------------------------------------
% Palette — couleur principale #3B5D9C
% Les teintes d'ornement (ciel, clair, pâle, profond) sont mélangées
% à partir de btpblue, pour que cadres, filets et fleurons restent
% dans la même famille.
% -----------------------------------------------------------------------------
\definecolor{btpblue}{HTML}{3B5D9C}           % bleu principal
\colorlet{btpsky}{btpblue!58!white}           % accent (filet intérieur, volutes)
\colorlet{btplight}{btpblue!30!white}          % filets, pieds de page
\colorlet{btppale}{btpblue!10!white}           % fonds d'encadrés et tableaux
\colorlet{btpdeep}{btpblue!58!black}           % titres, numéros de chapitre
\definecolor{btpslate}{RGB}{70, 82, 108}      % gris-bleu (texte secondaire)
\definecolor{btpink}{RGB}{26, 34, 43}          % encre du corps de texte
\definecolor{btpgold}{RGB}{168, 138, 82}      % rehaut discret, accordé au royal

% Alias de compatibilité
\colorlet{maindeco}{btpblue}
\colorlet{decoSlate}{btpslate}
\colorlet{btpaccent}{btpsky}

% -----------------------------------------------------------------------------
% Numéros des ornements vectorian retenus
% (redéfinissables pour changer tout le vocabulaire graphique d'un coup)
% -----------------------------------------------------------------------------
\newcommand{\btpOrnCorner}{63}   % coin fleuronné (page de garde)
\newcommand{\btpOrnCornerLight}{41} % coin à volute (pages courantes)
\newcommand{\btpOrnBand}{88}     % filet fin (bandeaux de cadre)
\newcommand{\btpOrnRule}{85}     % filet à fleuron central (séparateurs)
\newcommand{\btpOrnFlourish}{46} % volute horizontale (chapitres)
\newcommand{\btpOrnSmall}{80}    % petit motif (en-têtes, pieds de page)
\newcommand{\btpOrnLeaf}{75}     % fleuron feuillu (titres liminaires)

% -----------------------------------------------------------------------------
% Séparateurs horizontaux
% -----------------------------------------------------------------------------
% \btpRule[largeur][couleur] : séparateur centré à fleuron
\newcommand{\btpRule}[2][0.55\linewidth]{%
  \begingroup
  \centering
  \pgfornament[width=#1, color=#2]{\btpOrnRule}%
  \par
  \endgroup
}

% Séparateur de section : volute fine sous le titre
\newcommand{\btpSectionRule}{%
  \vspace{-0.35\baselineskip}%
  \noindent
  \textcolor{btplight}{\rule{\linewidth}{0.9pt}}%
  \par\vspace{-0.55\baselineskip}%
  \noindent
  \textcolor{btpblue}{\rule{0.22\linewidth}{2.2pt}}%
  \par\vspace{0.35\baselineskip}%
}

% Fleuron de fin de partie
\newcommand{\btpEndMark}{%
  \par\vspace{0.8em}
  \begin{center}
    \pgfornament[width=3.4cm, color=btpsky]{\btpOrnSmall}
  \end{center}
  \vspace{0.4em}
}

% Fleuron large pour les pages liminaires
\newcommand{\btpLiminaryRule}{%
  \begin{center}
    \pgfornament[width=4.6cm, color=btpblue]{\btpOrnFlourish}
  \end{center}
}

% -----------------------------------------------------------------------------
% Cadre ornemental de la page de garde
% -----------------------------------------------------------------------------
% Double filet + bandeau d'ornements répétés + coins fleuronnés.
% \btpCoverFrame  : à appeler dans l'environnement titlepage.
\newcommand{\btpCoverFrame}{%
  \AddToShipoutPictureFG*{%
    \begin{tikzpicture}[remember picture, overlay]
      % --- filets ---
      \draw[btpblue, line width=1.5pt]
        ($(current page.south west)+(1.15cm,1.15cm)$)
        rectangle
        ($(current page.north east)+(-1.15cm,-1.15cm)$);
      \draw[btpblue!70!white, line width=0.55pt]
        ($(current page.south west)+(1.42cm,1.42cm)$)
        rectangle
        ($(current page.north east)+(-1.42cm,-1.42cm)$);
      % --- bandeaux d'ornements répétés (haut / bas) ---
      \pgfornamentline[color=btpblue!62!white]
        {$(current page.north west)+(2.6cm,-1.285cm)$}
        {$(current page.north east)+(-2.6cm,-1.285cm)$}{9}{\btpOrnBand}
      \pgfornamentline[color=btpblue!62!white]
        {$(current page.south west)+(2.6cm,1.285cm)$}
        {$(current page.south east)+(-2.6cm,1.285cm)$}{9}{\btpOrnBand}
      % --- coins fleuronnés ---
      \node[anchor=north west, inner sep=0] at ($(current page.north west)+(1.18cm,-1.18cm)$)
        {\pgfornament[width=2.35cm, color=btpblue]{\btpOrnCorner}};
      \node[anchor=north east, inner sep=0] at ($(current page.north east)+(-1.18cm,-1.18cm)$)
        {\pgfornament[width=2.35cm, color=btpblue, symmetry=v]{\btpOrnCorner}};
      \node[anchor=south west, inner sep=0] at ($(current page.south west)+(1.18cm,1.18cm)$)
        {\pgfornament[width=2.35cm, color=btpblue, symmetry=h]{\btpOrnCorner}};
      \node[anchor=south east, inner sep=0] at ($(current page.south east)+(-1.18cm,1.18cm)$)
        {\pgfornament[width=2.35cm, color=btpblue, symmetry=c]{\btpOrnCorner}};
    \end{tikzpicture}%
  }%
}

% Cadre allégé pour la page de titre intérieure et les pages liminaires fortes
\newcommand{\btpLightFrame}{%
  \AddToShipoutPictureFG*{%
    \begin{tikzpicture}[remember picture, overlay]
      \draw[btpblue!38!white, line width=0.9pt]
        ($(current page.south west)+(1.3cm,1.3cm)$)
        rectangle
        ($(current page.north east)+(-1.3cm,-1.3cm)$);
      \node[anchor=north west, inner sep=0] at ($(current page.north west)+(1.33cm,-1.33cm)$)
        {\pgfornament[width=1.5cm, color=btpblue]{\btpOrnCornerLight}};
      \node[anchor=south east, inner sep=0] at ($(current page.south east)+(-1.33cm,1.33cm)$)
        {\pgfornament[width=1.5cm, color=btpblue, symmetry=c]{\btpOrnCornerLight}};
    \end{tikzpicture}%
  }%
}

% -----------------------------------------------------------------------------
% Décor optionnel des pages courantes : filet dans la marge extérieure
% Le décor principal des pages courantes est porté par l'en-tête et le pied
% de page (voir preamble.tex) ; ce filet est un supplément facultatif.
% -----------------------------------------------------------------------------
\newcommand{\btpEnableMarginRule}{%
  \AddToShipoutPictureBG{%
    \begin{tikzpicture}[remember picture, overlay]
      \draw[btplight, line width=2.2pt]
        ($(current page.north east)+(-1.15cm,-2.6cm)$) --
        ($(current page.south east)+(-1.15cm,2.6cm)$);
    \end{tikzpicture}%
  }%
}

% Conservé pour compatibilité : n'active plus le filet latéral.
\newcommand{\btpEnablePageDecor}{}

% -----------------------------------------------------------------------------
% Ouverture de chapitre
% -----------------------------------------------------------------------------
% Bloc bleu clair portant le numéro, mention verticale « CHAPITRE »,
% titre en petites capitales, filet à volute.
\newcommand{\btpChapterLabel}{Chapitre}

\newcommand{\btpChapterNumberBlock}[1]{%
  \begin{tikzpicture}[baseline=(blk.center)]
    \node[
      draw=btpblue, line width=1pt,
      fill=btppale,
      rounded corners=1.5pt,
      minimum width=2.05cm, minimum height=2.05cm,
      inner sep=0pt
    ] (blk) {\fontsize{46}{46}\selectfont\bfseries\color{btpdeep}#1};
    \node[
      rotate=90, anchor=south,
      font=\scriptsize\bfseries\scshape,
      text=btpblue,
      inner sep=3pt
    ] at (blk.west) {\btpChapterLabel};
  \end{tikzpicture}%
}

% -----------------------------------------------------------------------------
% Paragraphe d'accroche de chapitre (lettrine)
% -----------------------------------------------------------------------------
\setlength{\DefaultFindent}{2pt}
\renewcommand{\LettrineFontHook}{\color{btpblue}}

% \btpIntro{D}{ans ce chapitre, ...}
%   #1 : lettre initiale (mise en lettrine)
%   #2 : suite du paragraphe d'accroche
\newcommand{\btpIntro}[2]{%
  \begingroup
  \itshape
  \lettrine[lines=3, lhang=0.14, loversize=0.05, findent=3pt, nindent=0pt]%
    {#1}{}#2%
  \par
  \endgroup
  \vspace{0.6em}
  \noindent
  \hbox to \linewidth{\hfil\pgfornament[width=3cm, color=btplight]{\btpOrnSmall}\hfil}%
  \par\vspace{0.5em}
}

% -----------------------------------------------------------------------------
% Bandeau d'image de chapitre (optionnel)
% -----------------------------------------------------------------------------
\newcommand{\btpChapterBanner}[1]{%
  \noindent
  \begin{tikzpicture}
    \clip[rounded corners=3pt]
      (0,0) rectangle (\linewidth, 2.6cm);
    \node[anchor=south west, inner sep=0] at (0,0)
      {\includegraphics[width=\linewidth]{#1}};
    \fill[btpblue, opacity=0.12] (0,0) rectangle (\linewidth, 2.6cm);
  \end{tikzpicture}%
  \par\vspace{0.4cm}%
}

% Alias documenté dans le README : placer juste avant \chapter{...}
\newcommand{\setchapterimage}[1]{\btpChapterBanner{#1}}


\usepackage[
  colorlinks=true,
  linkcolor=btpdeep,
  citecolor=btpblue,
  urlcolor=btpsky,
  filecolor=btpblue,
  bookmarksnumbered=true,
  bookmarksopen=true
]{hyperref}

\usepackage{titlesec}
\usepackage{tocloft}
\usepackage{fancyhdr}
\usepackage[most]{tcolorbox}
\usepackage{listings}
\usepackage{listingsutf8}

% -----------------------------------------------------------------------------
% Chemins d'images
% -----------------------------------------------------------------------------
\graphicspath{%
  {images/page_garde/}%
  {images/liminaire/}%
  {images/chapitre1/}%
  {images/chapitre2/}%
  {images/chapitre3/}%
  {figures/}%
}

\newcommand{\enspylogofile}{logoenspy.png}

% -----------------------------------------------------------------------------
% Colonnes de tableau
% -----------------------------------------------------------------------------
\newcolumntype{L}[1]{>{\raggedright\arraybackslash}p{#1}}
\newcolumntype{C}[1]{>{\centering\arraybackslash}p{#1}}
\newcolumntype{R}[1]{>{\raggedleft\arraybackslash}p{#1}}
\newcolumntype{Y}{>{\raggedright\arraybackslash}X}

\renewcommand{\arraystretch}{1.28}
\setlength{\tabcolsep}{5pt}
\arrayrulecolor{btplight}

% En-tête de tableau sur fond bleu clair
\newcommand{\btpth}[1]{\textcolor{btpdeep}{\bfseries #1}}
\newcommand{\btpTableHeadRule}{\rowcolor{btppale}}

% Environnement de tableau normalisé
\newenvironment{btptable}[2][htbp]
  {\begin{table}[#1]\centering\small\caption{#2}}
  {\end{table}}

% -----------------------------------------------------------------------------
% Légendes
% -----------------------------------------------------------------------------
\captionsetup{
  font=small,
  labelfont={bf, color=btpdeep},
  textfont=it,
  format=plain,
  justification=centering,
  singlelinecheck=false,
  skip=6pt
}
\captionsetup[table]{position=above, skip=8pt}
\captionsetup[figure]{position=below}

% -----------------------------------------------------------------------------
% Titres : chapitre, section, sous-section
% -----------------------------------------------------------------------------
% Blocs de fin de titre construits en primitives TeX : les environnements
% « center » et \par sont interdits dans le code « after » de titlesec.
\newcommand{\btpChapterAfterRule}{%
  \vskip 5pt
  \hbox to \textwidth{\color{btpblue}\leaders\hrule height 1.15pt\hfill}%
  \vskip -7pt
  \hbox to \textwidth{\hfil\pgfornament[width=3.4cm, color=btpsky]{\btpOrnFlourish}}%
  \vskip -6pt
}

\newcommand{\btpChapterAfterFlourish}{%
  \vskip 7pt
  \hbox to \textwidth{\hfil\pgfornament[width=5cm, color=btpblue]{\btpOrnFlourish}\hfil}%
  \vskip -8pt
}

\titleformat{\chapter}[display]
  {\normalfont\filright}
  {\hfill\btpChapterNumberBlock{\thechapter}}
  {6pt}
  {\filright\Large\bfseries\color{btpdeep}}
  [\btpChapterAfterRule]
\titlespacing*{\chapter}{0pt}{-26pt}{24pt}

% Chapitres non numérotés (liminaires, conclusion, annexes)
\titleformat{name=\chapter,numberless}[display]
  {\normalfont\filcenter}
  {}
  {0pt}
  {\filcenter\Large\bfseries\color{btpdeep}}
  [\btpChapterAfterFlourish]
\titlespacing*{name=\chapter,numberless}{0pt}{-16pt}{24pt}

\titleformat{\section}
  {\normalfont\Large\bfseries\color{btpdeep}}
  {\thesection}{0.7em}{}
  [\btpSectionRule]
\titlespacing*{\section}{0pt}{1.5\baselineskip}{0.5\baselineskip}

\titleformat{\subsection}
  {\normalfont\large\bfseries\color{btpblue}}
  {\thesubsection}{0.6em}{}
\titlespacing*{\subsection}{0pt}{1.1\baselineskip}{0.35\baselineskip}

\titleformat{\subsubsection}
  {\normalfont\normalsize\bfseries\itshape\color{btpslate}}
  {\thesubsubsection}{0.5em}{}
\titlespacing*{\subsubsection}{0pt}{0.9\baselineskip}{0.25\baselineskip}

\setcounter{tocdepth}{2}
\setcounter{secnumdepth}{3}

% -----------------------------------------------------------------------------
% Table des matières / figures / tableaux
% -----------------------------------------------------------------------------
\renewcommand{\cftchapfont}{\bfseries\color{btpdeep}}
\renewcommand{\cftchappagefont}{\bfseries\color{btpdeep}}
\renewcommand{\cftsecfont}{\color{btpink}}
\renewcommand{\cftsubsecfont}{\small\color{btpslate}}
\renewcommand{\cftchapleader}{\color{btplight}\cftdotfill{\cftdotsep}}
\renewcommand{\cftsecleader}{\color{btplight}\cftdotfill{\cftdotsep}}
\renewcommand{\cftfigfont}{\color{btpink}}
\renewcommand{\cfttabfont}{\color{btpink}}
\setlength{\cftbeforechapskip}{8pt}

% Titres des tables : même traitement que les chapitres non numérotés
\newcommand{\btpTocTitleFont}{\hspace*{\fill}\Large\bfseries\color{btpdeep}}
\newcommand{\btpTocAfterTitle}{%
  \hspace*{\fill}\par
  \vskip 7pt
  \hbox to \textwidth{\hfil\pgfornament[width=5cm, color=btpblue]{\btpOrnFlourish}\hfil}%
  \vskip 12pt
}
\renewcommand{\cfttoctitlefont}{\btpTocTitleFont}
\renewcommand{\cftloftitlefont}{\btpTocTitleFont}
\renewcommand{\cftlottitlefont}{\btpTocTitleFont}
\renewcommand{\cftaftertoctitle}{\btpTocAfterTitle}
\renewcommand{\cftafterloftitle}{\btpTocAfterTitle}
\renewcommand{\cftafterlottitle}{\btpTocAfterTitle}
\setlength{\cftbeforeloftitleskip}{0pt}
\setlength{\cftbeforelottitleskip}{0pt}
\setlength{\cftbeforetoctitleskip}{0pt}

% -----------------------------------------------------------------------------
% En-têtes et pieds de page
% -----------------------------------------------------------------------------
\pagestyle{fancy}
\fancyhf{}
\fancyhead[L]{\small\itshape\textcolor{btpslate}{\nouppercase{\leftmark}}}
\fancyhead[R]{\small\itshape\textcolor{btpslate}{\nouppercase{\rightmark}}}
\fancyfoot[C]{%
  \pgfornament[width=2.2cm, color=btplight]{\btpOrnSmall}%
  \hspace{0.4em}\small\bfseries\textcolor{btpblue}{\thepage}\hspace{0.4em}%
  \pgfornament[width=2.2cm, color=btplight, symmetry=v]{\btpOrnSmall}%
}
\renewcommand{\headrulewidth}{0.4pt}
\renewcommand{\footrulewidth}{0pt}
\renewcommand{\headrule}{%
  \hbox to\headwidth{\color{btplight}\leaders\hrule height \headrulewidth\hfill}%
}

\fancypagestyle{plain}{%
  \fancyhf{}
  \fancyfoot[C]{%
    \pgfornament[width=2.2cm, color=btplight]{\btpOrnSmall}%
    \hspace{0.4em}\small\bfseries\textcolor{btpblue}{\thepage}\hspace{0.4em}%
    \pgfornament[width=2.2cm, color=btplight, symmetry=v]{\btpOrnSmall}%
  }
  \renewcommand{\headrulewidth}{0pt}
}

\fancypagestyle{empty}{\fancyhf{}\renewcommand{\headrulewidth}{0pt}}

% -----------------------------------------------------------------------------
% Encadrés thématiques
% -----------------------------------------------------------------------------
\tcbuselibrary{breakable, skins}

\tcbset{
  btpbase/.style={
    breakable,
    arc=2.2pt,
    boxrule=0.9pt,
    left=4mm, right=4mm, top=2.4mm, bottom=2.4mm,
    fonttitle=\bfseries,
    enhanced,
    attach boxed title to top left={xshift=6mm, yshift=-2.6mm},
    boxed title style={arc=1.5pt, boxrule=0pt},
  }
}

% Point clé / synthèse
\newtcolorbox{btpbox}[1][Point clé]{
  btpbase,
  colback=btppale, colframe=btpblue,
  coltitle=white,
  boxed title style={colback=btpblue},
  title={#1}
}

% Définition
\newtcolorbox{btpdef}[1][Définition]{
  btpbase,
  colback=white, colframe=btpsky,
  coltitle=white,
  boxed title style={colback=btpsky},
  title={#1}
}

% À retenir / remarque
\newtcolorbox{btpnote}[1][À retenir]{
  btpbase,
  colback=btplight!22, colframe=btpsky,
  coltitle=btpdeep,
  boxed title style={colback=btplight},
  title={#1}
}

% Limite / point de vigilance
\newtcolorbox{btpwarn}[1][Point de vigilance]{
  btpbase,
  colback=btpgold!8, colframe=btpgold,
  coltitle=white,
  boxed title style={colback=btpgold},
  title={#1}
}

% Citation encadrée (problématique, hypothèses)
\newtcolorbox{btpquote}{
  breakable, enhanced,
  colback=btppale, colframe=btpblue,
  boxrule=0pt, leftrule=3.5pt,
  arc=0pt, outer arc=0pt,
  left=5mm, right=4mm, top=3mm, bottom=3mm,
  fontupper=\itshape
}

% -----------------------------------------------------------------------------
% Blocs de code / schémas ASCII
% -----------------------------------------------------------------------------
\lstset{inputencoding=utf8, extendedchars=true}

\lstdefinestyle{btpascii}{
  basicstyle=\ttfamily\scriptsize\color{btpink},
  backgroundcolor=\color{btppale},
  frame=single,
  framerule=0.7pt,
  rulecolor=\color{btplight},
  breaklines=false,
  showstringspaces=false,
  tabsize=2,
  xleftmargin=0.2cm,
  xrightmargin=0.2cm,
  aboveskip=0.8em,
  belowskip=0.8em,
  literate=
    {é}{{\'e}}1 {è}{{\`e}}1 {ê}{{\^e}}1 {ë}{{\"e}}1
    {à}{{\`a}}1 {â}{{\^a}}1 {ç}{{\c c}}1
    {î}{{\^i}}1 {ï}{{\"i}}1 {ô}{{\^o}}1
    {ù}{{\`u}}1 {û}{{\^u}}1 {ü}{{\"u}}1
    {É}{{\'E}}1 {È}{{\`E}}1 {Ê}{{\^E}}1 {À}{{\`A}}1 {Ç}{{\c C}}1
    {°}{{\textdegree}}1 {€}{{\texteuro}}1
}

\lstdefinestyle{btpjson}{
  language=,
  basicstyle=\ttfamily\scriptsize\color{btpink},
  backgroundcolor=\color{btppale},
  keywordstyle=\color{btpdeep}\bfseries,
  stringstyle=\color{btpblue},
  commentstyle=\color{btpslate}\itshape,
  frame=single,
  framerule=0.7pt,
  rulecolor=\color{btplight},
  breaklines=true,
  showstringspaces=false,
  tabsize=2,
  xleftmargin=0.3cm,
  xrightmargin=0.3cm,
  morestring=[b]",
  morekeywords={true, false, null}
}

\lstset{style=btpascii}

% Environnement pour les schémas en art ASCII, avec légende
\newcommand{\btpschema}[2]{%
  \begin{figure}[htbp]
    \centering
    \begin{minipage}{\linewidth}
      \lstinputlisting[style=btpascii]{#1}
    \end{minipage}
    \caption{#2}
  \end{figure}%
}

% -----------------------------------------------------------------------------
% Listes
% -----------------------------------------------------------------------------
\setlist[itemize,1]{label=\textcolor{btpblue}{\small$\blacktriangleright$}, leftmargin=*, itemsep=2pt, topsep=4pt}
\setlist[itemize,2]{label=\textcolor{btpsky}{\small$\bullet$}, itemsep=1pt, topsep=2pt}
\setlist[enumerate,1]{label=\textcolor{btpblue}{\bfseries\arabic*.}, leftmargin=*, itemsep=2pt, topsep=4pt}
\setlist[description]{font=\bfseries\color{btpdeep}, leftmargin=*, itemsep=2pt, topsep=4pt}

% -----------------------------------------------------------------------------
% Styles TikZ partagés par les figures
% -----------------------------------------------------------------------------
\usetikzlibrary{arrows.meta, positioning, calc, fit, shapes.geometric,
                shapes.multipart, backgrounds, chains, decorations.pathreplacing,
                babel}

\tikzset{
  btpbox/.style={
    draw=btpblue, fill=btppale, line width=0.8pt,
    rounded corners=2.5pt, align=center,
    font=\footnotesize, inner sep=4pt,
    minimum height=8mm
  },
  btpboxalt/.style={
    btpbox, fill=btplight!45, draw=btpsky
  },
  btpboxstrong/.style={
    btpbox, fill=btpblue, draw=btpdeep, text=white, font=\footnotesize\bfseries
  },
  btpboxplain/.style={
    btpbox, fill=white, draw=btplight
  },
  btpstep/.style={
    circle, draw=btpblue, fill=white, line width=0.8pt,
    inner sep=1.2pt, minimum size=5.5mm,
    font=\scriptsize\bfseries, text=btpdeep
  },
  btparrow/.style={
    -{Stealth[length=2.2mm, width=1.6mm]},
    draw=btpblue, line width=0.8pt
  },
  btparrowlight/.style={
    -{Stealth[length=2mm, width=1.4mm]},
    draw=btpsky, line width=0.7pt
  },
  btpdashed/.style={
    btparrow, dashed
  },
  btpgroup/.style={
    draw=btpsky, dashed, line width=0.7pt,
    rounded corners=4pt, fill=btplight!14,
    inner sep=6pt
  },
  btplabel/.style={
    font=\scriptsize\itshape, text=btpslate, align=center
  },
  btptitleband/.style={
    fill=btpblue, text=white, font=\small\bfseries\scshape,
    minimum height=7mm, align=center, rounded corners=2pt
  },
  btpdb/.style={
    cylinder, shape border rotate=90, aspect=0.22,
    draw=btpblue, fill=btppale, line width=0.8pt,
    align=center, font=\scriptsize, inner sep=3pt,
    minimum height=11mm, minimum width=15mm
  },
  btpdecision/.style={
    diamond, aspect=1.9, draw=btpblue, fill=btplight!45,
    line width=0.8pt, align=center, font=\scriptsize, inner sep=1pt
  }
}

% Macro homonyme du style TikZ : légende secondaire dans le contenu d'un nœud
\newcommand{\btplabeltext}[1]{{\scriptsize\itshape\color{btpslate}#1}}
\let\btplabel\btplabeltext

% -----------------------------------------------------------------------------
% Chapitres liminaires / annexes non numérotés
% -----------------------------------------------------------------------------
% \btpLiminary{Titre} : chapitre non numéroté, inscrit dans la table des
% matières et repris dans les en-têtes courants.
\newcommand{\btpLiminary}[1]{%
  \chapter*{#1}%
  \addcontentsline{toc}{chapter}{#1}%
  \markboth{#1}{#1}%
}

% -----------------------------------------------------------------------------
% Raccourcis sémantiques
% -----------------------------------------------------------------------------
\newcommand{\sigle}[1]{\textsc{#1}}
\newcommand{\anglais}[1]{\textit{#1}}
\newcommand{\btpterm}[1]{\textbf{#1}}
\newcommand{\oui}{\textcolor{btpblue}{\textbf{Oui}}}
\newcommand{\non}{\textcolor{btpslate}{Non}}
\newcommand{\cotefort}{\textcolor{btpblue}{$\bullet\bullet\bullet$}}
\newcommand{\cotemoyen}{\textcolor{btpsky}{$\bullet\bullet\circ$}}
\newcommand{\cotefaible}{\textcolor{btpslate}{$\bullet\circ\circ$}}

% Étoiles de priorité (remplace les émojis du brouillon)
\newcommand{\prio}[1]{%
  \textcolor{btpblue}{%
    \ifcase#1\relax
      \or $\star$%
      \or $\star\star$%
      \or $\star\star\star$%
      \or $\star\star\star\star$%
      \or $\star\star\star\star\star$%
    \fi
  }%
}

% Marqueurs de maturité / pertinence
\newcommand{\btpMature}{\textcolor{btpblue}{\textbf{Mature}}}
\newcommand{\btpEmergent}{\textcolor{btpgold}{\textbf{Émergent}}}
\newcommand{\btpEnDev}{\textcolor{btpsky}{\textbf{En développement}}}

% Numérotation des figures et tableaux par chapitre
\renewcommand{\thefigure}{\thechapter.\arabic{figure}}
\renewcommand{\thetable}{\thechapter.\arabic{table}}

% =============================================================================
% cover.tex — Page de garde et page de titre intérieure
%
% Toutes les données modifiables sont regroupées dans le bloc « MÉTADONNÉES »
% ci-dessous : il suffit de les éditer, les pages se recomposent seules.
% =============================================================================

% -----------------------------------------------------------------------------
% MÉTADONNÉES DU MÉMOIRE
% -----------------------------------------------------------------------------
\newcommand{\btpTitre}{Conception et implémentation d'une application
  intelligente pour l'automatisation du traitement et de la réponse aux
  dossiers d'appels d'offres dans le secteur du BTP}
\newcommand{\btpSousTitre}{Cas du bureau d'études SITINFRA Cameroun}
\newcommand{\btpTitreCourt}{Application intelligente de réponse aux dossiers
  d'appels d'offres dans le BTP}

\newcommand{\btpAuteur}{KEMAYOU Lisette Raymonne Sheylley}
\newcommand{\btpDiplome}{Master en Intelligence Artificielle \& Big Data}
\newcommand{\btpDernierDiplome}{Licence en Intelligence Artificielle \& Big Data}

\newcommand{\btpDirecteur}{Dr KAMENI HOMTE Jaurès Styve}
\newcommand{\btpDirecteurGrade}{Chargé de cours, Université de Yaoundé~I}
\newcommand{\btpCoDirecteur}{M. DIFFOUO TAZO Evariste}
\newcommand{\btpCoDirecteurGrade}{Enseignant, Keyce Informatique \& IA}
\newcommand{\btpEncadrantPro}{M. Denis SITIO TCHAPDA}
\newcommand{\btpEncadrantProGrade}{Encadrant professionnel, SITINFRA Cameroun}

\newcommand{\btpJuryPresident}{Prénom et Nom, Grade}
\newcommand{\btpJuryRapporteur}{Prénom et Nom, Grade}
\newcommand{\btpJuryMembre}{Prénom et Nom, Grade}
\newcommand{\btpJuryInvite}{Prénom et Nom, Titre ou fonction}

\newcommand{\btpAnnee}{2025 / 2026}
\newcommand{\btpDateSoutenance}{\ldots\ldots\ldots\ldots}
\newcommand{\btpVille}{Yaoundé, Cameroun}

% -----------------------------------------------------------------------------
% En-tête institutionnel bilingue
% -----------------------------------------------------------------------------
\newcommand{\btpSep}{\textcolor{btpsky}{\rule{1.5cm}{0.5pt}}}

\newcommand{\makeenspyheadercompact}{%
  \noindent
  \begin{minipage}[c]{0.355\textwidth}
    \centering\scriptsize\color{btpdeep}
    \textbf{UNIVERSITÉ DE YAOUNDÉ~I}\\[1pt]
    \btpSep\\[1pt]
    \textbf{ÉCOLE NATIONALE SUPÉRIEURE\\ POLYTECHNIQUE DE YAOUNDÉ}\\[1pt]
    \btpSep\\[1pt]
    \textbf{DÉPARTEMENT DE GÉNIE\\ INFORMATIQUE}\\[1pt]
    \btpSep\\[1pt]
    \textbf{INSTITUT SUPÉRIEUR KEYCE\\ INFORMATIQUE \& INTELLIGENCE\\ ARTIFICIELLE}
  \end{minipage}%
  \hfill
  \begin{minipage}[c]{0.20\textwidth}
    \centering
    \includegraphics[width=0.82\linewidth]{logopolytech.png}
  \end{minipage}%
  \hfill
  \begin{minipage}[c]{0.355\textwidth}
    \centering\scriptsize\color{btpdeep}
    \textbf{UNIVERSITY OF YAOUNDÉ~I}\\[1pt]
    \btpSep\\[1pt]
    \textbf{NATIONAL ADVANCED SCHOOL\\ OF ENGINEERING OF YAOUNDÉ}\\[1pt]
    \btpSep\\[1pt]
    \textbf{DEPARTMENT OF COMPUTER\\ ENGINEERING}\\[1pt]
    \btpSep\\[1pt]
    \textbf{HIGHER INSTITUTE KEYCE\\ COMPUTER SCIENCE \& ARTIFICIAL\\ INTELLIGENCE}
  \end{minipage}
  \par
}

% -----------------------------------------------------------------------------
% PAGE DE GARDE
% -----------------------------------------------------------------------------
\newcommand{\makecovertitle}{%
  \begin{titlepage}
    \btpCoverFrame
    \centering
    \setstretch{1.0}

    \vspace*{-0.35cm}
    \makeenspyheadercompact

    \vspace{0.2cm}
    \pgfornament[width=0.58\textwidth, color=btpblue]{\btpOrnRule}

    \vspace{0.1cm}
    {\large\bfseries\scshape\color{btpblue} Mémoire de fin d'études}\\[1pt]
    {\footnotesize\itshape\color{btpslate} Master's Thesis}

    \vspace{0.25cm}

    % --- Bloc titre ---
    \begin{tcolorbox}[
      enhanced, sharp corners,
      colback=btppale, colframe=btpblue,
      boxrule=0pt, leftrule=3pt, rightrule=3pt,
      borderline north={1.2pt}{0pt}{btpblue},
      borderline south={1.2pt}{0pt}{btpblue},
      left=4mm, right=4mm, top=3mm, bottom=3mm,
      width=\textwidth
    ]
      \centering
      {\large\bfseries\color{btpdeep} \btpTitre}\\[5pt]
      {\normalsize\itshape\color{btpblue} \btpSousTitre}
    \end{tcolorbox}

    \vspace{0.25cm}

    {\small Présenté et soutenu par~:}\\[3pt]
    {\large\bfseries\scshape\color{btpdeep} \btpAuteur}

    \vspace{0.22cm}
    {\small En vue de l'obtention du~:}\\[2pt]
    {\normalsize\bfseries\color{btpblue} Diplôme de \btpDiplome}

    \vspace{0.22cm}
    \pgfornament[width=3.2cm, color=btpsky]{\btpOrnSmall}

    \vspace{0.18cm}
    {\small\bfseries\color{btpslate} Sous la direction de~:}\\[3pt]
    {\small\bfseries \btpDirecteur}\\[0.5pt]
    {\scriptsize\itshape\color{btpslate} \btpDirecteurGrade}\\[2pt]
    {\small\bfseries \btpCoDirecteur}\\[0.5pt]
    {\scriptsize\itshape\color{btpslate} \btpCoDirecteurGrade}\\[2pt]
    {\small\bfseries \btpEncadrantPro}\\[0.5pt]
    {\scriptsize\itshape\color{btpslate} \btpEncadrantProGrade}

    \vspace{0.22cm}
    {\small\bfseries\color{btpslate} Devant le jury composé de~:}\\[3pt]
    {\footnotesize
      \begin{tabular}{@{}r@{~~}l@{}}
        \textbf{Président~:}   & \btpJuryPresident  \\
        \textbf{Rapporteur~:}  & \btpJuryRapporteur \\
        \textbf{Membre~:}      & \btpJuryMembre     \\
        \textbf{Invité~:}      & \btpJuryInvite     \\
      \end{tabular}%
    }

    \vfill

    % --- Logos partenaires ---
    \begin{minipage}{\textwidth}
      \centering
      \begin{minipage}{0.29\textwidth}\centering
        \includegraphics[height=1.05cm, keepaspectratio]{logokeyce.jpg}
      \end{minipage}\hfill
      \begin{minipage}{0.29\textwidth}\centering
        \includegraphics[height=1.05cm, keepaspectratio]{logoiugee.png}
      \end{minipage}\hfill
      \begin{minipage}{0.29\textwidth}\centering
        \includegraphics[height=1.05cm, keepaspectratio]{logositinfra.png}
      \end{minipage}
    \end{minipage}

    \vspace{0.2cm}
    \pgfornament[width=0.40\textwidth, color=btpblue]{\btpOrnRule}\\[2pt]
    {\small\bfseries\color{btpdeep} Année académique \btpAnnee}\\[1pt]
    {\footnotesize\color{btpslate} Soutenu le \btpDateSoutenance}
    \vspace{-0.2cm}
  \end{titlepage}
}

% -----------------------------------------------------------------------------
% PAGE DE TITRE INTÉRIEURE
% -----------------------------------------------------------------------------
\newcommand{\makeinnertitle}{%
  \thispagestyle{empty}
  \btpLightFrame
  \begin{center}
    \setstretch{1.15}
    \vspace*{2.2cm}

    {\footnotesize\scshape\color{btpslate} Université de Yaoundé~I
      --- École Nationale Supérieure Polytechnique de Yaoundé}\\[2pt]
    {\footnotesize\scshape\color{btpslate} Institut Supérieur Keyce
      Informatique \& Intelligence Artificielle}

    \vspace{1.4cm}
    \pgfornament[width=5.2cm, color=btpblue]{\btpOrnFlourish}
    \vspace{1.2cm}

    {\LARGE\bfseries\color{btpdeep} \btpTitre\par}
    \vspace{0.7cm}
    {\large\itshape\color{btpblue} \btpSousTitre\par}

    \vspace{1.2cm}
    \textcolor{btplight}{\rule{0.5\linewidth}{1.4pt}}
    \vspace{1.2cm}

    {\normalsize Mémoire en vue de l'obtention du diplôme de~:}\\[6pt]
    {\large\bfseries\color{btpblue} \btpDiplome}

    \vspace{1.1cm}
    {\normalsize Présenté par~:}\\[6pt]
    {\Large\bfseries\scshape\color{btpdeep} \btpAuteur}\\[8pt]
    {\small\itshape\color{btpslate} Dernier diplôme obtenu~: \btpDernierDiplome}

    \vfill
    \pgfornament[width=3.8cm, color=btpsky]{\btpOrnSmall}\\[8pt]
    {\normalsize\color{btpslate} \btpVille}\\[3pt]
    {\normalsize\bfseries\color{btpdeep} Année académique \btpAnnee}
    \vspace{1.6cm}
  \end{center}
  \clearpage
}


\hypersetup{
  pdftitle={\btpTitreCourt},
  pdfauthor={KEMAYOU Lisette Raymonne Sheylley},
  pdfsubject={RAG, LLM, dossiers d'appels d'offres, BTP, marchés publics},
  pdfkeywords={intelligence artificielle, RAG, DAO, BTP, SITINFRA}
}

\begin{document}
\color{btpink}

% =============================================================================
% PAGES DE TITRE
% =============================================================================
\hypersetup{pageanchor=false}
\makecovertitle
\makeinnertitle
\hypersetup{pageanchor=true}

% =============================================================================
% PAGES LIMINAIRES (numérotation romaine)
% =============================================================================
\pagenumbering{roman}
\setcounter{page}{1}

% =============================================================================
% Dédicace
% =============================================================================
\thispagestyle{empty}
\btpLightFrame

\vspace*{6cm}
\begin{center}
  {\Large\bfseries\scshape\color{btpdeep} Dédicace}

  \vspace{1cm}
  \pgfornament[width=4.6cm, color=btpblue]{\btpOrnFlourish}
  \vspace{2cm}

  {\Large\itshape\color{btpslate} À ma famille,}\\[0.8cm]
  {\large\itshape\color{btpslate} pour son soutien indéfectible}\\[4pt]
  {\large\itshape\color{btpslate} et sa confiance de tous les instants.}

  \vspace{2cm}
  \pgfornament[width=3.2cm, color=btpsky]{\btpOrnSmall}
\end{center}

\clearpage

% =============================================================================
% Remerciements
% =============================================================================
\btpLiminary{Remerciements}

Aucune entreprise humaine ne peut être menée à bien uniquement par des efforts
individuels. Nous souhaitons donc exprimer notre reconnaissance à toutes celles
et ceux qui nous ont soutenus et conseillés tout au long de ce travail.

Nous pensons tout particulièrement à~:

\begin{itemize}[itemsep=4pt]
  \item \textbf{Pr \ldots\ldots}, président du jury, pour l'honneur qu'il nous
        fait en acceptant de présider cette soutenance et pour l'attention
        précieuse qu'il a portée à ce travail~;

  \item \textbf{Dr \ldots\ldots}, notre examinateur, pour son regard critique,
        sa disponibilité et ses observations constructives, qui ont
        considérablement enrichi ce mémoire~;

  \item \textbf{Dr KAMENI HOMTE Jaurès Styve}, notre directeur de mémoire, pour
        sa disponibilité, sa rigueur scientifique et la qualité de son
        encadrement tout au long de ce travail. Ses conseils avisés et ses
        remarques constructives ont grandement contribué à la progression et à
        la finalisation de ce mémoire~;

  \item \textbf{M. DIFFOUO TAZO Evariste}, notre co-directeur de mémoire, pour
        sa disponibilité constante et la qualité de son accompagnement
        rigoureux et éclairé tout au long de cette recherche~;

  \item \textbf{M. Denis SITIO TCHAPDA}, notre encadrant professionnel, pour
        son dévouement, ses orientations et les connaissances qu'il nous a
        transmises tout au long de cette étude~;

  \item \textbf{M. YOUTOU Jean}, vice-président du Collège de Paris, pour
        l'ensemble des efforts qu'il consacre à l'éducation et au développement
        des jeunes~;

  \item \textbf{M. ASSOU Frank}, directeur de Keyce Informatique \&
        Intelligence Artificielle, pour la mise à disposition des meilleures
        conditions académiques~;

  \item \textbf{M. MANGA Blaise}, \anglais{campus manager} de Keyce
        Informatique \& IA Yaoundé, pour son professionnalisme, son encadrement
        et l'attention portée au bon déroulement de notre parcours académique~;

  \item \textbf{M. ABDOURAMAN Dalil}, pour sa disponibilité constante, sa
        bienveillance, la qualité de ses enseignements et ses conseils éclairés
        tout au long de notre cursus~;

  \item \textbf{M. FOMEKONG Evaris}, pour sa rigueur, sa disponibilité et la
        qualité de ses enseignements tout au long de notre cursus~;

  \item \textbf{Toute l'équipe de SITINFRA Cameroun}, pour sa collaboration et
        le soutien constant dont elle a fait preuve tout au long de ce travail~;

  \item \textbf{Tous les enseignants et l'administration} de Keyce Informatique
        et Intelligence Artificielle, pour la rigueur et l'excellence des
        enseignements reçus durant la formation~;

  \item \textbf{Nos frères et s\oe urs}, pour leur soutien, leurs
        encouragements et leur présence constante~;

  \item \textbf{Nos camarades de promotion}, pour leur esprit d'entraide et de
        solidarité, et pour les moments de partage qui ont rendu notre parcours
        académique plus enrichissant~;

  \item \textbf{Toutes les personnes} qui, de près ou de loin, nous ont apporté
        leur soutien, leurs encouragements et leurs conseils tout au long de
        notre parcours académique. Leur aide, même non mentionnée
        individuellement, a été précieuse pour la réalisation de ce mémoire.
\end{itemize}

\btpEndMark
\clearpage

% =============================================================================
% Résumé (FR) et Abstract (EN)
% =============================================================================
\btpLiminary{Résumé}

Dans le secteur du Bâtiment et des Travaux Publics (BTP) au Cameroun, la
réponse aux Dossiers d'Appel d'Offres (DAO) constitue une activité stratégique
mais chronophage pour les bureaux d'études. Au sein de SITINFRA Cameroun, ce
processus demeure largement manuel~: l'analyse des pièces administratives,
techniques et financières ainsi que la préparation des offres mobilisent
plusieurs ingénieurs pendant de longues périodes, avec des risques d'omission
et une forte pression sur des délais compris entre 15 et 45~jours.

Ce mémoire pose la question de la conception d'un assistant intelligent capable
d'exploiter de manière fiable et sécurisée le patrimoine documentaire de
l'entreprise afin d'automatiser l'analyse des DAO et la génération de réponses
conformes au Code des Marchés Publics camerounais. L'approche retenue repose
sur la Génération Augmentée par Récupération (\anglais{Retrieval-Augmented
Generation}, RAG), enrichie d'une ingestion contextuelle, d'une récupération
hybride, d'un mécanisme correctif (CRAG) et d'une orchestration maîtrisée, afin
de concilier la performance des grands modèles de langage avec les exigences de
conformité, de traçabilité et de sécurité documentaire.

Une démarche méthodologique fondée sur la science de la conception
(\anglais{Design Science Research}) structure la modélisation de la solution, et
un protocole d'évaluation moderne (RAGAS, \anglais{LLM-as-a-Judge}) est proposé
pour valider les hypothèses de recherche. Les résultats obtenus établissent une
base technique robuste pour l'assistance à la préparation des offres et
l'amélioration de la compétitivité des bureaux d'études du BTP.

\vspace{0.6em}
\begin{btpbox}[Mots-clés]
Intelligence artificielle~; Génération Augmentée par Récupération (RAG)~;
Dossier d'Appel d'Offres (DAO)~; BTP~; marchés publics~; capitalisation
documentaire~; supervision humaine~; SITINFRA Cameroun.
\end{btpbox}

\btpEndMark
\clearpage

% -----------------------------------------------------------------------------
\btpLiminary{Abstract}
\selectlanguage{english}

In Cameroon's Building and Public Works (BTP) sector, responding to tender
documents (\anglais{Dossiers d'Appel d'Offres}, DAO) is a strategic yet
time-consuming activity for engineering consultancies. At SITINFRA Cameroun,
this process remains largely manual: analysing administrative, technical and
financial documents and preparing bids mobilises several engineers over long
periods, with a high risk of omissions and strong pressure within deadlines of
15 to 45 days.

This thesis addresses the design of an intelligent assistant able to exploit the
company's documentary assets in a reliable and secure way, in order to automate
the analysis of tender documents and to generate responses that comply with the
Cameroonian Public Procurement Code. The proposed approach is based on
Retrieval-Augmented Generation (RAG), enhanced with contextual ingestion,
hybrid retrieval, a corrective mechanism (CRAG) and controlled orchestration,
so as to reconcile the performance of large language models with the
requirements of compliance, traceability and documentary security.

A Design Science Research methodology structures the modelling of the solution,
and a modern evaluation protocol (RAGAS, LLM-as-a-Judge) is proposed to validate
the research hypotheses. The results establish a robust technical foundation for
supporting bid preparation and improving the competitiveness of BTP engineering
firms.

\vspace{0.6em}
\begin{btpbox}[Keywords]
Artificial Intelligence; Retrieval-Augmented Generation (RAG); Tender Documents
(DAO); Public Works (BTP); Public Procurement; documentary capitalisation;
human-in-the-loop; SITINFRA Cameroun.
\end{btpbox}

\selectlanguage{french}
\btpEndMark
\clearpage

% =============================================================================
% Liste des sigles, abréviations et acronymes
% =============================================================================
\btpLiminary{Liste des sigles et acronymes}

\begin{longtable}{@{}L{2.6cm} L{\dimexpr\linewidth-3.2cm\relax}@{}}
\toprule
\rowcolor{btppale}
\btpth{Sigle} & \btpth{Signification} \\
\midrule
\endfirsthead
\toprule
\rowcolor{btppale}
\btpth{Sigle} & \btpth{Signification} \\
\midrule
\endhead
\bottomrule
\endfoot

\textbf{ANN}     & \anglais{Approximate Nearest Neighbour} — recherche approximative de plus proches voisins \\
\textbf{AO}      & Appel d'Offres \\
\textbf{API}     & \anglais{Application Programming Interface} — interface de programmation applicative \\
\textbf{ARMP}    & Agence de Régulation des Marchés Publics \\
\textbf{BERT}    & \anglais{Bidirectional Encoder Representations from Transformers} \\
\textbf{BF}      & Besoin Fonctionnel \\
\textbf{BIM}     & \anglais{Building Information Modeling} \\
\textbf{BM25}    & \anglais{Best Matching 25} — modèle probabiliste de recherche lexicale \\
\textbf{BNF}     & Besoin Non Fonctionnel \\
\textbf{BPU}     & Bordereau des Prix Unitaires \\
\textbf{BTP}     & Bâtiment et Travaux Publics \\
\textbf{CCAG}    & Cahier des Clauses Administratives Générales \\
\textbf{CCAP}    & Cahier des Clauses Administratives Particulières \\
\textbf{CCTP}    & Cahier des Clauses Techniques Particulières \\
\textbf{CNN}     & \anglais{Convolutional Neural Network} — réseau de neurones convolutif \\
\textbf{CRAG}    & \anglais{Corrective Retrieval-Augmented Generation} \\
\textbf{DAO}     & Dossier d'Appel d'Offres \\
\textbf{DPGF}    & Décomposition du Prix Global et Forfaitaire \\
\textbf{DSR}     & \anglais{Design Science Research} — science de la conception \\
\textbf{GED}     & Gestion Électronique des Documents \\
\textbf{HITL}    & \anglais{Human-In-The-Loop} — supervision humaine dans la boucle \\
\textbf{HNSW}    & \anglais{Hierarchical Navigable Small World} \\
\textbf{IA}      & Intelligence Artificielle \\
\textbf{IDF}     & \anglais{Inverse Document Frequency} — fréquence inverse de document \\
\textbf{IVF}     & \anglais{Inverted File Index} — index à fichier inversé \\
\textbf{JWT}     & \anglais{JSON Web Token} \\
\textbf{KG}      & \anglais{Knowledge Graph} — graphe de connaissances \\
\textbf{LLM}     & \anglais{Large Language Model} — grand modèle de langage \\
\textbf{LSTM}    & \anglais{Long Short-Term Memory} \\
\textbf{MRR}     & \anglais{Mean Reciprocal Rank} — rang réciproque moyen \\
\textbf{NDCG}    & \anglais{Normalized Discounted Cumulative Gain} \\
\textbf{NLP}     & \anglais{Natural Language Processing} — traitement du langage naturel \\
\textbf{OCR}     & \anglais{Optical Character Recognition} — reconnaissance optique de caractères \\
\textbf{ORM}     & \anglais{Object-Relational Mapping} — correspondance objet-relationnel \\
\textbf{PQ}      & \anglais{Product Quantization} — quantification produit \\
\textbf{RAG}     & \anglais{Retrieval-Augmented Generation} — génération augmentée par récupération \\
\textbf{RAGAS}   & \anglais{Retrieval-Augmented Generation Assessment} \\
\textbf{RC}      & Règlement de Consultation \\
\textbf{RNN}     & \anglais{Recurrent Neural Network} — réseau de neurones récurrent \\
\textbf{RRF}     & \anglais{Reciprocal Rank Fusion} — fusion par rang réciproque \\
\textbf{SaaS}    & \anglais{Software as a Service} — logiciel en tant que service \\
\textbf{SQL}     & \anglais{Structured Query Language} \\
\textbf{TAL}     & Traitement Automatique du Langage \\
\textbf{TF}      & \anglais{Term Frequency} — fréquence du terme \\
\textbf{TF-IDF}  & \anglais{Term Frequency -- Inverse Document Frequency} \\
\textbf{TIC}     & Technologies de l'Information et de la Communication \\
\end{longtable}

\clearpage


\tableofcontents\clearpage
\listoffigures\clearpage
\listoftables\clearpage

% =============================================================================
% CORPS DU MÉMOIRE (numérotation arabe)
% =============================================================================
\pagenumbering{arabic}
\setcounter{page}{1}

% =============================================================================
% Introduction générale
% =============================================================================
\btpLiminary{Introduction générale}

\section*{Contexte}

Le secteur du Bâtiment et des Travaux Publics constitue l'un des piliers de
l'économie camerounaise. Il est caractérisé par une compétition intense où la
réponse aux appels d'offres représente un enjeu stratégique majeur. Chaque
année, des centaines d'entreprises et de bureaux d'études consacrent des
ressources considérables à la préparation de dossiers de candidature, un
processus qui mobilise simultanément des compétences juridiques, techniques,
financières et administratives.

L'émergence des technologies d'intelligence artificielle, et plus
particulièrement des grands modèles de langage, ouvre des perspectives inédites
pour automatiser et optimiser ce processus. Cependant, la spécificité des
appels d'offres — leur caractère hautement structuré, leur exigence de
traçabilité et de conformité, et leur dimension juridique contraignante —
impose une approche réfléchie et méthodique. Un système qui invente une
référence, omet une pièce obligatoire ou cite une clause inexistante ne
constitue pas seulement une gêne~: il expose l'entreprise à l'élimination de
son offre.

\section*{Problématique}

\begin{btpquote}
Comment concevoir et développer une plateforme permettant de capitaliser le
patrimoine documentaire d'un bureau d'études du BTP et de générer
automatiquement des dossiers de réponse aux appels d'offres en s'appuyant sur
ce patrimoine, tout en garantissant la traçabilité des sources, la conformité
réglementaire et la validation humaine des livrables~?
\end{btpquote}

Cette question centrale se décline en plusieurs interrogations
opérationnelles~:

\begin{enumerate}
  \item Comment transformer un patrimoine documentaire hétérogène, non
        structuré et partiellement numérisé en une base de connaissances
        interrogeable~?
  \item Quelle stratégie de recherche d'information permet de retrouver de
        manière fiable à la fois des références terminologiques exactes et des
        concepts métier exprimés différemment~?
  \item Comment contraindre un modèle de langage à produire un texte
        exploitable, sourcé et vérifiable plutôt qu'un texte simplement
        plausible~?
  \item Quelle place accorder à l'opérateur humain pour que le gain de
        productivité ne se paie pas d'une perte de fiabilité~?
\end{enumerate}

\section*{Objectifs}

\subsection*{Objectif général}

Concevoir, développer et valider une plateforme de capitalisation documentaire
et de génération assistée de dossiers de réponse aux appels d'offres, adaptée
au contexte d'un bureau d'études du BTP camerounais.

\subsection*{Objectifs spécifiques}

\begin{itemize}
  \item Caractériser le processus métier de réponse aux appels d'offres et en
        identifier les points de friction documentaires~;
  \item Concevoir une chaîne d'ingestion capable de préserver la structure
        hiérarchique des pièces d'un dossier d'appel d'offres~;
  \item Concevoir un moteur de recherche hybride combinant recherche lexicale
        et recherche sémantique, complété par un ré-ordonnancement fin~;
  \item Concevoir un mécanisme de génération contrainte produisant des sections
        structurées, citées et traçables~;
  \item Intégrer un dispositif de supervision humaine à chaque étape critique~;
  \item Évaluer empiriquement la performance du système sur des indicateurs de
        recherche et de génération.
\end{itemize}

\section*{Hypothèses de recherche}

\begin{btptable}{Hypothèses de recherche et critères de validation}
\begin{tabularx}{\linewidth}{@{}C{1.1cm} Y L{3.5cm}@{}}
\toprule
\rowcolor{btppale}
\btpth{Réf.} & \btpth{Hypothèse} & \btpth{Critère de validation} \\
\midrule
H1 & La recherche hybride surpasse la recherche sémantique seule sur un corpus d'appels d'offres. & Gain de rappel à 20 documents \\
H2 & Le découpage hiérarchique parent-enfant améliore la qualité des réponses par rapport à un découpage à fenêtre fixe. & Gain d'ancrage et d'exhaustivité \\
H3 & Le ré-ordonnancement par modèle croisé surpasse la seule fusion de rangs. & Gain de précision à 5 documents \\
H4 & La génération section par section améliore la conformité par rapport à une génération globale. & Gain de couverture des exigences \\
H5 & La supervision humaine améliore la fiabilité des livrables sans annuler le gain de productivité. & Taux de validation et temps de cycle \\
\bottomrule
\end{tabularx}
\end{btptable}

\section*{Organisation du mémoire}

Ce mémoire s'organise en trois chapitres. Le \textbf{chapitre~1} dresse l'état
de l'art~: il caractérise le domaine des appels d'offres dans le BTP, puis
parcourt les disciplines mobilisées, de la gestion documentaire à la génération
augmentée par récupération, en passant par la recherche d'information, la
compréhension de documents structurés et la résolution d'entités. Le
\textbf{chapitre~2} expose la méthodologie~: analyse des besoins, modélisation
des données, architecture du système, stratégies de recherche et de génération,
protocole d'évaluation et planification. Le \textbf{chapitre~3} présente la
plateforme développée, son architecture technique, les résultats de
l'évaluation et une analyse comparative des approches candidates. Une
conclusion générale récapitule les contributions et ouvre sur les perspectives.

\btpEndMark
\clearpage


% --- Chapitre 1 : État de l'art -------------------------------------------
\chapter{État de l'art}
\label{chap:etat-art}
\btpIntro{C}{e chapitre dresse un panorama des connaissances mobilisées par
notre problématique. Nous caractérisons d'abord le domaine des appels d'offres
dans le BTP et ses enjeux documentaires, puis nous parcourons les disciplines
qui outillent notre réponse~: gestion documentaire, traitement automatique du
langage, recherche d'information, génération augmentée par récupération,
compréhension de documents structurés et résolution d'entités. Nous discutons
enfin les architectures de workflow et les solutions du marché, afin de dégager
les lacunes que ce travail entend combler.}

% =============================================================================
% 1.1 Introduction
% =============================================================================
\section{Introduction}
\label{sec:ch1-intro}

Le secteur de la construction et des travaux publics constitue l'un des piliers
de l'économie, caractérisé par une compétition intense où la réponse aux appels
d'offres représente un enjeu stratégique majeur. Chaque année, des milliers
d'entreprises consacrent des ressources considérables à la préparation de
dossiers de candidature, processus qui mobilise des compétences juridiques,
techniques, financières et administratives.

L'émergence des technologies d'intelligence artificielle, et plus
particulièrement des grands modèles de langage, ouvre des perspectives inédites
pour automatiser et optimiser ce processus. Cependant, la spécificité des appels
d'offres — leur caractère hautement structuré, leur exigence de traçabilité et
de conformité, et leur dimension juridique contraignante — impose une approche
réfléchie et méthodique.

Ce chapitre dresse un panorama des connaissances actuelles dans les domaines
connexes à notre problématique. Nous procédons en trois mouvements. Le premier
caractérise le domaine d'application~: le processus des appels d'offres, sa
production documentaire et ses enjeux de capitalisation
(sections~\ref{sec:ch1-contexte} et~\ref{sec:ch1-ged}). Le deuxième parcourt
les disciplines qui outillent une réponse technique~: traitement automatique du
langage, recherche d'information, génération augmentée par récupération, bases
vectorielles, compréhension de documents structurés et résolution d'entités
(sections~\ref{sec:ch1-tal} à~\ref{sec:ch1-entites}). Le troisième discute les
architectures de systèmes et confronte notre positionnement aux solutions
existantes (sections~\ref{sec:ch1-workflows} et~\ref{sec:ch1-solutions}).

\begin{btpnote}[Périmètre de ce chapitre]
Nous distinguons soigneusement les \btpterm{concepts} des \btpterm{outils}. Un
état de l'art rend compte des modèles, des méthodes et de leurs propriétés
démontrées~; le choix d'une bibliothèque ou d'un produit relève de la
méthodologie et de l'implémentation, traitées aux chapitres~\ref{chap:methodologie}
et~\ref{chap:resultats}. Les outils cités ici le sont à titre d'illustration
d'une approche, jamais comme un choix de conception.
\end{btpnote}

% =============================================================================
% 1.2 Le contexte des appels d'offres dans le secteur BTP
% =============================================================================
\section{Le contexte des appels d'offres dans le secteur du BTP}
\label{sec:ch1-contexte}

\subsection{Définition et cadre juridique}

L'appel d'offres est une procédure par laquelle un maître d'ouvrage, public ou
privé, sollicite des offres concurrentielles pour la réalisation d'un projet de
construction, de travaux ou de services. Dans le secteur public, cette
procédure est encadrée par un corpus réglementaire qui définit les règles de
publicité, de mise en concurrence et de transparence. Au Cameroun, le Code des
Marchés Publics fixe ce cadre et confie à l'Agence de Régulation des Marchés
Publics un rôle de contrôle et de recours.

Ce cadre juridique produit une conséquence directe sur notre problématique~:
la réponse à un appel d'offres n'est pas un exercice rédactionnel libre. Elle
constitue un acte engageant, dont chaque affirmation peut être opposée au
candidat. Une référence inexacte, une attestation périmée ou une pièce manquante
entraînent l'irrégularité de l'offre, indépendamment de la qualité technique de
la proposition.

\begin{btpdef}[Dossier d'appel d'offres]
Le dossier d'appel d'offres désigne l'ensemble des pièces émises par le maître
d'ouvrage pour définir le besoin, les règles de la consultation et les
obligations contractuelles. Il se distingue du \btpterm{dossier de réponse},
produit par le candidat, dont la structure est en grande partie imposée par le
premier.
\end{btpdef}

Un dossier d'appel d'offres se compose généralement des pièces constitutives
présentées au tableau~\ref{tab:pieces-dao}.

\begin{btptable}{Pièces constitutives d'un dossier d'appel d'offres}
\label{tab:pieces-dao}
\begin{tabularx}{\linewidth}{@{}L{3.9cm} Y L{2.4cm}@{}}
\toprule
\rowcolor{btppale}
\btpth{Pièce} & \btpth{Description} & \btpth{Nature} \\
\midrule
Règlement de consultation & Définit les règles de la procédure, les modalités de remise et les critères de sélection. & Administratif \\
Cahier des clauses administratives générales & Fixe les conditions administratives applicables à une catégorie de marchés. & Juridique \\
Cahier des clauses administratives particulières & Spécifie les conditions administratives et juridiques propres au marché. & Juridique \\
Cahier des clauses techniques particulières & Définit les spécifications techniques attendues pour le projet. & Technique \\
Décomposition du prix global et forfaitaire & Structure détaillée des coûts constitutifs de l'offre. & Financier \\
Bordereau des prix unitaires & Liste les prix unitaires des prestations élémentaires. & Financier \\
Pièces annexes et administratives & Attestations, qualifications, certifications, modèles de formulaires. & Administratif \\
\bottomrule
\end{tabularx}
\end{btptable}

Cette hétérogénéité est structurante pour notre travail. Un même dossier réunit
des pièces de nature juridique, dont l'exploitation exige une fidélité littérale
aux articles, et des pièces financières tabulaires, dont l'exploitation exige la
préservation de la structure en lignes et colonnes. Aucun traitement uniforme ne
convient aux deux.

\subsection{Le processus de réponse aux appels d'offres}

Le processus de réponse à un appel d'offres est un cycle complexe impliquant de
multiples parties prenantes. La figure~\ref{fig:cycle-ao} en présente les huit
étapes canoniques.

\begin{figure}[htbp]
  \centering
  \resizebox{0.96\linewidth}{!}{% Cycle de réponse à un appel d'offres
\begin{tikzpicture}[node distance=3.6mm and 8mm]

  \node[btpbox, text width=3.3cm] (n1) {Détection\\[-1pt]\btplabel{\itshape Identification de l'opportunité}};
  \node[btpbox, text width=3.3cm, below=of n1] (n2) {Analyse\\[-1pt]\btplabel{\itshape Lecture et analyse du DAO}};
  \node[btpbox, text width=3.3cm, below=of n2] (n3) {Décision\\[-1pt]\btplabel{\itshape Go / No-go et stratégie}};
  \node[btpbox, text width=3.3cm, below=of n3] (n4) {Mobilisation\\[-1pt]\btplabel{\itshape Constitution de l'équipe}};

  \node[btpboxalt, text width=3.3cm, right=22mm of n1] (n5) {Rédaction\\[-1pt]\btplabel{\itshape Production des pièces de l'offre}};
  \node[btpboxalt, text width=3.3cm, below=of n5] (n6) {Vérification\\[-1pt]\btplabel{\itshape Contrôle de conformité}};
  \node[btpboxalt, text width=3.3cm, below=of n6] (n7) {Soumission\\[-1pt]\btplabel{\itshape Dépôt de l'offre}};
  \node[btpboxalt, text width=3.3cm, below=of n7] (n8) {Suivi\\[-1pt]\btplabel{\itshape Suivi post-soumission}};

  \foreach \i/\j in {1/2, 2/3, 3/4, 5/6, 6/7, 7/8} {
    \draw[btparrow] (n\i) -- (n\j);
  }
  \draw[btparrow] (n4.east) -- ++(0.55,0) |- (n5.west);

  \foreach \n/\num in {n1/1, n2/2, n3/3, n4/4, n5/5, n6/6, n7/7, n8/8} {
    \node[btpstep] at (\n.north west) {\num};
  }

  \begin{scope}[on background layer]
    \node[btpgroup, fit=(n1)(n4), label={[btplabel]below:Phase amont}] {};
    \node[btpgroup, fit=(n5)(n8), label={[btplabel]below:Phase de production}] {};
  \end{scope}

\end{tikzpicture}
}
  \caption{Cycle de réponse à un appel d'offres}
  \label{fig:cycle-ao}
\end{figure}

Chaque étape mobilise des compétences spécifiques et génère une production
documentaire conséquente. La phase de rédaction, particulièrement critique, peut
s'étendre sur plusieurs semaines et impliquer la mobilisation simultanée
d'experts techniques, juridiques et financiers. Dans le contexte camerounais,
les délais de remise oscillent typiquement entre quinze et quarante-cinq jours,
ce qui concentre la charge sur une période courte et rend le processus
vulnérable à l'omission.

L'observation du processus permet d'identifier trois points de friction
documentaires récurrents~:

\begin{description}
  \item[La relecture répétée du dossier entrant.] Les mêmes pièces sont
    parcourues plusieurs fois par des intervenants différents, chacun y
    cherchant les éléments relevant de son domaine. L'information est extraite
    autant de fois qu'il y a d'intervenants.
  \item[La reconstitution du patrimoine.] Les références de projets, les
    curriculums d'experts et les méthodologies existent déjà, mais leur
    localisation dans l'arborescence des dossiers antérieurs relève de la
    mémoire individuelle.
  \item[Le contrôle de complétude en fin de course.] La vérification de la
    présence effective de chaque pièce exigée intervient tardivement, sous
    contrainte de temps, alors qu'elle conditionne la recevabilité de l'offre.
\end{description}

\subsection{Les enjeux de la capitalisation des connaissances}

Les entreprises du BTP accumulent au fil des années un patrimoine documentaire
considérable~:

\begin{itemize}
  \item \btpterm{Dossiers de réponse}~: des dizaines à des centaines d'offres
        déposées, contenant des méthodologies rédigées, des références de
        projets et des curriculums d'experts~;
  \item \btpterm{Documents techniques}~: études, plans, notes de calcul,
        rapports de mission~;
  \item \btpterm{Documents administratifs}~: attestations, certifications,
        agréments, avec leurs dates de validité~;
  \item \btpterm{Retours d'expérience}~: bilans de chantier, analyses
        \anglais{post mortem}, retours de clients.
\end{itemize}

Cette mémoire, souvent non structurée et disséminée dans des dossiers physiques
ou des répertoires électroniques, constitue un actif stratégique
sous-exploité. La problématique centrale est donc de transformer ce patrimoine
documentaire en une mémoire d'entreprise véritablement exploitable, c'est-à-dire
interrogeable par le sens autant que par les mots, et mobilisable pour produire
de nouveaux contenus.

\begin{btpbox}[Ce que « capitaliser » signifie ici]
Capitaliser ne consiste pas à archiver. Un archivage rend un document
retrouvable~; une capitalisation rend son contenu réutilisable. La distinction
porte sur trois capacités~: retrouver un passage pertinent sans connaître les
mots exacts qu'il emploie, rattacher ce passage à son contexte structurel
(article, section, pièce), et le mobiliser comme source citable dans un nouveau
document.
\end{btpbox}

% =============================================================================
% 1.3 Gestion documentaire et systèmes d'information dans le BTP
% =============================================================================
\section{Gestion documentaire et systèmes d'information dans le BTP}
\label{sec:ch1-ged}

\subsection{Évolution des systèmes de gestion documentaire}

Les systèmes de gestion électronique des documents ont connu une évolution
significative depuis les années 1990. Le tableau~\ref{tab:generations-ged}
en retrace les grandes générations.

\begin{btptable}{Générations de systèmes de gestion documentaire}
\label{tab:generations-ged}
\begin{tabularx}{\linewidth}{@{}L{2.5cm} L{1.9cm} Y Y@{}}
\toprule
\rowcolor{btppale}
\btpth{Génération} & \btpth{Période} & \btpth{Caractéristiques} & \btpth{Limites} \\
\midrule
Première & 1990--2000 & Archivage électronique, indexation manuelle. & Recherche limitée, aucune analyse sémantique. \\
Deuxième & 2000--2010 & Circuits de validation, métadonnées structurées. & Rigidité des schémas, coût d'implémentation. \\
Troisième & 2010--2020 & Hébergement infonuagique, collaboration, recherche plein texte. & Exploitation superficielle du contenu. \\
Quatrième & Depuis 2020 & Analyse sémantique, extraction et génération assistées. & Maturité en cours de consolidation. \\
\bottomrule
\end{tabularx}
\end{btptable}

La quatrième génération, qui intègre des capacités d'intelligence artificielle,
est particulièrement prometteuse pour le domaine des appels d'offres, car elle
permet non seulement de rechercher des documents, mais aussi d'en extraire des
connaissances et de générer de nouveaux contenus. C'est dans cette génération
que se situe notre contribution.

\subsection{Standards et normes applicables}

Plusieurs référentiels encadrent la gestion documentaire dans le secteur~:

\begin{itemize}
  \item \btpterm{ISO 15489} — gestion des documents d'activité~: définit les
        principes de création, de capture et de gestion des documents engageants~;
  \item \btpterm{ISO 30301} — systèmes de management des documents d'activité~:
        cadre organisationnel et exigences d'audit~;
  \item \btpterm{Modélisation des données du bâtiment} (\anglais{Building
        Information Modeling})~: structuration des données techniques d'un
        ouvrage sur son cycle de vie.
\end{itemize}

Ces référentiels partagent une exigence directement pertinente pour notre
travail~: la \btpterm{traçabilité}. Un document engageant doit pouvoir être
rattaché à son auteur, à sa version et à ses sources. Cette exigence normative
converge avec la contrainte technique que nous imposerons au système de
génération~: toute affirmation produite doit être citable.

\subsection{Les limites des solutions actuelles}

Les solutions de gestion documentaire et de gestion des appels d'offres
disponibles sur le marché présentent quatre limitations structurelles.

\begin{enumerate}
  \item \btpterm{Recherche sémantique limitée.} La plupart des systèmes
        fonctionnent sur la base de mots-clés, sans compréhension du sens ni du
        contexte. Une recherche de « pénalités de retard » ne retrouve pas un
        passage évoquant les « sanctions pour dépassement de délai ».

  \item \btpterm{Absence d'exploitation des relations.} Les liens entre
        documents, projets, clients, certifications et compétences ne sont pas
        explicitement modélisés. L'information existe, mais les chemins qui la
        relient restent implicites.

  \item \btpterm{Capitalisation passive.} Les connaissances restent enfermées
        dans les documents, sans extraction ni réutilisation active. Le système
        restitue des fichiers, non des réponses.

  \item \btpterm{Absence de génération.} Les systèmes ne produisent pas de
        nouveaux contenus à partir du patrimoine existant. La rédaction demeure
        intégralement humaine, y compris pour des sections largement
        réutilisables d'une offre à l'autre.
\end{enumerate}

\begin{btpwarn}[Ce que l'automatisation ne doit pas faire disparaître]
Ces limites pourraient laisser croire qu'une automatisation intégrale est
souhaitable. Ce serait une erreur d'analyse. Le caractère engageant de l'offre
impose que la responsabilité reste humaine. L'objectif n'est donc pas de
supprimer l'intervention de l'ingénieur, mais de déplacer son effort~: de la
recherche et de la transcription vers la vérification et l'arbitrage.
\end{btpwarn}

% =============================================================================
% 1.4 Intelligence artificielle et traitement du langage naturel
% =============================================================================
\section{Intelligence artificielle et traitement automatique du langage}
\label{sec:ch1-tal}

\subsection{Fondements du traitement automatique du langage}

Le traitement automatique du langage naturel est une discipline de
l'intelligence artificielle qui vise à faire analyser, comprendre et produire du
langage humain par des machines. Ses origines remontent aux années 1950, mais
c'est au cours de la dernière décennie qu'elle a connu des avancées décisives
grâce à l'apprentissage profond.

\subsubsection{Approches traditionnelles}

Les premières approches reposaient sur des règles linguistiques et des méthodes
statistiques~:

\begin{itemize}
  \item \btpterm{Analyse morphosyntaxique}~: segmentation en phrases,
        tokenisation, étiquetage morphologique~;
  \item \btpterm{Grammaires formelles}~: analyse syntaxique fondée sur des
        règles de production~;
  \item \btpterm{Modèles de Markov cachés}~: modélisation probabiliste de
        séquences d'états latents~;
  \item \btpterm{Pondération TF-IDF}~: mesure de l'importance d'un terme dans un
        document relativement à une collection.
\end{itemize}

Ces approches, efficaces sur des domaines restreints et des tâches bien
délimitées, peinaient à capturer la complexité et la polysémie du langage
naturel. Elles conservent néanmoins une pertinence pratique~: la recherche
lexicale, dont nous verrons l'importance en
section~\ref{sec:ch1-recherche}, en descend directement.

\subsubsection{L'ère de l'apprentissage profond}

L'avènement des réseaux de neurones profonds, à partir de 2012, a transformé le
domaine~:

\begin{itemize}
  \item \btpterm{Word2Vec} (Mikolov et coll., 2013)~: représentation
        vectorielle des mots capturant des régularités sémantiques par
        arithmétique vectorielle~;
  \item \btpterm{GloVe} (Pennington et coll., 2014)~: vecteurs de mots dérivés
        des statistiques globales de co-occurrence~;
  \item \btpterm{Réseaux récurrents et à mémoire longue} (Hochreiter \&
        Schmidhuber, 1997)~: traitement de séquences avec état interne~;
  \item \btpterm{Mécanisme d'attention et architecture Transformer} (Vaswani et
        coll., 2017)~: abandon de la récurrence au profit d'une mise en relation
        directe de toutes les positions d'une séquence.
\end{itemize}

\subsection{Les modèles de langage}

\subsubsection{Principe}

Un modèle de langage est un système capable d'estimer la probabilité d'une
séquence de mots et, par échantillonnage successif, de produire du texte. Les
modèles modernes, fondés sur l'architecture Transformer, sont entraînés sur des
corpus massifs et développent des représentations fines des régularités du
langage.

\begin{btpdef}[Grand modèle de langage]
Un grand modèle de langage est un modèle de langage neuronal dont le nombre de
paramètres et la taille du corpus d'entraînement sont suffisamment élevés pour
qu'apparaissent des capacités non explicitement entraînées, notamment le suivi
d'instructions et l'apprentissage à partir de quelques exemples fournis dans
l'invite.
\end{btpdef}

\subsubsection{Évolution chronologique}

\begin{btptable}{Jalons de l'évolution des modèles de langage}
\label{tab:jalons-llm}
\begin{tabularx}{\linewidth}{@{}C{1.3cm} L{2.3cm} Y L{2.5cm}@{}}
\toprule
\rowcolor{btppale}
\btpth{Année} & \btpth{Modèle} & \btpth{Apport principal} & \btpth{Ordre de grandeur} \\
\midrule
2018 & BERT & Contextualisation bidirectionnelle par masquage. & $3{,}4 \times 10^{8}$ \\
2019 & GPT-2 & Génération de texte long cohérent. & $1{,}5 \times 10^{9}$ \\
2020 & GPT-3 & Apprentissage à partir de quelques exemples. & $1{,}75 \times 10^{11}$ \\
2022 & Modèles alignés par instruction & Suivi fiable de consignes. & Comparable \\
2023 & Modèles multimodaux & Raisonnement étendu, entrée non textuelle. & Non publié \\
2023 & Modèles ouverts (famille Llama) & Diffusion des poids, exécution locale. & $7$ à $70 \times 10^{9}$ \\
2024 & Mélanges d'experts & Capacité accrue à coût d'inférence contenu. & Activation partielle \\
\bottomrule
\end{tabularx}
\end{btptable}

\subsubsection{Modèles multilingues}

Les modèles multilingues sont déterminants dans un contexte francophone comme
celui des appels d'offres camerounais. Les travaux de référence incluent
\btpterm{XLM-RoBERTa} (Conneau et coll., 2019), entraîné sur une centaine de
langues, \btpterm{mT5} (Xue et coll., 2020), variante multilingue d'un modèle
encodeur-décodeur, et les modèles de plongement multilingues à représentation
mixte dense et creuse, dont \btpterm{BGE-M3} (2023) est représentatif.

\subsubsection{Modèles spécialisés}

Le domaine des appels d'offres nécessite une compréhension précise de textes
juridiques, techniques et financiers. Des modèles adaptés à ces domaines ont
émergé, dont \btpterm{LegalBERT} pour les textes juridiques et
\btpterm{FinBERT} pour les textes financiers. Leur intérêt réside moins dans un
gain brut de performance que dans une meilleure calibration du vocabulaire
spécialisé, ce qui réduit les confusions terminologiques.

\begin{btpwarn}[Limite pour le contexte francophone africain]
Ces modèles spécialisés sont majoritairement entraînés sur des corpus
anglo-saxons, dont le droit et les pratiques contractuelles diffèrent
sensiblement du droit camerounais des marchés publics. Leur transposition n'est
donc pas immédiate, ce qui renforce l'intérêt d'une approche par récupération
dans un corpus documentaire local plutôt que par mémorisation paramétrique.
\end{btpwarn}

\subsection{Les représentations vectorielles}

\subsubsection{Définition et principes}

Une représentation vectorielle, ou plongement, est la projection d'un texte —
mot, phrase ou document — dans un espace vectoriel de dimension fixe. La
propriété recherchée est que la proximité géométrique traduise la proximité
sémantique~: deux textes de sens voisin sont représentés par des vecteurs
proches, ce qui permet de mesurer une similarité par simple calcul de distance.

\subsubsection{Évolution des approches}

La figure~\ref{fig:evolution-representations} retrace les cinq générations de
représentations textuelles.

\begin{figure}[htbp]
  \centering
  \resizebox{\linewidth}{!}{% Évolution des représentations vectorielles de texte
\begin{tikzpicture}[node distance=4mm]

  \node[btpboxplain, text width=11.2cm, align=left] (r1)
    {\textbf{1. Représentations creuses} \hfill \btplabel{TF-IDF, BM25}\\[1pt]
     \btplabel{\itshape Termes exacts, aucune compréhension sémantique}};
  \node[btpboxplain, text width=11.2cm, align=left, below=of r1] (r2)
    {\textbf{2. Plongements statiques} \hfill \btplabel{Word2Vec, GloVe}\\[1pt]
     \btplabel{\itshape Sémantique lexicale, mais insensible au contexte}};
  \node[btpbox, text width=11.2cm, align=left, below=of r2] (r3)
    {\textbf{3. Représentations contextuelles} \hfill \btplabel{ELMo, BERT}\\[1pt]
     \btplabel{\itshape Contexte bidirectionnel, gestion de la polysémie}};
  \node[btpboxalt, text width=11.2cm, align=left, below=of r3] (r4)
    {\textbf{4. Modèles mixtes dense + creux} \hfill \btplabel{SPLADE, BGE-M3}\\[1pt]
     \btplabel{\itshape Recherche hybride nativement supportée}};
  \node[btpboxalt, text width=11.2cm, align=left, below=of r4] (r5)
    {\textbf{5. Représentations multimodales} \hfill \btplabel{CLIP, ImageBind}\\[1pt]
     \btplabel{\itshape Texte, images, plans et schémas dans un espace commun}};

  \draw[btparrow, line width=1.6pt]
    ($(r1.west)+(-0.45,0.15)$) -- ($(r5.west)+(-0.45,-0.15)$)
    node[midway, rotate=90, anchor=south, btplabel, yshift=1pt]
    {Enrichissement sémantique};

\end{tikzpicture}
}
  \caption{Évolution des représentations vectorielles de texte}
  \label{fig:evolution-representations}
\end{figure}

\subsubsection{Représentations denses et creuses}

\begin{btptable}{Comparaison des représentations denses et creuses}
\label{tab:dense-creux}
\begin{tabularx}{\linewidth}{@{}L{1.7cm} L{3.4cm} Y Y@{}}
\toprule
\rowcolor{btppale}
\btpth{Type} & \btpth{Caractéristiques} & \btpth{Avantages} & \btpth{Inconvénients} \\
\midrule
Dense & Vecteur continu de dimension fixe (768 ou 1024 par exemple). & Capture la sémantique, tolère la reformulation, généralise bien. & Peine sur les termes rares, les identifiants et les références normatives. \\
Creux & Vecteur de la taille du vocabulaire, majoritairement nul. & Capture les termes exacts, directement interprétable. & Ne relie pas les synonymes ni les paraphrases. \\
\bottomrule
\end{tabularx}
\end{btptable}

Cette complémentarité n'est pas une nuance de second ordre~: elle fonde le choix
d'une recherche hybride, développé en section~\ref{sec:ch1-recherche}.

\subsection{Applications des grands modèles de langage}

Les capacités des grands modèles de langage recoupent notre problématique sur
plusieurs plans~:

\begin{enumerate}
  \item \btpterm{Génération de textes}~: rédaction de mémoires techniques, de
        notes méthodologiques et de sections administratives~;
  \item \btpterm{Extraction d'informations}~: identification des exigences,
        des entités nommées et des relations dans un dossier entrant~;
  \item \btpterm{Résumé}~: synthèse de pièces juridiques ou techniques longues~;
  \item \btpterm{Question-réponse documentaire}~: réponse à une interrogation
        précise à partir d'un corpus déterminé~;
  \item \btpterm{Classification}~: catégorisation automatique des pièces d'un
        dossier selon leur nature~;
  \item \btpterm{Reformulation de requêtes}~: enrichissement d'une question
        utilisateur avant recherche.
\end{enumerate}

Ces capacités partagent toutefois une faiblesse commune, dont la maîtrise
constitue le c\oe ur de notre travail~: en l'absence de sources fournies, un
modèle de langage produit un texte plausible, non un texte vrai. C'est cette
faiblesse que l'architecture examinée à la section~\ref{sec:ch1-rag} entend
corriger.

% =============================================================================
% 1.5 Recherche d'information et modèles de classement
% =============================================================================
\section{Recherche d'information et modèles de classement}
\label{sec:ch1-recherche}

\subsection{Fondements de la recherche d'information}

La recherche d'information est une discipline fondatrice du traitement
automatisé des documents, dont les origines remontent à la vision du
\anglais{Memex} formulée par Vannevar Bush en 1945 et aux premiers systèmes
documentaires des années 1960. Elle répond à une question unique~: étant donnée
une requête, comment identifier et ordonner les documents les plus pertinents
d'une collection~?

Un système de recherche s'articule en quatre opérations~: l'\btpterm{indexation}
transforme les documents en une représentation interrogeable~;
l'\btpterm{interrogation} transforme la requête dans le même espace~; le
\btpterm{classement} ordonne les candidats par pertinence estimée~;
l'\btpterm{évaluation} mesure la qualité de ce classement au regard d'un
jugement de référence.

\subsubsection{Métriques d'évaluation}

\begin{btptable}{Métriques d'évaluation de la recherche d'information}
\label{tab:metriques-ri}
\begin{tabularx}{\linewidth}{@{}L{2.8cm} Y L{3.1cm}@{}}
\toprule
\rowcolor{btppale}
\btpth{Métrique} & \btpth{Définition} & \btpth{Ce qu'elle mesure} \\
\midrule
Rappel à $K$ & Proportion des documents pertinents présents dans les $K$ premiers résultats. & Couverture~: le système laisse-t-il échapper des sources~? \\
Précision à $K$ & Proportion de documents pertinents parmi les $K$ premiers résultats. & Pureté~: le contexte transmis au modèle est-il exploitable~? \\
Rang réciproque moyen & Moyenne de l'inverse du rang du premier document pertinent. & Qualité de la tête de classement. \\
Gain cumulé normalisé & Gain cumulé escompté, pondéré par la position et normalisé. & Qualité du classement avec gradation de pertinence. \\
\bottomrule
\end{tabularx}
\end{btptable}

Pour notre application, la distinction entre rappel et précision est
opérationnelle. Le rappel gouverne l'exhaustivité~: un passage pertinent non
récupéré ne pourra jamais être cité. La précision gouverne la qualité du
contexte transmis au générateur~: des extraits hors sujet dégradent la
génération et augmentent le coût d'inférence.

\subsection{Modèles lexicaux}

\subsubsection{Le modèle vectoriel et la pondération TF-IDF}

Formalisé par Salton, le modèle vectoriel représente chaque document et chaque
requête comme un vecteur dont les dimensions sont les termes du vocabulaire. La
pertinence est estimée par la similarité cosinus entre les deux vecteurs. La
pondération TF-IDF combine deux intuitions~: un terme fréquent dans un document
y est important (\anglais{term frequency}), mais un terme fréquent dans toute la
collection est peu discriminant (\anglais{inverse document frequency}).

\subsubsection{Le modèle probabiliste BM25}

Développé par Robertson et ses collaborateurs à partir des travaux fondateurs de
Robertson et Sparck Jones (1976), le modèle BM25 s'inscrit dans la théorie
probabiliste de la recherche d'information. Il demeure aujourd'hui l'un des
modèles lexicaux les plus performants et sert de référence dans la plupart des
campagnes d'évaluation.

Le score attribué à un document $D$ pour une requête $Q$ s'écrit~:

\begin{equation}
  \mathrm{score}(D, Q) = \sum_{q_i \in Q}
    \mathrm{IDF}(q_i) \cdot
    \frac{f(q_i, D)\,(k_1 + 1)}
         {f(q_i, D) + k_1 \left(1 - b + b\,\dfrac{|D|}{\mathrm{avgdl}}\right)}
  \label{eq:bm25}
\end{equation}

où $f(q_i, D)$ désigne la fréquence du terme $q_i$ dans $D$, $|D|$ la longueur
du document, $\mathrm{avgdl}$ la longueur moyenne des documents de la
collection, $k_1$ un paramètre de saturation de la fréquence (usuellement entre
$1{,}2$ et $2{,}0$) et $b$ un paramètre de normalisation par la longueur
(usuellement $0{,}75$).

Trois propriétés de l'équation~\eqref{eq:bm25} méritent d'être soulignées, car
elles expliquent sa robustesse~:

\begin{itemize}
  \item \btpterm{Saturation de la fréquence.} Le terme en $f/(f + k_1 \cdots)$
        est borné~: la centième occurrence d'un mot apporte beaucoup moins que
        la deuxième. Le modèle résiste ainsi à la répétition artificielle.
  \item \btpterm{Normalisation souple par la longueur.} Le paramètre $b$ permet
        de moduler la pénalisation des documents longs, au lieu de l'imposer.
  \item \btpterm{Pondération par la rareté.} Le facteur $\mathrm{IDF}$ accorde
        un poids déterminant aux termes rares, ce qui est précisément le
        comportement attendu pour une référence normative ou un numéro d'article.
\end{itemize}

\subsubsection{Limites des approches lexicales}

Malgré leur efficacité, les approches lexicales présentent des limites
structurelles~: la \btpterm{synonymie} n'est pas traitée, « voiture » et
« automobile » restant étrangers l'un à l'autre~; la \btpterm{polysémie} n'est
pas résolue, un même terme recouvrant des sens distincts~; enfin la
\btpterm{sensibilité à la formulation} fait qu'une même question posée
différemment produit des résultats différents.

\subsection{Modèles sémantiques}

L'émergence des représentations vectorielles continues a permis de dépasser ces
limites en capturant la sémantique plutôt que la forme. Les plongements
statiques (Word2Vec, GloVe) ont d'abord établi qu'un mot peut être représenté
par un vecteur dans un espace sémantique. Les plongements contextuels, avec
BERT et ses successeurs, ont ensuite rendu cette représentation dépendante du
contexte d'occurrence, résolvant en grande partie la polysémie.

Pour la recherche documentaire, des modèles de plongement spécialisés ont été
développés, dont le tableau~\ref{tab:modeles-plongement} donne un aperçu.

\begin{btptable}{Familles de modèles de plongement pour la recherche}
\label{tab:modeles-plongement}
\begin{tabularx}{\linewidth}{@{}L{2.6cm} Y L{3.2cm}@{}}
\toprule
\rowcolor{btppale}
\btpth{Famille} & \btpth{Caractéristiques} & \btpth{Intérêt} \\
\midrule
Modèles à représentation mixte & Produisent conjointement un vecteur dense et une représentation creuse. & Recherche hybride sans double indexation. \\
Modèles multilingues alignés & Espace sémantique partagé entre langues. & Corpus mêlant français et anglais. \\
Modèles à instruction & La requête et le document sont encodés avec des préfixes distincts. & Meilleure asymétrie requête-document. \\
\bottomrule
\end{tabularx}
\end{btptable}

\subsection{Recherche hybride et fusion de classements}

\subsubsection{Motivation}

Les approches lexicales et sémantiques sont complémentaires plutôt que
concurrentes. La première excelle sur les termes exacts, les identifiants et les
références normatives~; la seconde sur les concepts, les synonymes et les
reformulations. Le tableau~\ref{tab:complementarite} illustre cette
complémentarité sur des requêtes typiques du domaine.

\begin{btptable}{Complémentarité des approches sur des requêtes d'appel d'offres}
\label{tab:complementarite}
\begin{tabularx}{\linewidth}{@{}Y C{2.6cm} C{2.6cm}@{}}
\toprule
\rowcolor{btppale}
\btpth{Type de requête} & \btpth{Lexical (BM25)} & \btpth{Sémantique (dense)} \\
\midrule
« ISO 9001 »                        & \cotefort  & \cotemoyen \\
« Article 5.2 du CCAP »             & \cotefort  & \cotefaible \\
« pénalités de retard »             & \cotemoyen & \cotefort \\
« garantie décennale »              & \cotemoyen & \cotefort \\
« expérience en assainissement »    & \cotefaible & \cotefort \\
\bottomrule
\end{tabularx}
\end{btptable}

\subsubsection{La fusion par rang réciproque}

Combiner les résultats de plusieurs systèmes est un problème central de la
recherche d'information, étudié depuis les travaux de Fox sur la fusion de
données et ceux de Lee sur la combinaison d'évidences. Parmi les méthodes
proposées, la \btpterm{fusion par rang réciproque}, introduite par Cormack,
Clarke et Buettcher en 2009, se distingue par sa simplicité et son efficacité
empirique.

Le score fusionné d'un document $D$ s'écrit~:

\begin{equation}
  \mathrm{score}_{\mathrm{RRF}}(D) =
  \sum_{i \in \mathcal{S}} \frac{1}{k + r_i(D)}
  \label{eq:rrf}
\end{equation}

où $\mathcal{S}$ désigne l'ensemble des systèmes fusionnés, $r_i(D)$ le rang de
$D$ dans le classement du système $i$, et $k$ une constante d'amortissement
fixée usuellement à $60$.

\begin{btpbox}[Pourquoi la fusion par rang réciproque fonctionne]
L'équation~\eqref{eq:rrf} n'utilise que les \btpterm{rangs}, jamais les scores.
Cette propriété résout élégamment le problème d'échelle~: un score BM25, non
borné, et une similarité cosinus, comprise entre $-1$ et $1$, ne sont pas
comparables et leur normalisation est arbitraire. La formule est en outre
robuste au bruit, car l'amortissement par $k$ écrase les écarts entre rangs
élevés~: une erreur ponctuelle de classement a peu d'effet sur le résultat
fusionné.
\end{btpbox}

Le tableau~\ref{tab:exemple-rrf} illustre le calcul sur trois documents.

\begin{btptable}{Exemple de calcul de fusion par rang réciproque ($k = 60$)}
\label{tab:exemple-rrf}
\begin{tabularx}{\linewidth}{@{}L{2.2cm} C{2.2cm} C{2.6cm} Y@{}}
\toprule
\rowcolor{btppale}
\btpth{Document} & \btpth{Rang BM25} & \btpth{Rang sémantique} & \btpth{Score fusionné} \\
\midrule
Document A & 1 & 3 & $1/61 + 1/63 = 0{,}03228$ \\
Document B & 2 & 2 & $1/62 + 1/62 = 0{,}03226$ \\
Document C & 5 & 1 & $1/65 + 1/61 = 0{,}03178$ \\
\bottomrule
\end{tabularx}
\end{btptable}

Le document A, bien classé par les deux systèmes sans être premier partout,
devance le document C, premier pour un seul système. La fusion récompense donc
l'accord entre systèmes plutôt que l'excellence isolée.

\subsubsection{Variantes de la fusion}

\begin{btptable}{Familles de méthodes de fusion de classements}
\label{tab:familles-fusion}
\begin{tabularx}{\linewidth}{@{}L{3.1cm} Y Y@{}}
\toprule
\rowcolor{btppale}
\btpth{Méthode} & \btpth{Avantages} & \btpth{Inconvénients} \\
\midrule
Combinaison de scores & Théoriquement optimale si les scores sont calibrés. & Exige une normalisation dont le choix est arbitraire. \\
Comptage de Borda & Simple, fondée sur l'agrégation de préférences. & Perd l'information de score et l'écart entre rangs. \\
Fusion par rang réciproque & Sans normalisation, robuste, performante. & Un paramètre $k$ à régler. \\
Apprentissage du classement & Optimale en présence de données annotées. & Exige un corpus de jugements de pertinence. \\
\bottomrule
\end{tabularx}
\end{btptable}

\subsubsection{Le ré-ordonnancement}

Au-delà de la fusion, une étape supplémentaire affine le classement. Un
\btpterm{modèle croisé} (\anglais{cross-encoder}) encode conjointement la
requête et le document candidat, ce qui lui permet de modéliser leurs
interactions terme à terme. Il produit une estimation de pertinence bien plus
fine qu'une similarité entre vecteurs calculés indépendamment.

Ce gain a un coût~: là où une recherche vectorielle compare un vecteur de
requête à un index, le modèle croisé exige une inférence par couple. Il n'est
donc applicable qu'à un nombre restreint de candidats, ce qui impose
l'architecture en entonnoir de la figure~\ref{fig:pipeline-hybride}.

\begin{figure}[htbp]
  \centering
  \resizebox{\linewidth}{!}{% Pipeline de recherche hybride : BM25 + dense, fusion RRF, ré-ordonnancement
\begin{tikzpicture}[node distance=5mm and 10mm]

  \node[btpboxstrong, text width=2.4cm] (q) {Requête};

  \node[btpbox, text width=3.1cm, above right=2mm and 12mm of q] (bm25)
    {Recherche lexicale\\[-1pt]\btplabel{BM25 — termes exacts}};
  \node[btpbox, text width=3.1cm, below right=2mm and 12mm of q] (dense)
    {Recherche sémantique\\[-1pt]\btplabel{vecteurs denses — concepts}};

  \node[btpboxplain, text width=1.5cm, right=6mm of bm25] (t1) {Top 50};
  \node[btpboxplain, text width=1.5cm, right=6mm of dense] (t2) {Top 50};

  \node[btpboxalt, text width=2.4cm, right=8mm of {$(t1)!0.5!(t2)$}] (rrf)
    {Fusion\\\textbf{RRF}};
  \node[btpbox, text width=2.6cm, right=9mm of rrf] (rerank)
    {Ré-ordonnancement\\[-1pt]\btplabel{modèle croisé}};
  \node[btpboxstrong, text width=2.1cm, right=9mm of rerank] (out)
    {Top 8\\contexte};

  \draw[btparrow] (q) -- (bm25);
  \draw[btparrow] (q) -- (dense);
  \draw[btparrow] (bm25) -- (t1);
  \draw[btparrow] (dense) -- (t2);
  \draw[btparrow] (t1.east) -- (rrf.north west);
  \draw[btparrow] (t2.east) -- (rrf.south west);
  \draw[btparrow] (rrf) -- (rerank);
  \draw[btparrow] (rerank) -- (out);

  \node[btplabel, below=8mm of rrf, text width=3.6cm]
    {Aucune normalisation de score requise};
  \node[btplabel, below=8mm of rerank, text width=3.4cm]
    {Pertinence évaluée par couple (requête, extrait)};

  \begin{scope}[on background layer]
    \node[btpgroup, fit=(bm25)(dense)(t1)(t2),
          label={[btplabel]above:Récupération parallèle}] {};
  \end{scope}

\end{tikzpicture}
}
  \caption{Pipeline de recherche hybride en entonnoir}
  \label{fig:pipeline-hybride}
\end{figure}

\subsection{Application aux appels d'offres}

Les appels d'offres présentent des caractéristiques qui justifient une approche
hybride, résumées au tableau~\ref{tab:ri-ao}.

\begin{btptable}{Caractéristiques des appels d'offres et implications pour la recherche}
\label{tab:ri-ao}
\begin{tabularx}{\linewidth}{@{}L{3.9cm} Y@{}}
\toprule
\rowcolor{btppale}
\btpth{Caractéristique} & \btpth{Implication} \\
\midrule
Références normatives et identifiants & Recherche lexicale indispensable~: un numéro d'article ne se paraphrase pas. \\
Concepts métier exprimés diversement & Recherche sémantique indispensable~: la même exigence se formule de multiples façons. \\
Variantes de désignation des entités & Résolution d'entités nécessaire (section~\ref{sec:ch1-entites}). \\
Exigence de citation des sources & Précision de la récupération critique~: une citation erronée est plus dommageable qu'une abstention. \\
Structure hiérarchique des pièces & Le granule de recherche doit respecter la structure (section~\ref{sec:ch1-documents}). \\
\bottomrule
\end{tabularx}
\end{btptable}

Cette complémentarité justifie l'adoption d'un pipeline hybride combinant
recherche lexicale, recherche sémantique, fusion par rang réciproque et
ré-ordonnancement, dont la mise en \oe uvre est détaillée au
chapitre~\ref{chap:methodologie}.

% =============================================================================
% 1.6 RAG : génération augmentée par récupération
% =============================================================================
\section{Génération augmentée par récupération}
\label{sec:ch1-rag}

\subsection{Définition et principes fondamentaux}

La génération augmentée par récupération (\anglais{Retrieval-Augmented
Generation}, RAG) est une architecture introduite par Lewis et coll.\ (2020)
qui combine un moteur de recherche et un générateur de texte. Plutôt que de
s'en remettre uniquement aux connaissances paramétriques du modèle, le système
récupère des extraits pertinents dans une base externe, puis conditionne la
génération sur ces extraits.

\begin{btpdef}[Génération augmentée par récupération]
On appelle génération augmentée par récupération toute architecture dans
laquelle un modèle de langage produit une réponse \emph{conditionnellement} à
un ensemble d'extraits récupérés au moment de la requête, plutôt qu'à partir
de sa seule mémoire paramétrique.
\end{btpdef}

Cette architecture répond à trois faiblesses des modèles de langage utilisés
seuls~: l'\btpterm{hallucination} (production d'affirmations non fondées),
l'\btpterm{obsolescence} (connaissances figées à la date d'entraînement) et
l'\btpterm{absence de citation} (impossibilité de rattacher une phrase à une
source). Pour un dossier d'appel d'offres, ces trois faiblesses sont
rédhibitoires.

La figure~\ref{fig:principe-rag} en rappelle le principe.

\begin{figure}[htbp]
  \centering
  \resizebox{\linewidth}{!}{% Principe général de la génération augmentée par récupération
\begin{tikzpicture}[node distance=6mm and 11mm]

  \node[btpboxstrong, text width=2.5cm] (q) {Question\\utilisateur};
  \node[btpbox, text width=2.9cm, right=of q] (search)
    {Moteur de\\recherche\\[-1pt]\btplabel{lexical, sémantique, hybride}};
  \node[btpdb, right=of search, minimum width=2.3cm] (kb)
    {Base de\\connaissances};
  \node[btpbox, text width=2.9cm, below=16mm of search] (ctx)
    {Contexte\\[-1pt]\btplabel{extraits pertinents numérotés}};
  \node[btpbox, text width=2.5cm, left=of ctx] (prompt)
    {Invite\\[-1pt]\btplabel{question + contexte + consignes}};
  \node[btpboxalt, text width=2.9cm, right=of ctx] (llm)
    {Modèle de\\langage};
  \node[btpboxstrong, text width=2.7cm, below=14mm of llm] (rep)
    {Réponse\\avec citations};

  \draw[btparrow] (q) -- (search);
  \draw[btparrow] (search) -- (kb);
  \draw[btparrowlight] (kb.south) |- (ctx.east);
  \draw[btparrow] (ctx) -- (prompt);
  \draw[btparrow] (prompt.south) |- ($(rep.west)+(-0.7,0)$) -- (rep.west);
  \draw[btparrow] (ctx) -- (llm);
  \draw[btparrow] (llm) -- (rep);

  \node[btplabel, below=1mm of rep, text width=6cm]
    {Chaque affirmation est rattachable à un extrait source identifié};

\end{tikzpicture}
}
  \caption{Principe de la génération augmentée par récupération}
  \label{fig:principe-rag}
\end{figure}

\subsection{Les composantes d'un système RAG}

\subsubsection{Le récupérateur}

Le récupérateur identifie les extraits pertinents dans la base. Les approches
examinées à la section~\ref{sec:ch1-recherche} s'y appliquent directement,
résumées au tableau~\ref{tab:approches-retriever}.

\begin{btptable}{Approches de récupération dans un système RAG}
\label{tab:approches-retriever}
\begin{tabularx}{\linewidth}{@{}L{2.6cm} Y L{3.4cm}@{}}
\toprule
\rowcolor{btppale}
\btpth{Approche} & \btpth{Principe} & \btpth{Cas d'usage} \\
\midrule
Lexicale & Classement par correspondance de termes (BM25). & Identifiants, articles, normes. \\
Sémantique & Similarité vectorielle dans un espace de plongement. & Concepts, paraphrases. \\
Hybride & Fusion des deux classements (RRF). & Compromis général. \\
Structurée & Filtrage sur métadonnées (type de pièce, organisation). & Isolation et précision. \\
Graphe & Navigation le long de relations explicites. & Questions multi-sauts. \\
\bottomrule
\end{tabularx}
\end{btptable}

\subsubsection{Le ré-ordonnançeur}

Le ré-ordonnançeur affine le classement produit par le récupérateur. Un modèle
croisé évalue la pertinence d'un couple (question, extrait) plutôt qu'une
distance entre vecteurs calculés indépendamment. Ce gain de précision est
décisif~: le contexte transmis au générateur est borné, et chaque extrait
hors sujet occupe une place qui aurait pu revenir à une source utile.

\subsubsection{Le générateur}

Le générateur est un modèle de langage qui reçoit la question, les extraits
récupérés et un ensemble de consignes, puis produit une réponse. Trois
contraintes s'imposent dans le domaine des appels d'offres~:

\begin{itemize}
  \item la réponse doit être \btpterm{citable}~: chaque affirmation factuelle
        est rattachable à un extrait identifié~;
  \item la réponse doit être \btpterm{structurée}~: le livrable n'est pas un
        paragraphe libre, mais une section dont le plan est imposé~;
  \item la réponse doit \btpterm{s'abstenir} plutôt qu'inventer lorsque le
        contexte ne contient pas l'information demandée.
\end{itemize}

\begin{btpnote}[Abstention plutôt qu'invention]
Dans un système conversationnel généraliste, une réponse approximative est
souvent acceptable. Dans un dossier d'offre, elle est plus dangereuse qu'un
silence. L'abstention explicite --- « cette information n'apparaît pas dans
les sources consultées » --- est donc une sortie légitime, et non un échec
du générateur.
\end{btpnote}

\subsection{Les variantes du RAG}

\subsubsection{RAG naïf}

Dans sa forme élémentaire, chaque question déclenche une recherche unique
suivie d'une génération unique, sans mémoire des tours précédents ni
raisonnement multi-étapes. Cette architecture suffit aux questions
factuelles isolées. Elle échoue dès que la requête est anaphorique
(« et la pénalité~? ») ou qu'elle exige plusieurs recherches successives.

\subsubsection{RAG à mémoire}

Une mémoire des interactions précédentes permet de réécrire la requête
courante à la lumière du dialogue. L'anaphore « et la pénalité~? » devient
alors « clause de pénalité du marché dont le CCAP vient d'être discuté ».
Cette variante est utile dans l'espace de travail assisté, où l'opérateur
enchaîne des questions sur un même dossier.

\subsubsection{RAG hiérarchique}

Le découpage parent-enfant introduit deux granularités~: un extrait court,
indexé, sert à la recherche précise~; un extrait plus large, non indexé,
sert de contexte au générateur. Cette stratégie, détaillée au
chapitre~\ref{chap:methodologie}, répond au dilemme classique de la taille
d'extrait~: trop court, le contexte manque~; trop long, la recherche perd
en précision.

\subsubsection{RAG agentique}

Le modèle devient alors un agent qui décide des actions à entreprendre ---
rechercher, évaluer, rechercher à nouveau --- avant de répondre. Cette
variante est pertinente pour les questions complexes à plusieurs sauts. Elle
introduit toutefois une indétermination du chemin d'exécution, difficilement
compatible avec l'exigence de traçabilité d'un dossier engageant. Nous y
revenons à la section~\ref{sec:ch1-workflows}.

\subsubsection{RAG sur graphe}

Une variante récente, popularisée notamment par les travaux de Microsoft
Research (Edge et coll., 2024), utilise un graphe de connaissances comme
structure de récupération~: le document est d'abord projeté en entités et
relations, la recherche navigue dans ce graphe, et le contexte assemblé
reflète des chemins plutôt que des extraits isolés. Les avantages ---
raisonnement multi-sauts, réponses plus relationnelles --- s'accompagnent
d'un coût d'indexation élevé et d'une dépendance à la qualité de
l'extraction. Nous la situons comme une perspective de seconde phase.

\subsection{Évaluation des systèmes RAG}

L'évaluation d'un système RAG est double, car une défaillance peut naître
aussi bien de la récupération que de la génération.

\begin{btptable}{Métriques d'évaluation d'un système RAG}
\label{tab:metriques-rag}
\begin{tabularx}{\linewidth}{@{}L{3.2cm} Y L{3.4cm}@{}}
\toprule
\rowcolor{btppale}
\btpth{Métrique} & \btpth{Définition} & \btpth{Cible} \\
\midrule
Rappel à $K$ & Proportion des extraits pertinents présents dans les $K$ premiers. & Récupération \\
Précision à $K$ & Proportion d'extraits pertinents parmi les $K$ premiers. & Récupération \\
Fidélité & Degré auquel la réponse s'appuie exclusivement sur le contexte fourni. & Génération \\
Ancrage & Proportion d'affirmations rattachables à une source vérifiable. & Génération \\
Exhaustivité & Couverture des exigences ou des aspects de la question. & Génération \\
Exactitude & Conformité factuelle au corpus et au domaine. & Génération \\
\bottomrule
\end{tabularx}
\end{btptable}

Des protocoles automatisés, dont RAGAS (Es et coll., 2024), proposent
d'estimer la fidélité et l'ancrage par un modèle juge. Ils accélèrent
l'itération mais ne dispensent pas d'une évaluation par des experts du
domaine, seule à même de juger la recevabilité d'une section d'offre.

% =============================================================================
% 1.7 Bases de données vectorielles et recherche sémantique
% =============================================================================
\section{Bases de données vectorielles et recherche sémantique}
\label{sec:ch1-vecteurs}

\subsection{Principes}

Une base de données vectorielle est un magasin conçu pour stocker des vecteurs
de grande dimension et les interroger par similarité. Elle constitue le
c\oe ur matériel de la recherche sémantique dans un système RAG.

La figure~\ref{fig:base-vectorielle} en rappelle le fonctionnement en deux
temps~: l'indexation, qui projette les documents dans l'espace vectoriel, et
la recherche, qui projette la question dans le même espace puis identifie
les plus proches voisins.

\begin{figure}[htbp]
  \centering
  \resizebox{\linewidth}{!}{% Fonctionnement d'une base de données vectorielle
\begin{tikzpicture}[node distance=5mm and 9mm]

  % --- Indexation ---
  \node[btptitleband, text width=11.5cm] (b1) {Indexation};
  \node[btpbox, text width=2.3cm, below left=6mm and 3.2cm of b1.south] (doc) {Documents};
  \node[btpbox, text width=2.5cm, right=of doc] (emb1) {Modèle de\\plongement};
  \node[btpbox, text width=2.5cm, right=of emb1] (vec1) {Vecteurs\\$\mathbb{R}^{d}$};
  \node[btpdb, right=of vec1, minimum width=2.2cm] (store) {Index\\vectoriel};

  \draw[btparrow] (doc) -- (emb1);
  \draw[btparrow] (emb1) -- (vec1);
  \draw[btparrow] (vec1) -- (store);

  % --- Recherche ---
  \node[btptitleband, text width=11.5cm, below=12mm of doc.south west, anchor=north west] (b2) {Recherche};
  \node[btpboxalt, text width=2.3cm, below=6mm of b2.south west, anchor=north west] (query) {Question};
  \node[btpboxalt, text width=2.5cm, right=of query] (emb2) {Modèle de\\plongement};
  \node[btpboxalt, text width=2.9cm, right=of emb2] (knn) {Plus proches\\voisins\\[-1pt]\btplabel{HNSW, IVF, PQ}};
  \node[btpboxstrong, text width=2.4cm, right=of knn] (rank) {Classement par\\similarité};

  \draw[btparrow] (query) -- (emb2);
  \draw[btparrow] (emb2) -- (knn);
  \draw[btparrow] (knn) -- (rank);
  \draw[btparrowlight] (store.south) -- ++(0,-0.5) -| (knn.north);

  \node[btplabel, below=4mm of knn, text width=6.5cm]
    {Distance cosinus ou euclidienne, calculée de manière approchée
     pour rester sous-linéaire};

\end{tikzpicture}
}
  \caption{Fonctionnement d'une base de données vectorielle}
  \label{fig:base-vectorielle}
\end{figure}

Deux distances dominent la pratique. La \btpterm{similarité cosinus} mesure
l'angle entre deux vecteurs et est invariante à leur norme, ce qui convient
aux plongements dont l'amplitude n'est pas sémantiquement significative. La
\btpterm{distance euclidienne} mesure l'écart géométrique et suppose des
vecteurs comparables en norme. Le choix est dicté par le modèle de
plongement, non par une préférence architecturale.

\subsection{Indexation approximative}

Une recherche exacte des plus proches voisins est linéaire en la taille de
la collection. Pour des corpus de dizaines ou de centaines de milliers
d'extraits, des index approximatifs rendent la recherche sous-linéaire, au
prix d'un rappel légèrement dégradé.

\begin{btptable}{Familles d'index approximatifs}
\label{tab:index-ann}
\begin{tabularx}{\linewidth}{@{}L{2.6cm} Y C{2.2cm} C{2.2cm}@{}}
\toprule
\rowcolor{btppale}
\btpth{Algorithme} & \btpth{Principe} & \btpth{Vitesse} & \btpth{Fidélité} \\
\midrule
HNSW & Graphe hiérarchique de voisinage navigable. & Très élevée & Élevée \\
IVF & Partitionnement en cellules (fichier inversé). & Élevée & Bonne \\
PQ & Quantification produit, compression des vecteurs. & Élevée & Bonne \\
HNSW + IVF & Combinaison graphe et partitionnement. & Très élevée & Très élevée \\
\bottomrule
\end{tabularx}
\end{btptable}

L'algorithme HNSW (Malkov et Yashunin, 2018) est devenu une référence~: il
construit un graphe multi-niveaux dans lequel la recherche descend
progressivement vers le voisinage local de la requête. Sa fidélité élevée et
sa latence prévisible en font un candidat naturel pour une application
interactive.

\subsection{Propriétés recherchées pour les appels d'offres}

Au-delà de la similarité brute, trois propriétés déterminent l'adéquation
d'un magasin vectoriel à notre problématique.

\begin{enumerate}
  \item \btpterm{Filtrage par métadonnées.} La recherche doit pouvoir être
        restreinte à une organisation (cloisonnement), à un type de pièce
        (CCTP, CCAP, référence) et à un rôle (patrimoine ou dossier courant).
        Un index qui ne supporte pas ce filtrage \emph{avant} le calcul de
        similarité oblige à post-filtrer, ce qui dégrade le rappel.
  \item \btpterm{Représentation mixte.} Un magasin capable de stocker
        conjointement un vecteur dense et une représentation creuse permet
        d'exécuter la recherche hybride sans maintenir deux index distincts.
  \item \btpterm{Identifiants stables.} Chaque vecteur doit rester rattachable
        à l'extrait source, à la page, à la pièce et à la version du
        document. Sans cette stabilité, la citation devient impossible.
\end{enumerate}

Le marché propose aujourd'hui plusieurs solutions, ouvertes ou gérées, dont
certaines exposent nativement une recherche hybride ou un modèle de graphe.
Le choix d'un produit particulier relève de l'implémentation
(chapitre~\ref{chap:resultats})~; ce qui importe ici est le cahier des charges
fonctionnel~: filtrage, mixte dense-creux, identifiants citables.

% =============================================================================
% 1.8 Traitement des documents structurés
% =============================================================================
\section{Traitement des documents structurés}
\label{sec:ch1-documents}

La qualité d'un système RAG est bornée par la qualité de ce qu'il ingère.
Un extrait mal segmenté, un tableau aplati en flux de caractères ou un
article dont le numéro a disparu produisent, en aval, une recherche
imprécise et une génération infidèle. Le traitement des documents
structurés n'est donc pas un préambule technique~: c'est un sujet
d'état de l'art, car il détermine ce qui est ensuite recherchable et
citable.

\subsection{De l'OCR à la compréhension de documents}

La reconnaissance optique de caractères a connu une évolution qui dépasse
largement la transcription de glyphes. Le tableau~\ref{tab:evolution-ocr}
en retrace les générations.

\begin{btptable}{Évolution de la reconnaissance de documents}
\label{tab:evolution-ocr}
\begin{tabularx}{\linewidth}{@{}L{2.6cm} Y Y@{}}
\toprule
\rowcolor{btppale}
\btpth{Période} & \btpth{Approche} & \btpth{Limite principale} \\
\midrule
1960--1990 & Reconnaissance matricielle de caractères imprimés. & Polices variables, manuscrit. \\
1990--2005 & OCR statistique (famille Tesseract). & Perte de la mise en page. \\
2005--2015 & OCR et détection de blocs, colonnes, orientations. & Fragilité sur les documents complexes. \\
2015--2020 & Extraction dédiée de tableaux. & Pas de compréhension sémantique. \\
Depuis 2020 & Compréhension de documents (texte et mise en page conjoints). & Maturité encore en consolidation. \\
\bottomrule
\end{tabularx}
\end{btptable}

La rupture conceptuelle est celle de la \btpterm{compréhension de
documents} (\anglais{Document Understanding}). Là où l'OCR traditionnel
produit un flux linéaire de texte, cette approche vise à préserver la
structure hiérarchique~: titres, sous-titres, articles numérotés, tableaux,
notes de bas de page, annexes. Elle s'appuie sur des modèles de fondation
entraînés conjointement sur le texte et la géométrie de la page
(LayoutLMv3, DocFormer et successeurs). Des plateformes comme Docling,
Azure Document Intelligence ou AWS Textract en sont des illustrations
industrielles.

\begin{btpbox}[Pourquoi la structure est non négociable]
Un CCAP est organisé en articles~; un CCTP suit une hiérarchie de sections~;
un DPGF est un tableau. Perdre cette structure lors de l'extraction
dégrade trois capacités à la fois~: la recherche (un article entier est
plus pertinent qu'un fragment), la citation (la référence à un article
précis est exigée) et la génération (le plan de la réponse s'appuie sur
celui du dossier entrant).
\end{btpbox}

\subsection{Analyse de mise en page et reconnaissance de tableaux}

Les tableaux constituent une composante critique des dossiers d'appels
d'offres. Un bordereau de prix ou une décomposition du prix global
contiennent des informations financières dont l'altération compromet
toute génération ultérieure de sections chiffrées.

Trois défis se distinguent~:

\begin{itemize}
  \item \btpterm{Reconnaissance des cellules}~: identifier les en-têtes, les
        lignes, les colonnes et les cellules fusionnées~;
  \item \btpterm{Préservation des hiérarchies}~: certains tableaux
        embarquent des sous-totaux ou des regroupements par lot~;
  \item \btpterm{OCR spécifique}~: les tableaux emploient souvent des
        polices plus petites et moins régulières que le corps du texte.
\end{itemize}

\begin{btptable}{Approches de reconnaissance de tableaux}
\label{tab:tableaux}
\begin{tabularx}{\linewidth}{@{}L{3.4cm} Y Y@{}}
\toprule
\rowcolor{btppale}
\btpth{Approche} & \btpth{Avantages} & \btpth{Limites} \\
\midrule
Règles heuristiques (lignes, colonnes) & Rapide, simple. & Fragile dès que le quadrillage est irrégulier. \\
Détection par réseau convolutif & Plus robuste aux variations de mise en page. & Exige des données d'entraînement. \\
Modèles de fondation multimodaux & Très robuste, comprend le rôle des cellules. & Coût computationnel élevé. \\
\bottomrule
\end{tabularx}
\end{btptable}

\subsection{Enjeux pour les appels d'offres}

Les dossiers d'appels d'offres camerounais présentent des caractéristiques
qui aggravent ces défis. Une part non négligeable des pièces circule encore
sous forme de scans, parfois de photocopies de mauvaise qualité, portant
tampons, signatures et annotations manuscrites. Un pipeline unique, conçu
pour des PDF nativement numériques, échoue sur ce corpus.

Il en résulte une exigence de \btpterm{détection en amont}~: distinguer un
document numérique, pour lequel une analyse structurelle suffit, d'un
document numérisé, pour lequel une reconnaissance optique avec analyse de
mise en page s'impose. Le choix des outils concrets relève du
chapitre~\ref{chap:resultats}~; le constat d'état de l'art est que cette
bifurcation n'est plus optionnelle.

% =============================================================================
% 1.9 Correspondance approximative, résolution d'entités et graphes
% =============================================================================
\section{Résolution d'entités et graphes de connaissances}
\label{sec:ch1-entites}

\subsection{Correspondance approximative}

La correspondance approximative (\anglais{fuzzy matching}) désigne l'ensemble
des techniques permettant de comparer des chaînes malgré des variations
orthographiques, des fautes de frappe ou des artefacts d'OCR. Elle
s'inscrit dans un héritage ancien --- la distance d'édition de Levenshtein
(1965), les codes phonétiques de type Soundex, les coefficients de Jaccard
et les $n$-grammes --- et demeure indispensable dès que l'on traite des
textes produits par des humains ou par une reconnaissance optique imparfaite.

\begin{btptable}{Méthodes fondamentales de correspondance approximative}
\label{tab:fuzzy}
\begin{tabularx}{\linewidth}{@{}L{2.8cm} Y L{3.6cm}@{}}
\toprule
\rowcolor{btppale}
\btpth{Méthode} & \btpth{Principe} & \btpth{Application typique} \\
\midrule
Levenshtein & Nombre d'insertions, suppressions et substitutions. & Fautes, artefacts d'OCR. \\
Jaccard & Intersection sur union d'ensembles de tokens. & Ordre des mots secondaire. \\
$n$-grammes & Partage de sous-séquences de $n$ caractères. & Texte long, OCR. \\
Soundex et variantes & Encodage phonétique. & Noms propres (limité en français). \\
\bottomrule
\end{tabularx}
\end{btptable}

\subsection{Résolution d'entités}

La résolution d'entités consiste à décider que deux libellés désignent la
même entité réelle. Elle s'appuie sur la correspondance approximative, mais
l'intègre dans un processus~: normalisation (casse, accents, formes
juridiques), comparaison, seuillage, et, en cas de doute, arbitrage humain.

La figure~\ref{fig:resolution-entites} en illustre le déroulement sur deux
dénominations d'une même entreprise.

\begin{figure}[htbp]
  \centering
  \resizebox{0.92\linewidth}{!}{% Chaîne de résolution d'entités
\begin{tikzpicture}[node distance=7mm and 14mm]

  \node[btpboxplain, text width=3.4cm] (a) {\ttfamily\scriptsize «~SITINFRA Sarl~»};
  \node[btpboxplain, text width=3.4cm, right=26mm of a] (b) {\ttfamily\scriptsize «~Sitinfra Cameroun~»};

  \node[btpbox, text width=3.4cm, below=of a] (na) {Normalisation\\[-1pt]\btplabel{casse, accents, formes juridiques}};
  \node[btpbox, text width=3.4cm, below=of b] (nb) {Normalisation\\[-1pt]\btplabel{casse, accents, formes juridiques}};

  \node[btpboxplain, text width=3.4cm, below=of na] (ra) {\ttfamily\scriptsize sitinfra};
  \node[btpboxplain, text width=3.4cm, below=of nb] (rb) {\ttfamily\scriptsize sitinfra cameroun};

  \node[btpboxalt, text width=5.2cm, below=10mm of {$(ra)!0.5!(rb)$}] (cmp)
    {Comparaison approximative\\[-1pt]\btplabel{Levenshtein, Jaccard, $n$-grammes}};

  \node[btpdecision, text width=2.1cm, below=9mm of cmp] (dec)
    {Score $>$ seuil~?};

  \node[btpboxstrong, text width=3.0cm, below left=9mm and 4mm of dec] (same)
    {Même entité\\[-1pt]\btplabel{fusion}};
  \node[btpbox, text width=3.0cm, below right=9mm and 4mm of dec] (diff)
    {Entités distinctes\\[-1pt]\btplabel{arbitrage humain}};

  \draw[btparrow] (a) -- (na);
  \draw[btparrow] (b) -- (nb);
  \draw[btparrow] (na) -- (ra);
  \draw[btparrow] (nb) -- (rb);
  \draw[btparrow] (ra.south) -- (cmp.north west);
  \draw[btparrow] (rb.south) -- (cmp.north east);
  \draw[btparrow] (cmp) -- (dec);
  \draw[btparrow] (dec) -- node[btplabel, left, xshift=-1mm] {oui} (same);
  \draw[btparrow] (dec) -- node[btplabel, right, xshift=1mm] {non} (diff);

\end{tikzpicture}
}
  \caption{Chaîne de résolution d'entités}
  \label{fig:resolution-entites}
\end{figure}

Dans un corpus d'appels d'offres, les entités récurrentes sont nombreuses
et leurs dénominations instables~:

\begin{itemize}
  \item \btpterm{Entreprises}~: « SITINFRA », « Sitinfra Cameroun »,
        « SITINFRA Sarl »~;
  \item \btpterm{Projets}~: un même ouvrage désigné par son lieu, son maître
        d'ouvrage ou son acronyme interne~;
  \item \btpterm{Certifications}~: « ISO 9001 » et « ISO 9001:2015 »~;
  \item \btpterm{Personnes}~: variations d'état civil, de titre et d'ordre
        des noms.
\end{itemize}

Sans résolution, le système crée des doublons, manque des rapprochements
et, le cas échéant, peuple un graphe de n\oe uds redondants. Avec une
résolution trop agressive, il fusionne des entités distinctes. Le seuil
et l'arbitrage humain ne sont donc pas des raffinements~: ils sont
constitutifs de la méthode.

\subsection{Graphes de connaissances}

Un graphe de connaissances représente les entités comme des n\oe uds et
les relations comme des arêtes typées. La figure~\ref{fig:graphe}
en donne une instance schématique pour le domaine des appels d'offres.

\begin{figure}[htbp]
  \centering
  \resizebox{\linewidth}{!}{% Principe d'un graphe de connaissances appliqué aux appels d'offres
\begin{tikzpicture}[node distance=16mm and 26mm,
                    every edge quotes/.style={btplabel, fill=white, inner sep=1.5pt}]

  \node[btpboxstrong, text width=2.2cm] (ent) {ENTREPRISE};
  \node[btpbox, text width=2.0cm, right=of ent] (proj) {PROJET};
  \node[btpbox, text width=2.2cm, right=of proj] (client) {CLIENT};

  \node[btpboxalt, text width=2.4cm, below=of ent] (cert) {CERTIFICATION};
  \node[btpboxalt, text width=2.2cm, below=of proj] (pers) {PERSONNE};
  \node[btpboxalt, text width=2.2cm, below=of client] (exig) {EXIGENCE};

  \draw[btparrow] (ent) -- node[btplabel, fill=white, inner sep=1.5pt, above]
       {\textsc{a\_realise}} (proj);
  \draw[btparrow] (proj) -- node[btplabel, fill=white, inner sep=1.5pt, above]
       {\textsc{pour}} (client);
  \draw[btparrow] (ent) -- node[btplabel, fill=white, inner sep=1.5pt, left]
       {\textsc{detient}} (cert);
  \draw[btparrow] (ent) -- node[btplabel, fill=white, inner sep=1.5pt, pos=0.32, above, sloped]
       {\textsc{emploie}} (pers);
  \draw[btparrow] (pers) -- node[btplabel, fill=white, inner sep=1.5pt, pos=0.5, above, sloped]
       {\textsc{a\_travaille\_sur}} (proj);
  \draw[btparrow] (cert) -- node[btplabel, fill=white, inner sep=1.5pt, below]
       {\textsc{satisfait}} (exig);
  \draw[btparrowlight] (exig) -- node[btplabel, fill=white, inner sep=1.5pt, right]
       {\textsc{issue\_de}} (client);

  \node[btplabel, below=7mm of pers, text width=10cm]
    {Une question comme « quels ingénieurs certifiés ont travaillé sur un projet
     comparable pour ce maître d'ouvrage~? » se résout par un parcours de
     chemins dans le graphe, là où une recherche textuelle échouerait.};

\end{tikzpicture}
}
  \caption{Principe d'un graphe de connaissances appliqué aux appels d'offres}
  \label{fig:graphe}
\end{figure}

Quatre familles de questions, mal servies par une recherche textuelle,
deviennent formulables~:

\begin{enumerate}
  \item \btpterm{Traçabilité des compétences}~: quels ingénieurs ont
        travaillé sur des projets comparables~?
  \item \btpterm{Conformité relationnelle}~: quelles certifications
        couvrent quelles exigences du règlement de consultation~?
  \item \btpterm{Analyse transverse}~: quels maîtres d'ouvrage ont déjà
        attribué des marchés à l'entreprise~?
  \item \btpterm{Capitalisation des relations}~: projets, clients,
        équipements et compétences deviennent des chemins navigables.
\end{enumerate}

L'extraction automatique de tels graphes à partir de textes non structurés
reste toutefois un défi ouvert, résumé au tableau~\ref{tab:defis-graphe}.

\begin{btptable}{Défis de l'extraction de graphes de connaissances}
\label{tab:defis-graphe}
\begin{tabularx}{\linewidth}{@{}L{3.4cm} Y Y@{}}
\toprule
\rowcolor{btppale}
\btpth{Défi} & \btpth{Description} & \btpth{Piste} \\
\midrule
Résolution d'entités & Une même entité sous des formes variées. & Normalisation et correspondance approximative. \\
Extraction de relations & Identifier le type et le sens d'un lien. & Modèle contraint, schéma prédéfini. \\
Évolution & Le graphe doit s'enrichir sans se corrompre. & Fusion incrémentale, versions. \\
Évaluation & Mesurer l'exactitude des triplets extraits. & Validation humaine par échantillon. \\
\bottomrule
\end{tabularx}
\end{btptable}

\begin{btpwarn}[Positionnement]
Le graphe de connaissances n'est pas un prérequis de la première version
d'un système de génération d'offres. Il devient pertinent lorsque les
questions relationnelles multi-sauts constituent un besoin avéré, et
lorsque le corpus est suffisamment propre pour que l'extraction ne
produise pas plus de bruit que de signal. Nous le traitons donc comme une
perspective de seconde phase, et non comme un fondement de l'architecture
initiale.
\end{btpwarn}

% =============================================================================
% 1.10 Workflows déterministes, architectures agentiques et HITL
% =============================================================================
\section{Workflows, architectures agentiques et supervision humaine}
\label{sec:ch1-workflows}

\subsection{Systèmes déterministes et systèmes agentiques}

La littérature récente distingue deux régimes de conception des systèmes
fondés sur les grands modèles de langage.

Dans un \btpterm{workflow déterministe}, le chemin d'exécution est connu à
l'avance. Le modèle intervient \emph{dans} une étape --- extraire, classer,
rédiger --- mais ne décide pas de l'enchaînement des étapes. Le code
orchestre, filtre, valide et interrompt.

Dans un \btpterm{système agentique}, le modèle décide de la prochaine
action, peut boucler, reformuler, changer d'outil et d'objectif. Cette
flexibilité est précieuse lorsque le chemin n'est pas prédictible. Elle a
un coût~: l'indétermination du parcours, la difficulté d'audit et le risque
d'une action non autorisée.

Anthropic (2024) souligne que les systèmes efficaces sont souvent construits
à partir de composants simples et composables, et que l'ajout d'agents doit
être justifié par la complexité réelle du problème, non par une préférence
esthétique pour l'autonomie. Wang et coll.\ (2023) recensent par ailleurs
les défaillances typiques des agents autonomes~: boucles infinies, dérive
d'objectif, appels d'outils injustifiés.

\begin{btptable}{Comparaison des deux régimes}
\label{tab:workflow-vs-agent}
\begin{tabularx}{\linewidth}{@{}L{3.6cm} Y Y@{}}
\toprule
\rowcolor{btppale}
\btpth{Critère} & \btpth{Workflow déterministe} & \btpth{Agent autonome} \\
\midrule
Chemin d'exécution & Connu, reproductible. & Décidé à l'exécution. \\
Traçabilité & Chaque étape est journalisable. & Le journal est un récit, pas un contrat. \\
Contrôle d'erreur & Validations déterministes entre étapes. & Dépend de l'auto-critique du modèle. \\
Coût d'inférence & Borné par le graphe d'étapes. & Potentiellement non borné. \\
Pertinence pour les AO & Élevée~: le processus métier est connu. & Ponctuelle~: questions multi-sauts. \\
\bottomrule
\end{tabularx}
\end{btptable}

Pour les appels d'offres, le processus métier est connu~: déposer, analyser,
extraire, rechercher, rédiger une section, faire valider, exporter. Ce
n'est pas un espace d'exploration. Une approche workflow, dans laquelle
l'intelligence est \emph{locale} à chaque étape, est donc le régime par
défaut. Un n\oe ud agentique peut, plus tard, être encapsulé dans une
étape précise --- par exemple une recherche multi-sauts --- sans changer le
reste du processus.

\subsection{Supervision humaine dans la boucle}

La supervision humaine (\anglais{Human-In-The-Loop}, HITL) désigne
l'insertion d'un opérateur dans le cycle de production, à des points
d'interruption définis. Elle n'est pas un aveu d'échec de l'automatisation~:
elle est la traduction opérationnelle du caractère engageant de l'offre.

Trois modalités se distinguent~:

\begin{description}
  \item[Approbation.] L'opérateur accepte la section telle quelle.
  \item[Correction.] L'opérateur modifie le texte~; la version corrigée
        devient la référence, et le système en conserve la trace.
  \item[Relance.] L'opérateur rejette la section et redéclenche la
        génération, éventuellement avec une consigne supplémentaire.
\end{description}

\begin{btpbox}[Ce que la supervision humaine garantit]
Elle garantit trois propriétés que aucun modèle ne peut certifier seul~:
la \btpterm{responsabilité} (un humain signe l'offre), la
\btpterm{conformité} (un expert vérifie que la section répond à
l'exigence) et l'\btpterm{apprentissage organisationnel} (les corrections
alimentent le patrimoine pour les dossiers suivants).
\end{btpbox}

\subsection{Architectures pour la génération de documents}

La génération d'un document long par un modèle de langage se heurte à une
limite désormais bien documentée~: la dégradation de l'attention au milieu
des contextes longs (Liu et coll., 2023). Plus le document à produire est
étendu, plus une génération globale perd en fidélité, en structure et en
couverture des exigences.

Trois stratégies s'opposent~:

\begin{btptable}{Stratégies de génération de documents longs}
\label{tab:strategies-generation}
\begin{tabularx}{\linewidth}{@{}L{3.6cm} Y Y@{}}
\toprule
\rowcolor{btppale}
\btpth{Stratégie} & \btpth{Avantage} & \btpth{Limite} \\
\midrule
Génération globale & Simple, un seul appel. & Perte de structure, citations diluées, pas de versions intermédiaires. \\
Génération par section & Contrôle local, citations ciblées, versions. & Exige un plan préalable et une orchestration. \\
Génération par section avec revue & Idem, plus responsabilité humaine. & Temps de cycle plus long. \\
\bottomrule
\end{tabularx}
\end{btptable}

Pour un dossier d'offre, dont le plan est largement imposé par le dossier
entrant, la génération par section n'est pas un raffinement~: elle est la
seule stratégie qui rende possibles la citation, la revue et le versionnage.
Nous en ferons le c\oe ur de la méthodologie au chapitre~\ref{chap:methodologie}.

% =============================================================================
% 1.11 Solutions existantes et synthèse de l'état de l'art
% =============================================================================
\section{Solutions existantes et synthèse}
\label{sec:ch1-solutions}

\subsection{Solutions du marché}

Le marché de la gestion des appels d'offres et de la gestion de projet BTP
propose plusieurs familles d'outils. Le tableau~\ref{tab:solutions-marche}
en donne un aperçu représentatif, sans prétendre à l'exhaustivité.

\begin{btptable}{Familles de solutions existantes}
\label{tab:solutions-marche}
\begin{tabularx}{\linewidth}{@{}L{2.8cm} Y C{1.6cm} C{2.0cm}@{}}
\toprule
\rowcolor{btppale}
\btpth{Famille} & \btpth{Fonctionnalités typiques} & \btpth{IA générative} & \btpth{Capitalisation} \\
\midrule
Gestion d'appels d'offres & Suivi des consultations, collaboration, modèles de réponse. & Limitée & Passive \\
Gestion de projet BTP & Planning, coûts, documents de chantier. & Absente & Passive \\
Veille des marchés publics & Détection d'opportunités, alertes. & Absente & Sans objet \\
GED généraliste & Archivage, circuits, recherche plein texte. & Émergente & Passive \\
\bottomrule
\end{tabularx}
\end{btptable}

Ces solutions partagent quatre limites, déjà identifiées à la
section~\ref{sec:ch1-ged} et que l'état de l'art technologique permet
désormais d'adresser~: recherche sémantique absente ou rudimentaire,
absence de modélisation des relations, capitalisation réduite à
l'archivage, et génération automatique inexistante ou non sourcée.

Aucune, à notre connaissance, n'articule dans un même système une ingestion
structurelle des pièces d'un DAO, une recherche hybride sur un patrimoine
d'entreprise, une génération section par section avec citations, et une
supervision humaine obligatoire avant export --- le tout dans le cadre
juridique des marchés publics camerounais.

\subsection{Tendances émergentes}

Quatre tendances se dessinent néanmoins dans la littérature et chez les
éditeurs~:

\begin{enumerate}
  \item l'usage des grands modèles de langage pour produire un premier
        brouillon de réponse~;
  \item l'extraction automatique des exigences d'un dossier entrant~;
  \item le RAG comme architecture d'aide à la rédaction, plutôt que comme
        chatbot générique~;
  \item l'orchestration de workflows couvrant le processus de bout en bout.
\end{enumerate}

Ces tendances confirment la pertinence de notre problématique. Elles en
confirment aussi le risque~: une génération non contrainte, sans citation
ni revue, n'est pas une aide, c'est un nouveau vecteur d'irrégularité.

\subsection{Synthèse de l'état de l'art}

Le tableau~\ref{tab:synthese-eoa} croise la maturité des domaines parcourus
et leur pertinence pour la réponse aux appels d'offres.

\begin{btptable}{Maturité des domaines et pertinence pour les appels d'offres}
\label{tab:synthese-eoa}
\begin{tabularx}{\linewidth}{@{}L{4.4cm} C{3.0cm} Y@{}}
\toprule
\rowcolor{btppale}
\btpth{Domaine} & \btpth{Maturité} & \btpth{Pertinence} \\
\midrule
Traitement automatique du langage & \btpMature & Très élevée \\
Grands modèles de langage & \btpMature{} (évolutive) & Très élevée \\
Recherche d'information lexicale (BM25) & \btpMature & Très élevée \\
Fusion de classements (RRF) & \btpMature & Très élevée \\
Génération augmentée par récupération & \btpMature & Très élevée \\
Bases vectorielles & \btpMature & Très élevée \\
Compréhension de documents & \btpEnDev & Très élevée \\
Résolution d'entités & \btpMature & Élevée \\
Graphes de connaissances & \btpEnDev & Élevée en seconde phase \\
Agents fondés sur les LLM & \btpEmergent & Ponctuelle, non centrale \\
Supervision humaine dans la boucle & \btpMature & Constitutive \\
\bottomrule
\end{tabularx}
\end{btptable}

De ce panorama se dégagent cinq lacunes que le présent travail entend
combler~:

\begin{enumerate}
  \item l'absence de solution intégrée couvrant le cycle analyse--rédaction
        --validation--export d'un dossier d'offre~;
  \item le manque de traçabilité des sources dans les approches génératives
        existantes~;
  \item une capitalisation encore passive, réduite à l'archivage~;
  \item l'absence de supervision humaine comme étape obligatoire, et non
        comme option~;
  \item la non-adaptation au cadre juridique et documentaire des marchés
        publics camerounais.
\end{enumerate}

\begin{btpbox}[Ce que le chapitre suivant doit trancher]
L'état de l'art justifie une architecture de génération augmentée par
récupération, adossée à une recherche hybride, à une ingestion
structurelle et à une revue humaine. Il ne dit pas \emph{comment}
articuler ces briques, ni selon quels critères les évaluer. C'est l'objet
du chapitre~\ref{chap:methodologie}.
\end{btpbox}

\btpEndMark


% --- Chapitre 2 : Méthodologie --------------------------------------------
\chapter{Méthodologie}
\label{chap:methodologie}
\btpIntro{L}{a méthodologie exposée dans ce chapitre s'inscrit dans une
démarche de recherche et développement appliquée. Elle articule l'analyse du
domaine métier, la spécification des besoins, la modélisation des données,
la conception de l'architecture et la définition des stratégies de recherche,
de génération et d'évaluation. Chaque choix est explicité et justifié au
regard des exigences de traçabilité, de conformité et de validation humaine
propres aux dossiers d'appels d'offres.}

% =============================================================================
% 2.1 Introduction
% =============================================================================
\section{Introduction}
\label{sec:ch2-intro}

La méthodologie adoptée s'inscrit dans une démarche de recherche et
développement appliquée, visant à concevoir, développer et valider une
plateforme de capitalisation et de génération assistée de dossiers d'appels
d'offres. Elle combine des principes du génie logiciel, des méthodes de
l'intelligence artificielle et une analyse du domaine métier des marchés
publics camerounais.

Ce chapitre n'est pas un journal de bord du développement. Il explicite les
choix, leurs justifications et les garanties qu'ils produisent. Un lecteur
compétent doit pouvoir y suivre le problème métier, le flux de traitement
de bout en bout, la place de chaque brique, et la frontière entre le
périmètre de la première version et ce qui est explicitement reporté.

Nous procédons en huit temps~: formulation de la problématique et de
l'approche générale~; analyse des besoins et critères d'évaluation~;
modélisation des données~; architecture du système~; stratégies de
recherche et de génération~; workflow de traitement~; protocole
d'évaluation~; planification.

% =============================================================================
% 2.2 Problématique et approche générale
% =============================================================================
\section{Problématique et approche générale}
\label{sec:ch2-problematique}

\subsection{Problématique métier}

La question centrale, déjà posée en introduction générale, se reformule ici
en termes opératoires~:

\begin{btpquote}
Comment concevoir et développer une plateforme permettant de capitaliser le
patrimoine documentaire d'un bureau d'études du BTP et de générer
automatiquement des dossiers de réponse aux appels d'offres en s'appuyant
sur ce patrimoine, tout en garantissant la traçabilité des sources, la
conformité réglementaire et la validation humaine des livrables~?
\end{btpquote}

Trois tensions structurent la réponse, et chacune interdit une solution
naïve.

\begin{description}
  \item[Couverture contre précision.] Trop peu d'extraits et l'on rate une
        clause~; trop d'extraits et le modèle se perd dans un contexte
        bruité.
  \item[Texte libre contre structure.] Une partie du savoir est narrative
        (méthodologie d'exécution)~; une autre est tabulaire (moyens,
        montants, certifications). Un traitement uniforme échoue sur l'une
        ou l'autre.
  \item[Autonomie contre responsabilité.] Une offre engage l'entreprise.
        Toute architecture qui délègue la décision à un agent autonome
        déplace la responsabilité sans l'assumer.
\end{description}

\subsection{Approche proposée}

L'approche générale repose sur cinq piliers, qui découlent directement de
l'état de l'art du chapitre~\ref{chap:etat-art}.

\begin{enumerate}
  \item \btpterm{Capitalisation structurée.} Transformer le patrimoine
        documentaire en une base de connaissances interrogeable, en
        préservant la hiérarchie des pièces.
  \item \btpterm{Recherche hybride.} Combiner recherche lexicale et
        recherche sémantique, fusionner par rang réciproque, puis
        ré-ordonner.
  \item \btpterm{Génération contrôlée.} Produire des sections, non des
        documents entiers, sous schéma contraint, avec citations et
        abstentions explicites.
  \item \btpterm{Workflow validé.} Orchestrer le processus avec
        interruption obligatoire pour la revue humaine.
  \item \btpterm{Architecture modulaire.} Concevoir par composants
        indépendants, afin qu'une brique --- le graphe, un n\oe ud
        agentique --- puisse s'ajouter sans reconstruire le reste.
\end{enumerate}

\begin{btpnote}[Ce que le système n'est pas]
Le système visé n'est pas un assistant conversationnel qui rédige un appel
d'offres d'un seul tenant. Ce n'est pas non plus un essaim d'agents
autonomes. C'est une plateforme de capitalisation et de génération
assistée, organisée en quatre capacités~: mémoire documentaire, intelligence
de l'appel d'offres, espace de travail assisté, génération documentaire
versionnée.
\end{btpnote}

% =============================================================================
% 2.3 Analyse des besoins et spécification
% =============================================================================
\section{Analyse des besoins et spécification}
\label{sec:ch2-besoins}

\subsection{Besoins fonctionnels}

L'analyse des besoins a été menée en collaboration avec des praticiens du
domaine des appels d'offres au sein de SITINFRA Cameroun. Elle a permis
d'identifier les besoins fonctionnels du tableau~\ref{tab:besoins-f}.

\begin{btptable}{Besoins fonctionnels}
\label{tab:besoins-f}
\begin{tabularx}{\linewidth}{@{}C{1.1cm} L{3.6cm} C{2.0cm} Y@{}}
\toprule
\rowcolor{btppale}
\btpth{Id.} & \btpth{Besoin} & \btpth{Priorité} & \btpth{Description} \\
\midrule
BF1  & Ingestion de documents     & Haute    & Dépôt de fichiers (PDF, DOCX, tableurs, scans). \\
BF2  & Indexation                 & Haute    & Indexation sémantique et lexicale des extraits. \\
BF3  & Recherche                  & Haute    & Recherche hybride, filtrée par organisation et type. \\
BF4  & Extraction d'exigences     & Haute    & Extraction structurée des exigences du DAO. \\
BF5  & Génération de sections     & Haute    & Rédaction assistée d'une section, avec citations. \\
BF6  & Validation humaine         & Haute    & Revue, correction ou relance avant intégration. \\
BF7  & Export                     & Moyenne  & Assemblage et export du dossier final. \\
BF8  & Graphe de connaissances    & Basse    & Liens entre entités (seconde phase). \\
BF9  & Cloisonnement              & Moyenne  & Isolation des données par organisation. \\
BF10 & Historique                 & Moyenne  & Versions des sections et trace des décisions. \\
\bottomrule
\end{tabularx}
\end{btptable}

\subsection{Besoins non fonctionnels}

\begin{btptable}{Besoins non fonctionnels}
\label{tab:besoins-nf}
\begin{tabularx}{\linewidth}{@{}C{1.3cm} L{2.8cm} Y@{}}
\toprule
\rowcolor{btppale}
\btpth{Id.} & \btpth{Besoin} & \btpth{Description} \\
\midrule
BNF1 & Performance     & Temps de réponse inférieur à cinq secondes pour une recherche. \\
BNF2 & Fiabilité       & Taux d'erreur opérationnelle inférieur à un pour cent. \\
BNF3 & Sécurité        & Authentification, cloisonnement strict des corpus. \\
BNF4 & Traçabilité     & Chaque génération est rattachée à des sources et à une version. \\
BNF5 & Évolutivité     & Architecture modulaire, ajout de briques sans rupture. \\
BNF6 & Maintenabilité  & Code structuré, documenté, couvert par des tests. \\
\bottomrule
\end{tabularx}
\end{btptable}

\subsection{Spécifications fonctionnelles}

\paragraph{Ingestion et indexation (BF1, BF2).}
Le dépôt d'un fichier via l'interface web déclenche une détection du type
(numérique ou numérisé), une extraction préservant titres et tableaux, un
découpage hiérarchique parent-enfant, puis une indexation conjointe dans
un magasin vectoriel et une base relationnelle. Le document n'est marqué
« prêt » que lorsque les deux magasins ont confirmé l'écriture.

\paragraph{Extraction d'exigences (BF4).}
Le dossier d'appel d'offres courant est analysé pour en extraire les
exigences --- catégorie, caractère obligatoire, pièce source --- stockées
en base relationnelle et liées au dossier en cours. Cette extraction est
contrainte par un schéma~: le modèle ne produit pas un texte, il produit
des enregistrements validables.

\paragraph{Génération de sections (BF5, BF6).}
Un plan est d'abord construit à partir du dossier entrant. Pour chaque
section, le système construit une requête, récupère un contexte, génère
une sortie structurée, la valide automatiquement, puis s'interrompt pour
la revue humaine. Aucune section n'est intégrée au livrable sans décision
explicite.

\subsection{Critères d'évaluation}

Les critères du tableau~\ref{tab:criteres-eval} seront utilisés au
chapitre~\ref{chap:resultats} pour valider le système. Ils dérivent des
hypothèses H1 à H5 formulées en introduction générale.

\begin{btptable}{Critères d'évaluation et cibles}
\label{tab:criteres-eval}
\begin{tabularx}{\linewidth}{@{}L{3.4cm} Y C{2.6cm}@{}}
\toprule
\rowcolor{btppale}
\btpth{Critère} & \btpth{Mesure} & \btpth{Cible} \\
\midrule
Rappel à 20        & Proportion d'extraits pertinents dans les vingt premiers. & $> 0{,}75$ \\
Fidélité           & Proportion d'affirmations fondées sur le contexte. & $> 0{,}85$ \\
Temps d'ingestion  & Document numérique de cinquante pages. & $< 60$\,s \\
Temps de génération & Une section, hors revue humaine. & $< 30$\,s \\
Taux de validation & Sections approuvées sans modification. & $> 70$\,\% \\
\bottomrule
\end{tabularx}
\end{btptable}

% =============================================================================
% 2.4 Modélisation des données
% =============================================================================
\section{Modélisation des données}
\label{sec:ch2-donnees}

\subsection{Modèle conceptuel}

Le modèle conceptuel identifie les entités du domaine, représentées à la
figure~\ref{fig:modele-conceptuel}.

\begin{figure}[htbp]
  \centering
  \resizebox{\linewidth}{!}{% Modèle conceptuel des données
\begin{tikzpicture}[node distance=11mm and 16mm,
                    card/.style={btplabel, fill=white, inner sep=1.2pt}]

  \node[btpboxstrong, text width=2.0cm] (tenant) {ORGANISATION};

  \node[btpbox, text width=2.0cm, below left=12mm and 20mm of tenant] (user) {UTILISATEUR};
  \node[btpbox, text width=2.0cm, below=12mm of tenant] (doc) {DOCUMENT};
  \node[btpbox, text width=2.0cm, below right=12mm and 20mm of tenant] (dossier) {DOSSIER};

  \node[btpboxalt, text width=2.0cm, below=of doc] (chunk) {EXTRAIT};

  \node[btpboxalt, text width=2.0cm, below left=of dossier] (exig) {EXIGENCE};
  \node[btpboxalt, text width=2.0cm, below=of dossier] (plan) {PLAN};
  \node[btpboxalt, text width=2.0cm, below right=of dossier] (export) {EXPORT};

  \node[btpbox, text width=2.0cm, below=of plan] (section) {SECTION};
  \node[btpbox, text width=2.0cm, below=of section] (revue) {REVUE};

  \draw[btparrow] (tenant) -- node[card, pos=0.75] {1..n} (user);
  \draw[btparrow] (tenant) -- node[card, pos=0.75] {1..n} (doc);
  \draw[btparrow] (tenant) -- node[card, pos=0.75] {1..n} (dossier);
  \draw[btparrow] (doc) -- node[card, pos=0.6] {1..n} (chunk);
  \draw[btparrow] (dossier) -- node[card, pos=0.65] {1..n} (exig);
  \draw[btparrow] (dossier) -- node[card, pos=0.6] {1..n} (plan);
  \draw[btparrow] (dossier) -- node[card, pos=0.65] {1..n} (export);
  \draw[btparrow] (plan) -- node[card, pos=0.6] {1..n} (section);
  \draw[btparrow] (section) -- node[card, pos=0.6] {1..n} (revue);

  % relation réflexive parent-enfant
  \draw[btparrowlight] (chunk.west) .. controls +(-1.2,0.4) and +(-1.2,-0.4) .. (chunk.south west)
    node[card, left, xshift=-9mm, yshift=2mm] {parent / enfant};

  % citation d'un extrait par une section
  \draw[btpdashed] (section.west) -- node[card, above, sloped] {cite} (chunk.south east);

\end{tikzpicture}
}
  \caption{Modèle conceptuel des données}
  \label{fig:modele-conceptuel}
\end{figure}

Les entités se regroupent en quatre grappes.

\begin{description}
  \item[Organisation.] L'organisation cliente de la plateforme isole
        utilisateurs, documents, dossiers et conversations. Ce
        cloisonnement n'est pas une option de déploiement~: c'est une
        contrainte de confidentialité.
  \item[Mémoire documentaire.] Un document déposé est découpé en extraits.
        Un extrait peut être parent d'autres extraits. L'identifiant de
        l'extrait est la clé de citation.
  \item[Dossier en cours.] Un dossier de réponse porte des exigences
        extraites, un plan, des sections générées et des exports. Chaque
        section accumule des revues.
  \item[Dialogue.] Une conversation relie des messages à un dossier et,
        éventuellement, à une section.
\end{description}

\subsection{Relations structurantes}

Trois relations méritent d'être soulignées, car elles portent des
invariants du système.

\begin{itemize}
  \item \btpterm{Parent-enfant entre extraits.} La recherche s'exécute sur
        les enfants~; le contexte transmis au modèle est le parent. Un
        enfant sans parent est une anomalie.
  \item \btpterm{Citation d'un extrait par une section.} Une section
        générée référence des identifiants d'extraits. Une citation dont
        l'identifiant n'existe pas est rejetée par la validation
        déterministe.
  \item \btpterm{Revue d'une section.} Toute décision humaine ---
        approbation, correction, relance --- est une entité, pas un
        attribut écrasé. L'historique est appendice, jamais destructif.
\end{itemize}

\subsection{Modèle logique}

Le modèle logique, implanté dans une base relationnelle, traduit ces
entités en tables avec contraintes d'intégrité. Les extraits y sont
stockés en clair, avec leurs métadonnées (pièce, page, rôle, identifiant
parent). Les vecteurs correspondants vivent dans l'index vectoriel, liés
par le même identifiant. Cette dualité --- texte en relationnel, vecteur
en index --- est volontaire~: elle permet de reconstruire un contexte
citable même si l'index est reconstruit, et d'auditer le texte effectivement
transmis au modèle.

Le détail des colonnes et des contraintes figure au
chapitre~\ref{chap:resultats}, où il accompagne les choix d'implémentation.

% =============================================================================
% 2.5 Architecture du système
% =============================================================================
\section{Architecture du système}
\label{sec:ch2-architecture}

\subsection{Vue d'ensemble : deux horloges}

L'architecture s'organise autour de deux horloges indépendantes, représentées
à la figure~\ref{fig:deux-horloges}. Cette séparation n'est pas un détail
d'implémentation~: elle exprime que l'ingestion et l'exploitation n'ont ni
la même durée, ni le même contrat vis-à-vis de l'utilisateur.

\begin{figure}[htbp]
  \centering
  \resizebox{\linewidth}{!}{% Architecture générale : deux horloges indépendantes
\begin{tikzpicture}[node distance=4mm and 7mm]

  % --- Horloge A ---
  \node[btptitleband, text width=12.2cm, fill=btpblue] (ta)
    {Horloge A — Ingestion \textnormal{\small(asynchrone, non bloquante)}};
  \node[btpbox, text width=1.9cm, below=7mm of ta.south west, anchor=north west] (a1) {Dépôt};
  \node[btpbox, text width=1.9cm, right=of a1] (a2) {Stockage};
  \node[btpbox, text width=1.9cm, right=of a2] (a3) {Analyse\\structurelle};
  \node[btpbox, text width=1.9cm, right=of a3] (a4) {Découpage\\en extraits};
  \node[btpbox, text width=1.9cm, right=of a4] (a5) {Indexation};
  \foreach \i/\j in {1/2, 2/3, 3/4, 4/5} {\draw[btparrow] (a\i) -- (a\j);}
  \node[btplabel, below=2.5mm of a3, text width=8cm]
    {Traitement long, exécuté en arrière-plan par des travailleurs dédiés};

  % --- Horloge B ---
  \node[btptitleband, text width=12.2cm, fill=btpsky,
        below=16mm of ta.south west, anchor=north west] (tb)
    {Horloge B — Exploitation \textnormal{\small(synchrone, interactive)}};
  \node[btpboxalt, text width=1.9cm, below=7mm of tb.south west, anchor=north west] (b1) {Question};
  \node[btpboxalt, text width=1.9cm, right=of b1] (b2) {Recherche\\hybride};
  \node[btpboxalt, text width=1.9cm, right=of b2] (b3) {Contexte};
  \node[btpboxalt, text width=1.9cm, right=of b3] (b4) {Génération\\contrainte};
  \node[btpboxalt, text width=1.9cm, right=of b4] (b5) {Validation\\humaine};
  \foreach \i/\j in {1/2, 2/3, 3/4, 4/5} {\draw[btparrow] (b\i) -- (b\j);}
  \node[btplabel, below=2.5mm of b3, text width=8cm]
    {Traitement court, borné en temps, restituant un résultat vérifiable};

  % --- Liaison ---
  \draw[btparrowlight, dashed] (a5.south) -- ++(0,-0.55)
    node[btplabel, right, xshift=1mm, yshift=1mm] {base de connaissances partagée}
    -| (b2.north);

\end{tikzpicture}
}
  \caption{Architecture générale~: deux horloges indépendantes}
  \label{fig:deux-horloges}
\end{figure}

L'\btpterm{horloge A} capitalise. Elle est asynchrone, non bloquante, et
peut durer des dizaines de secondes. L'utilisateur dépose un fichier et
reçoit une confirmation immédiate~; le traitement se poursuit en
arrière-plan.

L'\btpterm{horloge B} exploite. Elle est synchrone, bornée en temps, et
doit restituer un résultat vérifiable en quelques secondes. Elle ne
s'exécute que sur des documents dont l'horloge A a confirmé l'indexation.

Les deux horloges partagent la base de connaissances. Elles ne partagent
pas le fil d'exécution~: une ingestion longue ne doit jamais geler une
recherche.

\subsection{Architecture en couches}

La figure~\ref{fig:architecture-couches} décline cette organisation en
cinq couches de responsabilité.

\begin{figure}[htbp]
  \centering
  \resizebox{0.92\linewidth}{!}{% Architecture en couches de la plateforme
\begin{tikzpicture}[node distance=4mm]

  \newcommand{\layerw}{12.4cm}

  \node[btptitleband, text width=\layerw, fill=btpdeep] (l1)
    {Couche présentation};
  \node[btpboxplain, text width=\layerw, below=1.5mm of l1, align=left, font=\scriptsize]
    (c1) {Interface web~: dépôt de documents, conversation, revue des sections
           générées, export des livrables};

  \node[btptitleband, text width=\layerw, fill=btpblue, below=of c1] (l2)
    {Couche service et orchestration};
  \node[btpboxplain, text width=\layerw, below=1.5mm of l2, align=left, font=\scriptsize]
    (c2) {Interface applicative HTTP, authentification et cloisonnement par
           organisation, routage des intentions, machine à états du workflow,
           points d'interruption pour la supervision humaine};

  \node[btptitleband, text width=\layerw, fill=btpsky, below=of c2] (l3)
    {Couche traitement};
  \node[btpboxplain, text width=\layerw, below=1.5mm of l3, align=left, font=\scriptsize]
    (c3) {Analyse structurelle des documents et reconnaissance optique,
           découpage hiérarchique, calcul des plongements, recherche hybride,
           ré-ordonnancement, extraction d'exigences, génération contrainte};

  \node[btptitleband, text width=\layerw, fill=btplight, text=btpdeep, below=of c3] (l4)
    {Couche persistance};
  \node[btpboxplain, text width=\layerw, below=1.5mm of l4, align=left, font=\scriptsize]
    (c4) {Base relationnelle (métadonnées, extraits, exigences, versions,
           décisions), index vectoriel (dense et creux), stockage objet
           (fichiers d'origine et exports), file de travaux};

  \node[btptitleband, text width=\layerw, fill=btpslate, below=of c4] (l5)
    {Couche infrastructure};
  \node[btpboxplain, text width=\layerw, below=1.5mm of l5, align=left, font=\scriptsize]
    (c5) {Conteneurisation des services, journalisation structurée,
           supervision des métriques, sauvegardes};

  \draw[btparrow, line width=1.4pt] ($(l1.west)+(-0.4,0)$) -- ($(c5.west)+(-0.4,0)$)
    node[midway, rotate=90, anchor=south, btplabel, yshift=1pt] {Descente des appels};

\end{tikzpicture}
}
  \caption{Architecture en couches de la plateforme}
  \label{fig:architecture-couches}
\end{figure}

Cette décomposition autorise le remplacement d'une brique --- le modèle de
plongement, le magasin vectoriel, le moteur d'orchestration --- sans
remettre en cause les contrats des couches adjacentes. Elle prépare aussi
la seconde phase~: un magasin graphe s'insère dans la couche de
persistance, un n\oe ud agentique dans la couche de traitement, sans
modifier l'interface ni le workflow de revue.

\subsection{Flux de données}

Trois flux principaux traversent l'architecture.

\paragraph{Flux d'ingestion.}
Fichier $\rightarrow$ stockage objet $\rightarrow$ analyse structurelle
$\rightarrow$ extraits $\rightarrow$ plongements $\rightarrow$ index
vectoriel et base relationnelle.

\paragraph{Flux de recherche.}
Question $\rightarrow$ plongement $\rightarrow$ recherche hybride filtrée
$\rightarrow$ ré-ordonnancement $\rightarrow$ contexte numéroté.

\paragraph{Flux de génération.}
Contexte et consigne $\rightarrow$ modèle de langage $\rightarrow$ section
structurée $\rightarrow$ validation déterministe $\rightarrow$ revue humaine
$\rightarrow$ export.

Le chapitre~\ref{chap:resultats} nomme les composants qui incarnent ces
flux. La méthodologie s'arrête aux rôles, afin que les choix de produits
restent substituables.

% =============================================================================
% 2.6 Stratégie de recherche et de génération
% =============================================================================
\section{Stratégie de recherche et de génération}
\label{sec:ch2-recherche}

\subsection{Recherche hybride}

La stratégie de recherche retenue combine, dans cet ordre, les opérations
justifiées au chapitre~\ref{chap:etat-art}.

\begin{enumerate}
  \item \btpterm{Recherche lexicale} (BM25), pour les termes exacts, les
        identifiants et les références normatives.
  \item \btpterm{Recherche sémantique} (vecteurs denses), pour les concepts
        et les paraphrases.
  \item \btpterm{Filtrage par métadonnées}, appliqué \emph{avant} le calcul
        de similarité~: organisation, type de pièce, rôle du document
        (patrimoine ou dossier courant).
  \item \btpterm{Fusion par rang réciproque}, qui combine les deux
        classements sans normaliser les scores.
  \item \btpterm{Ré-ordonnancement} par modèle croisé, qui resserre le
        contexte transmis au générateur.
\end{enumerate}

Le pipeline en entonnoir --- une cinquantaine de candidats fusionnés, huit
extraits ré-ordonnés --- est celui de la figure~\ref{fig:pipeline-hybride}.
Les effectifs (cinquante, huit) sont des paramètres, non des constantes
théoriques~; ils seront ajustés empiriquement au
chapitre~\ref{chap:resultats}.

\subsection{Découpage parent-enfant}

La figure~\ref{fig:parent-enfant} précise le compromis retenu entre
précision de la recherche et richesse du contexte.

\begin{figure}[htbp]
  \centering
  \resizebox{\linewidth}{!}{% Stratégie de découpage parent-enfant
\begin{tikzpicture}[node distance=6mm and 10mm]

  \node[btpbox, text width=9.6cm] (parent)
    {\textbf{Extrait parent} — 800 à 1500 unités lexicales\\[-1pt]
     \btplabel{Conservé en base relationnelle, transmis au modèle comme contexte}};

  \node[btpboxalt, text width=2.7cm, below=12mm of parent.south west, anchor=north west] (c1)
    {Enfant 1\\[-1pt]\btplabel{250--400 unités}};
  \node[btpboxalt, text width=2.7cm, right=of c1] (c2)
    {Enfant 2\\[-1pt]\btplabel{250--400 unités}};
  \node[btpboxalt, text width=2.7cm, right=of c2] (c3)
    {Enfant 3\\[-1pt]\btplabel{250--400 unités}};

  \foreach \c in {c1, c2, c3} {
    \draw[btparrowlight] (parent.south) -- (\c.north);
  }

  \node[btplabel, below=5mm of c2, text width=10cm]
    {Les enfants seuls sont indexés dans l'espace vectoriel~: ils rendent la
     recherche précise. Les parents correspondants sont ensuite récupérés,
     dédupliqués, puis envoyés au modèle~: ils fournissent un contexte
     suffisant pour rédiger.};

  \node[btpboxstrong, text width=2.4cm, right=14mm of parent] (idx) {Recherche\\sur les enfants};
  \draw[btparrow] (idx.south) |- ($(c3.east)+(0.35,0)$) -- (c3.east);

\end{tikzpicture}
}
  \caption{Stratégie de découpage parent-enfant}
  \label{fig:parent-enfant}
\end{figure}

Les bornes --- 250 à 400 unités lexicales pour un enfant, 800 à 1500 pour
un parent --- s'alignent sur deux contraintes. En dessous, un enfant ne
porte plus assez de signal pour être retrouvé. Au-dessus, un parent
réintroduit le phénomène de « perte au milieu » documenté par Liu et
coll.\ (2023). Le découpage suit en outre la structure du document
(titres, articles) plutôt qu'une fenêtre glissante, afin qu'un extrait
coincide autant que possible avec une unité juridique ou technique.

\subsection{Génération par section}

La génération s'exécute section par section, jamais sur le dossier entier.
Pour chaque section du plan~:

\begin{enumerate}
  \item construction d'une requête de recherche à partir du titre, des
        exigences liées et, le cas échéant, du tour de dialogue~;
  \item recherche hybride et assemblage d'un contexte numéroté~;
  \item appel au modèle avec une sortie contrainte par un schéma~;
  \item validation déterministe (schéma, citations existantes, absence de
        titres inventés, faits chiffrés cités)~;
  \item interruption pour la revue humaine~;
  \item enregistrement de la version.
\end{enumerate}

La sortie n'est pas un texte libre. Elle obéit à un schéma dont l'annexe
donne une instance, et dont les champs obligatoires sont~: la clé de
section, le titre, les blocs (paragraphes, listes, tableaux), les
identifiants d'extraits cités, les exigences couvertes, et les
abstentions explicites.

\begin{lstlisting}[style=btpjson, caption={Schéma de sortie d'une section générée},
                   label={lst:schema-section}]
{
  "section_key": "references",
  "titre": "5. References",
  "blocs": [
    {"type": "paragraphe", "texte": "..."},
    {"type": "liste", "texte": "- Projet A\\n- Projet B"}
  ],
  "citations": ["chunk-001", "chunk-045"],
  "exigences_couvertes": ["ex-001", "ex-003"],
  "abstentions": ["Pas de reference superieure a 20 M EUR"]
}
\end{lstlisting}

Les caractères accentués sont volontairement dépouillés dans cet exemple
de schéma, afin d'éviter toute ambiguïté d'encodage. Le système de
production accepte le français complet.

\subsection{Stratégie graphe (seconde phase)}

La stratégie graphe est reportée. Les types d'entités envisagés
(entreprise, projet, client, certification, personne, exigence) et les
relations (\textsc{a\_realise}, \textsc{detient}, \textsc{emploie},
\textsc{satisfait}) sont déjà identifiés. Leur extraction n'entre dans le
périmètre qu'après validation empirique du RAG hybride, afin de ne pas
superposer deux sources d'erreur --- retrieval et extraction de triplets
--- avant d'avoir mesuré chacune.

% =============================================================================
% 2.7 Workflow de traitement
% =============================================================================
\section{Workflow de traitement}
\label{sec:ch2-workflow}

\subsection{Processus d'ingestion}

Le processus d'ingestion, porté par l'horloge A, est détaillé à la
figure~\ref{fig:ingestion}.

\begin{figure}[htbp]
  \centering
  \resizebox{\linewidth}{!}{% Chaîne d'ingestion documentaire
\begin{tikzpicture}[node distance=5mm and 9mm]

  \node[btpboxstrong, text width=2.1cm] (up) {Dépôt du\\fichier};
  \node[btpdb, right=of up, minimum width=2.0cm] (obj) {Stockage\\objet};
  \node[btpbox, text width=2.3cm, right=of obj] (detect)
    {Détection\\du type\\[-1pt]\btplabel{numérique / numérisé}};

  \node[btpbox, text width=2.5cm, above right=4mm and 9mm of detect] (parse1)
    {Analyse\\structurelle\\[-1pt]\btplabel{PDF numérique}};
  \node[btpbox, text width=2.5cm, below right=4mm and 9mm of detect] (parse2)
    {Reconnaissance\\optique\\[-1pt]\btplabel{document numérisé}};

  \node[btpbox, text width=2.4cm, right=9mm of {$(parse1)!0.5!(parse2)$}] (chunk)
    {Découpage\\hiérarchique};
  \node[btpboxalt, text width=2.3cm, below=14mm of chunk] (embed)
    {Calcul des\\plongements};
  \node[btpdb, left=12mm of embed, minimum width=2.1cm] (vec) {Index\\vectoriel};
  \node[btpdb, right=12mm of embed, minimum width=2.3cm] (sql) {Base\\relationnelle};

  \draw[btparrow] (up) -- (obj);
  \draw[btparrow] (obj) -- (detect);
  \draw[btparrow] (detect) -- (parse1);
  \draw[btparrow] (detect) -- (parse2);
  \draw[btparrow] (parse1) -- (chunk);
  \draw[btparrow] (parse2) -- (chunk);
  \draw[btparrow] (chunk) -- (embed);
  \draw[btparrow] (embed) -- (vec);
  \draw[btparrow] (embed) -- (sql);

  \node[btplabel, below=4mm of embed, text width=9.5cm]
    {Le statut du document passe de « en attente » à « prêt » seulement
     lorsque l'indexation est confirmée dans les deux magasins};

\end{tikzpicture}
}
  \caption{Chaîne d'ingestion documentaire}
  \label{fig:ingestion}
\end{figure}

Cinq étapes s'enchaînent.

\begin{enumerate}
  \item \btpterm{Dépôt.} L'utilisateur transmet un fichier. Le système
        enregistre le binaire dans le stockage objet, crée un
        enregistrement « en attente » et place un travail dans la file.
        La réponse HTTP est immédiate.
  \item \btpterm{Analyse.} Un travailleur détecte si le document est
        numérique ou numérisé, puis applique l'analyse structurelle ou la
        reconnaissance optique avec préservation de la mise en page.
  \item \btpterm{Découpage.} Les titres et articles structurent le
        découpage parent-enfant. Un préfixe de métadonnées (type de pièce,
        page, identifiant) est injecté dans le texte indexé, afin que la
        recherche lexicale bénéficie de ces ancres.
  \item \btpterm{Indexation.} Les plongements sont calculés et écrits dans
        l'index vectoriel~; le texte et les métadonnées le sont dans la
        base relationnelle. Le statut passe à « prêt » seulement après
        confirmation des deux écritures.
  \item \btpterm{Graphe (phase ultérieure).} Un travail distinct extraire
        entités et relations, fusionnées dans le magasin graphe.
\end{enumerate}

\subsection{Processus de génération}

Le processus de génération, porté par l'horloge B, est représenté à la
figure~\ref{fig:workflow-generation}.

\begin{figure}[htbp]
  \centering
  \resizebox{0.92\linewidth}{!}{% Workflow de génération d'une section avec supervision humaine
\begin{tikzpicture}[node distance=5.5mm and 12mm]

  \node[btpboxstrong, text width=2.6cm] (s1) {1. Routage de\\l'intention};
  \node[btpbox, text width=2.6cm, below=of s1] (s2) {2. Construction\\de la requête};
  \node[btpbox, text width=2.6cm, below=of s2] (s3) {3. Recherche\\hybride};
  \node[btpbox, text width=2.6cm, below=of s3] (s4) {4. Assemblage\\du contexte};
  \node[btpboxalt, text width=2.6cm, below=of s4] (s5) {5. Génération\\structurée};

  \node[btpdecision, text width=2.3cm, right=16mm of s5] (v)
    {Validation\\automatique};

  \node[btpbox, text width=2.7cm, above=12mm of v] (retry)
    {Reprise bornée\\[-1pt]\btplabel{2 tentatives au plus}};

  \node[btpboxstrong, text width=2.7cm, above=12mm of retry, fill=btpsky] (hitl)
    {6. Revue humaine\\[-1pt]\btplabel{approuver / corriger / relancer}};

  \node[btpbox, text width=2.7cm, above=12mm of hitl] (save)
    {7. Enregistrement\\de la version};

  \foreach \i/\j in {1/2, 2/3, 3/4, 4/5} {\draw[btparrow] (s\i) -- (s\j);}
  \draw[btparrow] (s5) -- (v);
  \draw[btparrow] (v) -- node[btplabel, right, xshift=1mm] {échec} (retry);
  \draw[btparrow] (retry) -- node[btplabel, right, xshift=1mm] {succès} (hitl);
  \draw[btparrow] (v.east) -- ++(0.7,0) |- node[btplabel, pos=0.25, right] {succès} (hitl.east);
  \draw[btparrow] (hitl) -- (save);
  \draw[btparrow] (retry.west) -- ++(-0.6,0) |- (s5.west);

  \node[btplabel, below=3mm of s5, text width=4.4cm]
    {Sortie contrainte par un schéma~: blocs, citations, exigences couvertes,
     abstentions explicites};

\end{tikzpicture}
}
  \caption{Workflow de génération d'une section avec supervision humaine}
  \label{fig:workflow-generation}
\end{figure}

La validation automatique qui précède la revue humaine n'est pas un
second modèle juge. Ce sont des contrôles déterministes~: le document
produit est-il conforme au schéma~? Les identifiants cités existent-ils~?
Aucun titre de premier ou deuxième niveau n'a-t-il été inventé hors plan~?
Tout fait chiffré est-il accompagné d'une citation~? Un échec déclenche
au plus deux reprises~; au-delà, la section est présentée à l'humain avec
le diagnostic.

\subsection{Routage des intentions}

L'entrée conversationnelle de l'espace de travail n'est pas un dialogue
libre. Un routeur, combinant règles et modèle, classe l'intention avant
tout appel coûteux.

\begin{btptable}{Routage des intentions}
\label{tab:intentions}
\begin{tabularx}{\linewidth}{@{}L{2.6cm} Y L{3.4cm} L{2.4cm}@{}}
\toprule
\rowcolor{btppale}
\btpth{Intention} & \btpth{Indices} & \btpth{Traitement} & \btpth{Sortie} \\
\midrule
Export & « exporter », « télécharger », « pdf » & Assemblage immédiat des sections approuvées. & Fichier \\
Générer une section & « génère », « rédige », « section » & Recherche, génération, persistance, revue. & Section \\
Question documentaire & Autre question ancrée dans le dossier & Recherche hybride et réponse citée. & Réponse \\
Échange social & Salutation, absence d'ancrage & Dialogue court, sans recherche. & Réponse \\
\bottomrule
\end{tabularx}
\end{btptable}

Ce routage évite deux écueils~: lancer une génération coûteuse pour une
simple salutation, et lancer un export alors qu'aucune section n'est
encore approuvée. Dans le second cas, le système refuse plutôt que de
produire un livrable incomplet.

% =============================================================================
% 2.8 Stratégie d'évaluation et de validation
% =============================================================================
\section{Stratégie d'évaluation et de validation}
\label{sec:ch2-evaluation}

\subsection{Évaluation de la récupération}

L'évaluation du moteur de recherche s'appuie sur un jeu de référence
(\anglais{golden set}) de trente questions, dont les extraits pertinents
auront été identifiés manuellement par des praticiens. Les métriques
retenues sont le rappel à 20, la précision à 5 et le rang réciproque
moyen.

Quatre configurations seront comparées~: recherche dense seule, BM25 seul,
hybride par fusion de rangs, hybride suivi d'un ré-ordonnancement. Les
seuils d'interprétation sont les suivants~: un rappel à 20 inférieur à
$0{,}70$ signale un problème d'ingestion ou de découpage~; un rappel
supérieur à $0{,}80$ indique un système opérationnel.

\subsection{Évaluation de la génération}

L'évaluation de la génération est confiée à des experts du domaine, sur
un dossier de test. Chaque section est notée selon quatre critères~:
fidélité au contexte, ancrage des affirmations, exhaustivité au regard
des exigences, exactitude factuelle. Un protocole complémentaire, de type
modèle juge (RAGAS), accélère l'itération mais ne substitue pas le
jugement expert.

Les seuils d'alerte sont une fidélité inférieure à $0{,}80$ (problème de
contexte ou d'invite) et un ancrage inférieur à $0{,}70$ (problème de
citation).

\subsection{Validation des hypothèses}

Les hypothèses H1 à H5, formulées en introduction générale, se traduisent
en comparaisons contrôlées.

\begin{btptable}{Validation empirique des hypothèses}
\label{tab:validation-h}
\begin{tabularx}{\linewidth}{@{}C{1.1cm} Y L{3.8cm} C{1.8cm}@{}}
\toprule
\rowcolor{btppale}
\btpth{Réf.} & \btpth{Énoncé opérationnel} & \btpth{Comparaison} & \btpth{Gain visé} \\
\midrule
H1 & La recherche hybride surpasse la recherche dense seule. & Rappel à 20 & $+10$\,\% \\
H2 & Le découpage parent-enfant surpasse la fenêtre fixe. & Ancrage et exhaustivité & $+15$\,\% \\
H3 & Le ré-ordonnancement surpasse la seule fusion de rangs. & Précision à 5 & $+20$\,\% \\
H4 & La génération par section surpasse la génération globale. & Couverture des exigences & $+25$\,\% \\
H5 & La revue humaine améliore la fiabilité sans annuler le gain. & Taux de validation, temps de cycle & $+30$\,\% \\
\bottomrule
\end{tabularx}
\end{btptable}

\subsection{Robustesse}

Cinq dispositifs bornent le comportement du système en cas de défaillance.

\begin{enumerate}
  \item des validations déterministes avant et après chaque appel au
        modèle~;
  \item des reprises bornées (deux essais au plus) avant interruption
        humaine~;
  \item des délais d'expiration sur les travaux longs, avec arrêt propre~;
  \item des transactions pour la cohérence des écritures duales (relationnel
        et vectoriel)~;
  \item des journaux structurés, exploitables pour l'audit et le diagnostic.
\end{enumerate}

% =============================================================================
% 2.9 Planification et synthèse de la méthodologie
% =============================================================================
\section{Planification et synthèse}
\label{sec:ch2-planification}

\subsection{Phases de développement}

Le développement s'organise en cinq phases, représentées à la
figure~\ref{fig:phases}. Les quatre premières constituent le périmètre de
la première version. La cinquième est optionnelle et conditionnée par les
résultats des précédentes.

\begin{figure}[htbp]
  \centering
  \resizebox{\linewidth}{!}{% Phases de développement
\begin{tikzpicture}[x=1.05cm, y=1cm]

  % axe des semaines
  \draw[btplight, line width=0.8pt] (0,0) -- (10.4,0);
  \foreach \x/\lab in {0/S1, 2/S3, 4/S5, 6/S7, 8/S9, 10/S11} {
    \draw[btplight, line width=0.8pt] (\x,0) -- (\x,-0.14);
    \node[btplabel, below=0pt] at (\x,-0.14) {\lab};
  }

  \newcommand{\phase}[5]{%
    \fill[#5, rounded corners=2pt] (#1,#3-0.24) rectangle (#2,#3+0.24);
    \node[font=\scriptsize\bfseries, text=white, anchor=west] at (#1+0.12,#3) {#4};
  }

  \phase{0}{2}{3.2}{Phase 1 — Socle}{btpdeep}
  \phase{2}{4}{2.5}{Phase 2 — Ingestion}{btpblue}
  \phase{4}{6}{1.8}{Phase 3 — Recherche hybride}{btpblue!85}
  \phase{6}{8}{1.1}{Phase 4 — Workflow et supervision}{btpsky}
  \phase{8}{10}{0.4}{Phase 5 — Graphe (option)}{btpslate}

  \node[btplabel, anchor=west] at (10.5,3.2) {Infrastructure, schéma, sécurité};
  \node[btplabel, anchor=west] at (10.5,2.5) {Analyse, découpage, indexation};
  \node[btplabel, anchor=west] at (10.5,1.8) {Fusion, ré-ordonnancement};
  \node[btplabel, anchor=west] at (10.5,1.1) {Génération, revue, export};
  \node[btplabel, anchor=west] at (10.5,0.4) {Entités et relations};

\end{tikzpicture}
}
  \caption{Phasage du développement}
  \label{fig:phases}
\end{figure}

\begin{description}
  \item[Phase 1 --- Socle.] Infrastructure, schéma relationnel,
        authentification, points d'entrée minimaux.
  \item[Phase 2 --- Ingestion.] Analyse structurelle, reconnaissance
        optique, découpage parent-enfant, indexation.
  \item[Phase 3 --- Recherche hybride.] BM25, vecteurs denses, fusion,
        ré-ordonnancement, construction de contexte.
  \item[Phase 4 --- Workflow.] Orchestration, revue humaine, export,
        interface.
  \item[Phase 5 --- Graphe (optionnelle).] Extraction d'entités et de
        relations, requêtes multi-sauts.
\end{description}

\subsection{Livrables}

Les livrables attendus sont~: le code source testé, un environnement
reproductible, une documentation technique (interface applicative et
architecture), une documentation utilisateur, un jeu de test (documents et
questions) et les rapports d'évaluation correspondant aux hypothèses H1
à H5.

\subsection{Synthèse de la méthodologie}

La méthodologie se caractérise par une approche modulaire et incrémentale,
une évaluation continue, une validation humaine constitutive, une
traçabilité complète et une adaptabilité volontairement préservée. Les
choix suivants sont particulièrement notables.

\begin{btptable}{Choix méthodologiques et justifications}
\label{tab:choix-methodo}
\begin{tabularx}{\linewidth}{@{}L{4.2cm} Y@{}}
\toprule
\rowcolor{btppale}
\btpth{Choix} & \btpth{Justification} \\
\midrule
Recherche hybride & Meilleur compromis entre termes exacts et concepts. \\
Découpage parent-enfant & Recherche précise et contexte suffisant. \\
Génération par section & Contrôle de qualité, citations, versions. \\
Supervision humaine & Responsabilité, conformité, apprentissage. \\
Workflow déterministe & Traçabilité, coûts bornés, auditabilité. \\
Ingestion asynchrone & L'utilisateur n'attend pas l'indexation. \\
Relationnel et vecteurs & Faits métier d'un côté, similarité de l'autre. \\
\bottomrule
\end{tabularx}
\end{btptable}

Le chapitre suivant incarne ces choix dans une plateforme, en nomme les
composants, et en mesure les performances.

\btpEndMark


% --- Chapitre 3 : Résultats et architecture technique ---------------------
\chapter{Résultats et architecture technique}
\label{chap:resultats}
\btpIntro{C}{e chapitre présente la plateforme développée, son architecture
technique et les résultats obtenus. Nous détaillons les composants retenus,
mesurons les performances d'ingestion, de recherche et de génération, puis
comparons les approches candidates. La discussion met en regard les forces du
système, ses limites actuelles et les pistes d'amélioration identifiées.}

% =============================================================================
% 3.1 Présentation de la plateforme
% =============================================================================
\section{Présentation de la plateforme développée}
\label{sec:ch3-plateforme}

La plateforme implémente la méthodologie du chapitre~\ref{chap:methodologie}.
Elle permet l'ingestion de documents hétérogènes, l'indexation conjointe
lexicale et vectorielle, la recherche hybride, la génération contrainte
de sections, la revue humaine et l'export d'un livrable versionné.

\subsection{Interface utilisateur}

L'interface s'organise en quatre espaces, qui correspondent aux quatre
capacités identifiées en méthodologie.

\begin{itemize}
  \item un espace de dépôt et de gestion des documents, avec indication du
        statut d'indexation (en attente, en cours, prêt, en échec)~;
  \item un espace de conversation, point d'entrée unique pour les questions,
        les demandes de génération et les exports~;
  \item un espace de revue, où chaque section générée est présentée avec
        ses citations, ses exigences couvertes et ses abstentions~;
  \item un espace d'export, qui n'assemble que les sections approuvées.
\end{itemize}

\subsection{Périmètre de la version évaluée}

La version évaluée dans ce chapitre couvre les phases 1 à 4 du phasage.
Le graphe de connaissances n'est pas activé. Les mesures portent donc sur
le RAG hybride, le découpage parent-enfant, le ré-ordonnancement et la
revue humaine, et non sur des questions relationnelles multi-sauts.

% =============================================================================
% 3.2 Architecture technique
% =============================================================================
\section{Architecture technique}
\label{sec:ch3-technique}

Ce chapitre nomme les composants qui incarnent les rôles définis en
méthodologie. Le tableau~\ref{tab:composants} en donne la correspondance.

\begin{btptable}{Correspondance rôles méthodologiques --- composants techniques}
\label{tab:composants}
\begin{tabularx}{\linewidth}{@{}L{4.0cm} Y L{3.6cm}@{}}
\toprule
\rowcolor{btppale}
\btpth{Rôle} & \btpth{Composant} & \btpth{Justification brève} \\
\midrule
Interface applicative & FastAPI & Contrats HTTP explicites, validation Pydantic. \\
File de travaux & ARQ et Redis & Ingestion asynchrone, reprises. \\
Analyse de documents & Docling / PaddleOCR & Numérique vs numérisé. \\
Indexation et recherche & LlamaIndex et Qdrant & Hybride dense et creux natif. \\
Plongements & BGE-M3 & Mixte, multilingue, français inclus. \\
Ré-ordonnancement & BGE-reranker-v2-m3 & Modèle croisé, couple (requête, extrait). \\
Orchestration & LangGraph & Graphes d'étapes et interruptions HITL. \\
Modèle de langage & Ollama (local) ou API & Exécution locale prioritaire. \\
Faits métier & PostgreSQL & Versions, citations, revues. \\
Fichiers & MinIO & Originaux et exports. \\
Graphe (phase 5) & Neo4j & Reporté. \\
\bottomrule
\end{tabularx}
\end{btptable}

\subsection{Choix d'extraction documentaire}

Deux outils coexistent, conformément à la bifurcation identifiée à la
section~\ref{sec:ch1-documents}. Docling est retenu pour les PDF
numériques, car il préserve titres, listes et tableaux. PaddleOCR, avec
analyse de mise en page, est retenu pour les scans. Le basculement est
automatique, fondé sur la détection d'une couche texte extractible.

\subsection{Choix d'indexation}

Qdrant est retenu pour trois propriétés alignées sur le cahier des charges
de la section~\ref{sec:ch1-vecteurs}~: filtrage par métadonnées avant
similarité, représentation mixte dense et creuse, identifiants stables
rattachables aux extraits stockés dans PostgreSQL. BGE-M3 produit
conjointement les deux représentations, ce qui évite d'entretenir deux
modèles de plongement.

\subsection{Choix d'orchestration}

LangGraph est retenu parce qu'il expose nativement des graphes d'étapes
interruptibles. Le point d'interruption de revue n'est pas un rappel
applicatif greffé après coup~: il est un n\oe ud du graphe, ce qui
garantit qu'aucune section n'est persistée comme approuvée sans décision
humaine. LlamaIndex fournit les abstractions de récupération (enfants,
parents, fusion, ré-ordonnancement) sans dicter l'orchestration.

\begin{btpnote}[Substitutabilité]
Ces choix sont des incarnations, pas des fondements. Remplacer Qdrant par
un autre magasin hybride, ou Ollama par une API, ne change pas le
workflow. C'est précisément le bénéfice de l'architecture en couches.
\end{btpnote}

% =============================================================================
% 3.3 Performances
% =============================================================================
\section{Performances et évaluation}
\label{sec:ch3-perfs}

Les mesures ci-dessous portent sur un corpus de dossiers d'appels d'offres
et de pièces patrimoniales fournis par SITINFRA Cameroun, complété par des
documents publics de marchés. Elles constituent une première campagne, à
interpréter comme une démonstration de faisabilité plutôt que comme une
évaluation définitive.

\subsection{Performance de l'ingestion}

\begin{btptable}{Temps d'ingestion observés}
\label{tab:perf-ingestion}
\begin{tabularx}{\linewidth}{@{}C{2.4cm} L{3.6cm} Y@{}}
\toprule
\rowcolor{btppale}
\btpth{Taille} & \btpth{Type} & \btpth{Durée observée} \\
\midrule
20 pages  & PDF numérique & 8 à 12\,s \\
50 pages  & PDF numérique & 15 à 25\,s \\
20 pages  & Scan          & 30 à 45\,s \\
50 pages  & Scan          & 60 à 90\,s \\
\bottomrule
\end{tabularx}
\end{btptable}

La cible méthodologique (moins de soixante secondes pour cinquante pages
numériques) est tenue. Elle ne l'est pas pour les scans de même volume,
ce qui confirme que la reconnaissance optique reste le goulot
d'étranglement de l'horloge A. Ce dépassement n'affecte pas l'horloge B,
précisément parce que les deux horloges sont découplées.

\subsection{Performance de la recherche}

\begin{btptable}{Comparaison des approches de recherche}
\label{tab:perf-recherche}
\begin{tabularx}{\linewidth}{@{}L{4.0cm} C{2.4cm} C{2.6cm} C{2.2cm}@{}}
\toprule
\rowcolor{btppale}
\btpth{Approche} & \btpth{Rappel@20} & \btpth{Précision@5} & \btpth{Temps} \\
\midrule
Dense seul            & $0{,}72$ & $0{,}45$ & $0{,}3$\,s \\
BM25 seul             & $0{,}68$ & $0{,}38$ & $0{,}1$\,s \\
Hybride (RRF)         & $0{,}78$ & $0{,}52$ & $0{,}4$\,s \\
Hybride et ré-ordonnancement & $0{,}81$ & $0{,}61$ & $0{,}6$\,s \\
\bottomrule
\end{tabularx}
\end{btptable}

L'hypothèse H1 est soutenue~: le hybride gagne $0{,}06$ de rappel à 20
sur le dense seul, soit un peu plus de huit pour cent relatifs, proche
de la cible de dix pour cent. L'hypothèse H3 l'est également~: le
ré-ordonnancement gagne $0{,}09$ de précision à 5 sur la seule fusion,
soit environ dix-sept pour cent relatifs, légèrement sous la cible de
vingt pour cent, mais dans le bon sens et à un surcoût de latence
acceptable ($0{,}2$\,s).

\subsection{Performance de la génération}

\begin{btptable}{Scores de génération (moyenne sur cinq sections, trois évaluateurs)}
\label{tab:perf-generation}
\begin{tabularx}{\linewidth}{@{}L{4.4cm} C{3.0cm}@{}}
\toprule
\rowcolor{btppale}
\btpth{Métrique} & \btpth{Score} \\
\midrule
Fidélité au contexte           & $0{,}87$ \\
Ancrage des affirmations       & $0{,}82$ \\
Exhaustivité                   & $0{,}78$ \\
Taux de validation sans retouche & $72$\,\% \\
\bottomrule
\end{tabularx}
\end{btptable}

Les cibles de fidélité ($> 0{,}85$) et de taux de validation ($> 70$\,\%)
sont atteintes. L'ancrage ($0{,}82$) dépasse le seuil d'alerte de $0{,}70$
mais laisse une marge de progrès, principalement sur les faits chiffrés
dont la citation est parfois omise par le modèle malgré la consigne.

\subsection{Coûts}

\begin{btptable}{Coûts estimés pour dix dossiers par mois}
\label{tab:couts}
\begin{tabularx}{\linewidth}{@{}L{5.2cm} Y@{}}
\toprule
\rowcolor{btppale}
\btpth{Poste} & \btpth{Ordre de grandeur} \\
\midrule
Plongements (BGE-M3, local)     & Négligeable une fois le modèle chargé \\
Modèle de langage local (Ollama) & Nul en licence, coût matériel déjà amorti \\
Modèle de langage par API (option) & 10 à 30\,€ \\
Stockage objet et bases         & 5 à 10\,€ \\
\bottomrule
\end{tabularx}
\end{btptable}

L'exécution locale du modèle de langage, lorsqu'un accélérateur est
disponible, maintient le coût marginal proche de zéro et évite de faire
sortir des pièces confidentielles vers un prestataire externe. C'est un
argument de sécurité autant que de budget.

% =============================================================================
% 3.4 Analyse comparative
% =============================================================================
\section{Analyse comparative}
\label{sec:ch3-comparaison}

\subsection{Approches de recherche}

Le tableau~\ref{tab:perf-recherche} se relit comme une confirmation de la
complémentarité argumentée au chapitre~\ref{chap:etat-art}. BM25 seul est
le plus rapide et le plus faible en précision~: il retrouve les
identifiants, mais rate les reformulations. Le dense seul inverse le
profil. La fusion récupère le meilleur des deux classements. Le
ré-ordonnancement n'améliore que faiblement le rappel --- les extraits
pertinents étaient déjà dans les cinquante candidats --- mais améliore
nettement la tête de liste, ce qui est précisément ce que le générateur
consomme.

\subsection{Approches de découpage}

\begin{btptable}{Comparaison des approches de découpage}
\label{tab:comp-chunking}
\begin{tabularx}{\linewidth}{@{}L{3.8cm} C{3.0cm} C{3.0cm}@{}}
\toprule
\rowcolor{btppale}
\btpth{Approche} & \btpth{Ancrage} & \btpth{Exhaustivité} \\
\midrule
Fenêtre fixe          & $0{,}65$ & $0{,}55$ \\
Découpage structurel  & $0{,}78$ & $0{,}70$ \\
Parent-enfant         & $0{,}82$ & $0{,}78$ \\
\bottomrule
\end{tabularx}
\end{btptable}

L'hypothèse H2 est soutenue. Le passage de la fenêtre fixe au
parent-enfant gagne $0{,}17$ d'ancrage, soit plus de vingt-cinq pour cent
relatifs, au-delà de la cible de quinze pour cent. Le découpage
structurel, sans hiérarchie parent-enfant, capte déjà une grande part du
gain~: respecter les articles vaut mieux que découper au hasard. Le
parent-enfant ajoute le reste, en donnant au modèle un contexte plus
large que l'unité de recherche.

\subsection{Approches de génération}

\begin{btptable}{Comparaison des approches de génération}
\label{tab:comp-generation}
\begin{tabularx}{\linewidth}{@{}L{4.2cm} C{2.6cm} C{2.4cm} C{2.2cm}@{}}
\toprule
\rowcolor{btppale}
\btpth{Approche} & \btpth{Conformité} & \btpth{Qualité} & \btpth{Versions} \\
\midrule
Tout en un                 & $0{,}55$ & $0{,}60$ & Non \\
Section par section        & $0{,}78$ & $0{,}75$ & Oui \\
Section et revue humaine   & $0{,}92$ & $0{,}88$ & Oui \\
\bottomrule
\end{tabularx}
\end{btptable}

L'hypothèse H4 est nettement soutenue~: la génération par section gagne
$0{,}23$ de conformité sur la génération globale. L'hypothèse H5 l'est
aussi~: la revue humaine porte la conformité à $0{,}92$, au prix d'un
allongement du temps de cycle. Ce surcoût est acceptable --- et, dans le
domaine, non négociable --- dès lors que l'offre engage l'entreprise.

\begin{btpbox}[Lecture croisée]
Les trois comparaisons racontent la même histoire. Chaque fois que l'on
remplace une opération globale et opaque par une opération locale,
contrôlée et, le cas échéant, interrompue, la qualité mesurable
s'améliore. L'architecture en entonnoir, le parent-enfant et la revue
humaine sont trois instances de ce même principe.
\end{btpbox}

% =============================================================================
% 3.5 Discussion
% =============================================================================
\section{Discussion}
\label{sec:ch3-discussion}

\subsection{Forces du système}

Quatre propriétés distinguent la plateforme des outils recensés à la
section~\ref{sec:ch1-solutions}.

\begin{enumerate}
  \item \btpterm{Capitalisation active.} Le patrimoine n'est plus seulement
        archivé~: il est interrogeable et génératif.
  \item \btpterm{Traçabilité.} Chaque génération est ancrée dans des
        extraits identifiés, eux-mêmes rattachés à une pièce, une page et
        une version.
  \item \btpterm{Fiabilité procédurale.} La revue humaine n'est pas une
        option de l'interface~: c'est une étape du graphe d'exécution.
  \item \btpterm{Évolutivité.} L'architecture en couches permet d'ajouter
        un graphe ou un n\oe ud agentique sans reconstruire le workflow.
\end{enumerate}

\subsection{Limites actuelles}

Ces forces ne masquent pas quatre limites, qu'il faut énoncer clairement.

\begin{enumerate}
  \item \btpterm{Dépendance au corpus.} La qualité des réponses est bornée
        par la qualité et la couverture des documents patrimoniaux. Un
        bureau d'études dont la mémoire est pauvre n'obtiendra pas, par
        magie, des références qu'il n'a pas.
  \item \btpterm{OCR imparfait.} Les scans de mauvaise qualité dégradent
        l'indexation, donc la recherche, donc la génération. Le goulot
        d'étranglement mesuré à la section~\ref{sec:ch3-perfs} le confirme.
  \item \btpterm{Asymétrie linguistique.} Les modèles restent moins
        performants en français qu'en anglais, a fortiori sur le français
        juridique des marchés publics camerounais.
  \item \btpterm{Spécificité de certains dossiers.} Des consultations très
        atypiques, hors du patrimoine, dépassent les capacités du système,
        qui doit alors s'abstenir --- ce qu'il fait, mais au prix d'une
        aide réduite.
\end{enumerate}

\subsection{Pistes d'amélioration}

\begin{enumerate}
  \item \btpterm{Adaptation de domaine.} Un affinage ou, à défaut, des
        invites et des glossaires spécialisés réduiraient l'asymétrie
        linguistique.
  \item \btpterm{Récupération agentique localisée.} Encapsuler une
        recherche multi-sauts dans un n\oe ud du workflow, sans rendre
        l'ensemble agentique.
  \item \btpterm{Graphe de connaissances.} Phase 5, une fois le RAG hybride
        consolidé.
  \item \btpterm{Hygiène du corpus.} Outils de détection des doublons, des
        attestations périmées et des scans illisibles, en amont de
        l'indexation.
  \item \btpterm{Multimodalité.} Intégration des plans, schémas et images,
        aujourd'hui hors périmètre.
\end{enumerate}

\subsection{Conclusion du chapitre}

La plateforme démontre la faisabilité d'un système de capitalisation et de
génération assistée de dossiers d'appels d'offres dans le secteur du BTP
camerounais. Les résultats obtenus montrent qu'une recherche hybride
améliore significativement la pertinence, qu'une génération contrainte
produit des sections traçables, et qu'un workflow avec revue humaine
porte la conformité à un niveau compatible avec le caractère engageant
d'une offre. Les limites identifiées dessinent un programme de travail,
non un obstacle dirimant.

\btpEndMark


% =============================================================================
% PARTIES FINALES
% =============================================================================
% =============================================================================
% Conclusion générale
% =============================================================================
\btpLiminary{Conclusion générale}

\section*{Résumé des contributions}

Ce travail de recherche a contribué sur cinq plans.

\begin{enumerate}
  \item \btpterm{Analyse du domaine.} Une caractérisation du processus de
        réponse aux appels d'offres dans le BTP camerounais, de sa
        production documentaire et de ses points de friction~: relecture
        répétée, reconstitution du patrimoine, contrôle de complétude tardif.
  \item \btpterm{État de l'art.} Un panorama raisonné des disciplines
        mobilisées --- traitement automatique du langage, recherche
        d'information, génération augmentée par récupération, compréhension
        de documents, résolution d'entités, workflows et supervision
        humaine --- en distinguant les concepts des outils.
  \item \btpterm{Méthodologie.} Une approche structurée, fondée sur deux
        horloges, une recherche hybride, un découpage parent-enfant, une
        génération par section et une revue humaine constitutive.
  \item \btpterm{Architecture.} Une plateforme modulaire qui incarne ces
        choix, évaluable et substituable pièce par pièce.
  \item \btpterm{Validation.} Une première campagne empirique qui soutient
        les hypothèses H1 à H5 et démontre la faisabilité opérationnelle
        du système.
\end{enumerate}

Le résultat n'est pas « un modèle qui rédige des offres ». C'est une
chaîne de traitement dans laquelle l'intelligence artificielle accélère
la recherche et la rédaction, et dans laquelle l'ingénieur conserve la
responsabilité de l'acte engageant.

\section*{Limites du travail}

Les limites énoncées au chapitre~\ref{chap:resultats} valent pour
l'ensemble du mémoire. La campagne d'évaluation reste de taille
modeste~; le graphe de connaissances n'est pas activé~; le français
juridique camerounais demeure un terrain difficile pour des modèles
entraînés ailleurs. Ces limites bornent la portée des conclusions, elles
ne les invalident pas.

\section*{Perspectives}

Cinq perspectives se dessinent.

\begin{enumerate}
  \item \btpterm{Industrialisation.} Passage à une solution multi-organisations,
        avec cloisonnement, journalisation d'audit et engagements de
        service.
  \item \btpterm{Graphe de connaissances.} Activation de la phase 5, une
        fois le RAG hybride consolidé, pour les questions relationnelles.
  \item \btpterm{Intelligence localement agentique.} Encapsuler une
        recherche multi-sauts dans un n\oe ud du workflow, sans rendre
        l'ensemble autonome.
  \item \btpterm{Multimodalité.} Intégration des plans, schémas et images,
        aujourd'hui absents de l'index.
  \item \btpterm{Adaptation de domaine.} Affinage ou spécialisation des
        modèles au français des marchés publics et au vocabulaire du BTP.
\end{enumerate}

Au-delà de ces prolongements techniques, le bénéfice attendu est
organisationnel~: transformer un patrimoine documentaire dormant en une
mémoire d'entreprise effectivement mobilisable, dossier après dossier,
sous le contrôle de ceux qui en portent la responsabilité.

\btpEndMark
\clearpage

% =============================================================================
% Bibliographie
% =============================================================================
\begin{thebibliography}{99}
\addcontentsline{toc}{chapter}{Bibliographie}

\bibitem{bush1945}
V. Bush,
«~As We May Think~»,
\emph{The Atlantic Monthly}, vol.~176, n\textsuperscript{o}~1, p.~101--108, 1945.

\bibitem{robertson1976}
S.~E. Robertson and K. Sparck Jones,
«~Relevance Weighting of Search Terms~»,
\emph{Journal of the American Society for Information Science},
vol.~27, n\textsuperscript{o}~3, p.~129--146, 1976.

\bibitem{salton1983}
G. Salton and M.~J. McGill,
\emph{Introduction to Modern Information Retrieval},
McGraw-Hill, 1983.

\bibitem{robertson2009bm25}
S. Robertson and H. Zaragoza,
«~The Probabilistic Relevance Framework: BM25 and Beyond~»,
\emph{Foundations and Trends in Information Retrieval},
vol.~3, n\textsuperscript{o}~4, p.~333--389, 2009.

\bibitem{cormack2009rrf}
G.~V. Cormack, C.~L.~A. Clarke and S. Buettcher,
«~Reciprocal Rank Fusion outperforms Condorcet and individual Rank Learning Methods~»,
\emph{Proceedings of SIGIR}, p.~758--759, 2009.

\bibitem{hochreiter1997lstm}
S. Hochreiter and J. Schmidhuber,
«~Long Short-Term Memory~»,
\emph{Neural Computation}, vol.~9, n\textsuperscript{o}~8, p.~1735--1780, 1997.

\bibitem{mikolov2013}
T. Mikolov \textit{et al.},
«~Efficient Estimation of Word Representations in Vector Space~»,
arXiv:1301.3781, 2013.

\bibitem{pennington2014glove}
J. Pennington, R. Socher and C.~D. Manning,
«~GloVe: Global Vectors for Word Representation~»,
\emph{Proceedings of EMNLP}, p.~1532--1543, 2014.

\bibitem{vaswani2017}
A. Vaswani \textit{et al.},
«~Attention Is All You Need~»,
\emph{Advances in Neural Information Processing Systems}, 2017.

\bibitem{devlin2019bert}
J. Devlin \textit{et al.},
«~BERT: Pre-training of Deep Bidirectional Transformers for Language Understanding~»,
\emph{Proceedings of NAACL-HLT}, p.~4171--4186, 2019.

\bibitem{conneau2020xlmr}
A. Conneau \textit{et al.},
«~Unsupervised Cross-lingual Representation Learning at Scale~»,
\emph{Proceedings of ACL}, p.~8440--8451, 2020.

\bibitem{xue2021mt5}
L. Xue \textit{et al.},
«~mT5: A Massively Multilingual Pre-trained Text-to-Text Transformer~»,
\emph{Proceedings of NAACL}, p.~483--498, 2021.

\bibitem{lewis2020rag}
P. Lewis \textit{et al.},
«~Retrieval-Augmented Generation for Knowledge-Intensive NLP Tasks~»,
\emph{Advances in Neural Information Processing Systems}, 2020.

\bibitem{gao2024survey}
Y. Gao \textit{et al.},
«~Retrieval-Augmented Generation for Large Language Models: A Survey~»,
arXiv:2312.10997, 2024.

\bibitem{liu2023lost}
N.~F. Liu \textit{et al.},
«~Lost in the Middle: How Language Models Use Long Contexts~»,
\emph{Transactions of the ACL}, 2023.

\bibitem{es2024ragas}
S. Es \textit{et al.},
«~RAGAS: Automated Evaluation of Retrieval Augmented Generation~»,
arXiv:2309.15217, 2024.

\bibitem{malkov2018hnsw}
Y.~A. Malkov and D.~A. Yashunin,
«~Efficient and robust approximate nearest neighbor search using Hierarchical Navigable Small World graphs~»,
\emph{IEEE Transactions on Pattern Analysis and Machine Intelligence},
vol.~42, n\textsuperscript{o}~4, p.~824--836, 2020.

\bibitem{chen2020layoutlmv3}
Y. Xu \textit{et al.},
«~LayoutLMv2: Multi-modal Pre-training for Visually-Rich Document Understanding~»,
\emph{Proceedings of ACL}, 2021.

\bibitem{edge2024graphrag}
D. Edge \textit{et al.},
«~From Local to Global: A Graph RAG Approach to Query-Focused Summarization~»,
Microsoft Research, arXiv:2404.16130, 2024.

\bibitem{asai2023selfrag}
A. Asai \textit{et al.},
«~Self-RAG: Learning to Retrieve, Generate, and Critique through Self-Reflection~»,
arXiv:2310.11511, 2023.

\bibitem{yao2023react}
S. Yao \textit{et al.},
«~ReAct: Synergizing Reasoning and Acting in Language Models~»,
\emph{ICLR}, 2023.

\bibitem{wang2023agents}
L. Wang \textit{et al.},
«~A Survey on Large Language Model based Autonomous Agents~»,
arXiv:2308.11432, 2023.

\bibitem{anthropic2024agents}
Anthropic,
«~Building Effective Agents~»,
note d'ingénierie, 2024.

\bibitem{iso15489}
ISO 15489-1:2016,
\emph{Information et documentation --- Gestion des documents d'activité --- Partie 1 : Concepts et principes},
Organisation internationale de normalisation, 2016.

\bibitem{iso30301}
ISO 30301:2019,
\emph{Information et documentation --- Systèmes de management des documents d'activité --- Exigences},
Organisation internationale de normalisation, 2019.

\bibitem{cmpcameroun}
République du Cameroun,
\emph{Code des marchés publics},
décret n\textsuperscript{o} 2018/366 du 20 juin 2018.

\end{thebibliography}
\clearpage


\appendix
% =============================================================================
% Annexes
% =============================================================================
\chapter{Schéma de sortie d'une section}
\label{annexe:schema}

Le schéma ci-dessous est le contrat entre le générateur et le reste du
système. Toute sortie qui s'en écarte est rejetée par la validation
déterministe, avant même la revue humaine.

\begin{lstlisting}[style=btpjson, caption={Contrat JSON d'une section générée}]
{
  "section_key": "memoire_technique",
  "titre": "Memoire technique d'execution",
  "blocs": [
    {
      "type": "paragraphe",
      "texte": "La methode d'execution retenue s'appuie sur..."
    },
    {
      "type": "liste",
      "texte": "- Phase preparatoire\n- Phase chantier\n- Reception"
    }
  ],
  "citations": ["chunk-001", "chunk-014", "chunk-027"],
  "exigences_couvertes": ["ex-012", "ex-018"],
  "abstentions": [
    "Aucune reference patrimoniale superieure a 20 M EUR n'a ete trouvee."
  ]
}
\end{lstlisting}

Les champs \texttt{citations} et \texttt{exigences\_couvertes} sont des
identifiants, jamais des libellés libres. C'est cette contrainte qui rend
la citation auditable.

\chapter{Protocole du jeu de référence}
\label{annexe:golden}

Le jeu de référence utilisé pour l'évaluation de la récupération
comprend trente questions, réparties en trois familles.

\begin{btptable}{Composition du jeu de référence}
\begin{tabularx}{\linewidth}{@{}L{4.2cm} C{2.2cm} Y@{}}
\toprule
\rowcolor{btppale}
\btpth{Famille} & \btpth{Nombre} & \btpth{Exemple} \\
\midrule
Terme exact / identifiant & 10 & « Que dit l'article 12.3 du CCAP~? » \\
Concept métier paraphrasé & 10 & « Quelles sanctions en cas de retard~? » \\
Question mixte            & 10 & « Quelles références en assainissement, norme ISO 14001~? » \\
\bottomrule
\end{tabularx}
\end{btptable}

Pour chaque question, deux à cinq extraits pertinents ont été identifiés
manuellement par un praticien. Un extrait est jugé pertinent s'il
contient une information nécessaire à une réponse complète, pas seulement
s'il partage des mots avec la question.

\chapter{Glossaire des pièces d'un dossier}
\label{annexe:pieces}

Ce glossaire complète le tableau~\ref{tab:pieces-dao}. Il n'a pas vocation
à se substituer aux définitions réglementaires.

\begin{description}[style=nextline, leftmargin=1.2em, labelindent=0pt, font=\bfseries\color{btpdeep}]
  \item[Règlement de consultation.] Pièce qui fixe les règles du jeu~:
        calendrier, modalités de remise, critères de jugement, pièces
        exigées.
  \item[CCAG.] Clauses administratives générales applicables à une
        famille de marchés (travaux, fournitures, prestations
        intellectuelles).
  \item[CCAP.] Clauses administratives particulières au marché, qui
        précisent ou dérogent au CCAG.
  \item[CCTP.] Spécifications techniques de l'ouvrage ou de la
        prestation.
  \item[DPGF.] Décomposition du prix global et forfaitaire, tableau
        structuré des coûts.
  \item[BPU.] Bordereau des prix unitaires, utilisé notamment dans les
        marchés à prix unitaires.
\end{description}

\btpEndMark
\clearpage


\end{document}
